\section{Experiments}
\label{sec:experiments}

\subsection{Setup}
\label{sec:setup}
All runs used \code{wolframscript} with \emph{15.0.1 for Microsoft Windows
(64-bit) (July 2, 2026)} on the machine holding the repository; the license
allows one kernel at a time, so every script ran in a fresh kernel in
sequence. A small harness (\code{runner.wl}, Appendix~\ref{app:harness}) loads
the program, evaluates each probe under \code{TimeConstrained} (120\,s unless
stated), captures every message that is switched on through
\code{Internal`HandlerBlock}, and records:
\begin{itemize}
\item \emph{Equal}: exact equality, \code{RootReduce[result - expected] === 0},
      between the result and the independently known value of the input (for
      probes without a known closed form, 40-digit numeric agreement with the
      input is shown instead); \textbf{NO} therefore means a wrong value;
\item \emph{Depth}: $\rdepth$ before and after;
\item wall-clock time and the list of messages.
\end{itemize}
The original program prints its diagnostics; \code{Strad} suppresses them but
not messages, and both are visible in the logs. Tables are generated from the
recorded runs by \code{export_tables.wl}.

\subsection{The original program}
\label{sec:exporig}

\paragraph{Classical identities (Table~\ref{tab:origA}).}
The program denests all 26 identities in group A that admit a denesting,
including cube, fourth, fifth, sixth and seventh roots, sums and products of
nested radicals, complex radicands, and Ramanujan's $\sqrt[3]{\sqrt[3]2-1}$,
each in well under a second. A13, $\sqrt[3]{\sqrt[3]5-\sqrt[3]4}$, is correctly
left alone: the kernel shows that its minimal polynomial has degree 27, so it
does not lie in the degree-9 field $\Q(\sqrt[3]2,\sqrt[3]5)$, and the closed
form $(\sqrt[3]2+\sqrt[3]{20}-\sqrt[3]{25})/3$ that we first expected for it is
not an identity (\code{RootReduce} shows that the cube of that expression is a
cubic irrational). Two observations temper the success:
\begin{itemize}
\item 19 of the 31 rows carry the message \code{PossibleZeroQ::ztest1}
      (``Unable to decide \ldots\ Assuming it is'') together with
      \code{N::meprec}. The candidate filter did not \emph{decide} equality; it
      assumed it.
\item A27 $(3+2\sqrt2)^{1/4}\to\sqrt{1+\sqrt2}$ is equal and smaller but not
      shallower; A30/A31, the depth-three example of Borodin, Fagin, Hopcroft
      and Tompa, comes back as
      \[
      \frac{\sqrt{319-58\sqrt{29}}-10\sqrt{11-2\sqrt{29}}+\sqrt5(3\sqrt{29}-16)}
           {-5+\sqrt{29}-\sqrt{55-10\sqrt{29}}},
      \]
      which is equal to the input and one level shallower, but far from the
      answer $\sqrt5+\sqrt{11-2\sqrt{29}}$ (found by the corrected version,
      Table~\ref{tab:fixedA}).
\end{itemize}
\begin{longtable}{@{}p{0.055\textwidth}>{\ttfamily\footnotesize\raggedright\arraybackslash}p{0.33\textwidth}>{\ttfamily\footnotesize\raggedright\arraybackslash}p{0.33\textwidth}p{0.05\textwidth}p{0.05\textwidth}p{0.05\textwidth}@{}}
\caption{Original program on classical denesting identities. ``Equal'' is exact equality (\texttt{RootReduce}) with the expected value; where an expected value was not supplied, numeric agreement with the input is shown. Depth is the syntactic radical depth before and after.}\label{tab:origA}\\
\toprule ID & Input & Output & Equal & Depth & Time (s)\\\midrule\endfirsthead
\toprule ID & Input & Output & Equal & Depth & Time (s)\\\midrule\endhead
\bottomrule\endfoot
A01 & Strad[\allowbreak{}Sqrt[\allowbreak{}3 +\allowbreak{} 2*\allowbreak{}Sqrt[\allowbreak{}2]\allowbreak{}]\allowbreak{}]\allowbreak{} & 1 +\allowbreak{} Sqrt[\allowbreak{}2]\allowbreak{} & yes & 2$\to$1 & 0.146\\
A02 & Strad[\allowbreak{}Sqrt[\allowbreak{}5 +\allowbreak{} 2*\allowbreak{}Sqrt[\allowbreak{}6]\allowbreak{}]\allowbreak{}]\allowbreak{} & Sqrt[\allowbreak{}2]\allowbreak{} +\allowbreak{} Sqrt[\allowbreak{}3]\allowbreak{} & yes & 2$\to$1 & 0.040\\
A03 & Strad[\allowbreak{}Sqrt[\allowbreak{}2 +\allowbreak{} Sqrt[\allowbreak{}3]\allowbreak{}]\allowbreak{}]\allowbreak{} & (\allowbreak{}1 +\allowbreak{} Sqrt[\allowbreak{}3]\allowbreak{})\allowbreak{}/\allowbreak{}Sqrt[\allowbreak{}2]\allowbreak{}\newline{\scriptsize\itshape PossibleZeroQ::\allowbreak{}ztest1} & yes & 2$\to$1 & 0.072\\
A04 & Strad[\allowbreak{}Sqrt[\allowbreak{}11 +\allowbreak{} 6*\allowbreak{}Sqrt[\allowbreak{}2]\allowbreak{}]\allowbreak{}]\allowbreak{} & 3 +\allowbreak{} Sqrt[\allowbreak{}2]\allowbreak{} & yes & 2$\to$1 & 0.016\\
A05 & Strad[\allowbreak{}Sqrt[\allowbreak{}3 +\allowbreak{} Sqrt[\allowbreak{}5]\allowbreak{}]\allowbreak{}]\allowbreak{} & (\allowbreak{}1 +\allowbreak{} Sqrt[\allowbreak{}5]\allowbreak{})\allowbreak{}/\allowbreak{}Sqrt[\allowbreak{}2]\allowbreak{}\newline{\scriptsize\itshape PossibleZeroQ::\allowbreak{}ztest1} & yes & 2$\to$1 & 0.030\\
A06 & Strad[\allowbreak{}Sqrt[\allowbreak{}7 +\allowbreak{} 2*\allowbreak{}Sqrt[\allowbreak{}10]\allowbreak{}]\allowbreak{}]\allowbreak{} & Sqrt[\allowbreak{}2]\allowbreak{} +\allowbreak{} Sqrt[\allowbreak{}5]\allowbreak{} & yes & 2$\to$1 & 0.030\\
A07 & Strad[\allowbreak{}Sqrt[\allowbreak{}6 +\allowbreak{} 2*\allowbreak{}Sqrt[\allowbreak{}2]\allowbreak{} +\allowbreak{} 2*\allowbreak{}Sqrt[\allowbreak{}3]\allowbreak{} +\allowbreak{} 2*\allowbreak{}Sqrt[\allowbreak{}6]\allowbreak{}]\allowbreak{}]\allowbreak{} & 1 +\allowbreak{} Sqrt[\allowbreak{}2]\allowbreak{} +\allowbreak{} Sqrt[\allowbreak{}3]\allowbreak{}\newline{\scriptsize\itshape PossibleZeroQ::\allowbreak{}ztest1, General::\allowbreak{}stop} & yes & 2$\to$1 & 0.054\\
A08 & Strad[\allowbreak{}Sqrt[\allowbreak{}12 +\allowbreak{} 2*\allowbreak{}Sqrt[\allowbreak{}6]\allowbreak{} +\allowbreak{} 2*\allowbreak{}Sqrt[\allowbreak{}14]\allowbreak{} +\allowbreak{} 2*\allowbreak{}Sqrt[\allowbreak{}21]\allowbreak{}]\allowbreak{}]\allowbreak{} & Sqrt[\allowbreak{}2]\allowbreak{} +\allowbreak{} Sqrt[\allowbreak{}3]\allowbreak{} +\allowbreak{} Sqrt[\allowbreak{}7]\allowbreak{}\newline{\scriptsize\itshape PossibleZeroQ::\allowbreak{}ztest1} & yes & 2$\to$1 & 0.096\\
A09 & Strad[\allowbreak{}(\allowbreak{}Sqrt[\allowbreak{}5]\allowbreak{} -\allowbreak{} 2)\allowbreak{}\textasciicircum{}\allowbreak{}(\allowbreak{}1/\allowbreak{}3)\allowbreak{}]\allowbreak{} & (\allowbreak{}-1 +\allowbreak{} Sqrt[\allowbreak{}5]\allowbreak{})\allowbreak{}/\allowbreak{}2\newline{\scriptsize\itshape PossibleZeroQ::\allowbreak{}ztest1} & yes & 2$\to$1 & 0.027\\
A10 & Strad[\allowbreak{}(\allowbreak{}2 +\allowbreak{} Sqrt[\allowbreak{}5]\allowbreak{})\allowbreak{}\textasciicircum{}\allowbreak{}(\allowbreak{}1/\allowbreak{}3)\allowbreak{}]\allowbreak{} & (\allowbreak{}1 +\allowbreak{} Sqrt[\allowbreak{}5]\allowbreak{})\allowbreak{}/\allowbreak{}2\newline{\scriptsize\itshape PossibleZeroQ::\allowbreak{}ztest1} & yes & 2$\to$1 & 0.019\\
A11 & Strad[\allowbreak{}(\allowbreak{}10 +\allowbreak{} 6*\allowbreak{}Sqrt[\allowbreak{}3]\allowbreak{})\allowbreak{}\textasciicircum{}\allowbreak{}(\allowbreak{}1/\allowbreak{}3)\allowbreak{}]\allowbreak{} & 1 +\allowbreak{} Sqrt[\allowbreak{}3]\allowbreak{}\newline{\scriptsize\itshape PossibleZeroQ::\allowbreak{}ztest1} & yes & 2$\to$1 & 0.029\\
A12 & Strad[\allowbreak{}(\allowbreak{}2\textasciicircum{}\allowbreak{}(\allowbreak{}1/\allowbreak{}3)\allowbreak{} -\allowbreak{} 1)\allowbreak{}\textasciicircum{}\allowbreak{}(\allowbreak{}1/\allowbreak{}3)\allowbreak{}]\allowbreak{} & (\allowbreak{}1 -\allowbreak{} 2\textasciicircum{}\allowbreak{}(\allowbreak{}1/\allowbreak{}3)\allowbreak{} +\allowbreak{} 2\textasciicircum{}\allowbreak{}(\allowbreak{}2/\allowbreak{}3)\allowbreak{})\allowbreak{}/\allowbreak{}3\textasciicircum{}\allowbreak{}(\allowbreak{}2/\allowbreak{}3)\allowbreak{}\newline{\scriptsize\itshape PossibleZeroQ::\allowbreak{}ztest1} & yes & 2$\to$1 & 0.052\\
A13 & Strad[\allowbreak{}(\allowbreak{}5\textasciicircum{}\allowbreak{}(\allowbreak{}1/\allowbreak{}3)\allowbreak{} -\allowbreak{} 4\textasciicircum{}\allowbreak{}(\allowbreak{}1/\allowbreak{}3)\allowbreak{})\allowbreak{}\textasciicircum{}\allowbreak{}(\allowbreak{}1/\allowbreak{}3)\allowbreak{}]\allowbreak{} & (\allowbreak{}-2\textasciicircum{}\allowbreak{}(\allowbreak{}2/\allowbreak{}3)\allowbreak{} +\allowbreak{} 5\textasciicircum{}\allowbreak{}(\allowbreak{}1/\allowbreak{}3)\allowbreak{})\allowbreak{}\textasciicircum{}\allowbreak{}(\allowbreak{}1/\allowbreak{}3)\allowbreak{} & yes & 2$\to$2 & 1.063\\
A14 & Strad[\allowbreak{}(\allowbreak{}(\allowbreak{}11 +\allowbreak{} 5*\allowbreak{}Sqrt[\allowbreak{}5]\allowbreak{})\allowbreak{}/\allowbreak{}2)\allowbreak{}\textasciicircum{}\allowbreak{}(\allowbreak{}1/\allowbreak{}5)\allowbreak{}]\allowbreak{} & (\allowbreak{}1 +\allowbreak{} Sqrt[\allowbreak{}5]\allowbreak{})\allowbreak{}/\allowbreak{}2\newline{\scriptsize\itshape PossibleZeroQ::\allowbreak{}ztest1} & yes & 2$\to$1 & 0.031\\
A15 & Strad[\allowbreak{}Sqrt[\allowbreak{}1 +\allowbreak{} 2*\allowbreak{}2\textasciicircum{}\allowbreak{}(\allowbreak{}1/\allowbreak{}3)\allowbreak{} +\allowbreak{} 2\textasciicircum{}\allowbreak{}(\allowbreak{}2/\allowbreak{}3)\allowbreak{}]\allowbreak{}]\allowbreak{} & 1 +\allowbreak{} 2\textasciicircum{}\allowbreak{}(\allowbreak{}1/\allowbreak{}3)\allowbreak{} & yes & 2$\to$1 & 0.004\\
A16 & Strad[\allowbreak{}(\allowbreak{}17 +\allowbreak{} 12*\allowbreak{}Sqrt[\allowbreak{}2]\allowbreak{})\allowbreak{}\textasciicircum{}\allowbreak{}(\allowbreak{}1/\allowbreak{}4)\allowbreak{}]\allowbreak{} & 1 +\allowbreak{} Sqrt[\allowbreak{}2]\allowbreak{}\newline{\scriptsize\itshape PossibleZeroQ::\allowbreak{}ztest1} & yes & 2$\to$1 & 0.018\\
A17 & Strad[\allowbreak{}(\allowbreak{}99 +\allowbreak{} 70*\allowbreak{}Sqrt[\allowbreak{}2]\allowbreak{})\allowbreak{}\textasciicircum{}\allowbreak{}(\allowbreak{}1/\allowbreak{}6)\allowbreak{}]\allowbreak{} & 1 +\allowbreak{} Sqrt[\allowbreak{}2]\allowbreak{}\newline{\scriptsize\itshape PossibleZeroQ::\allowbreak{}ztest1} & yes & 2$\to$1 & 0.019\\
A18 & Strad[\allowbreak{}(\allowbreak{}3 +\allowbreak{} 2*\allowbreak{}Sqrt[\allowbreak{}2]\allowbreak{})\allowbreak{}\textasciicircum{}\allowbreak{}(\allowbreak{}3/\allowbreak{}2)\allowbreak{}]\allowbreak{} & 7 +\allowbreak{} 5*\allowbreak{}Sqrt[\allowbreak{}2]\allowbreak{}\newline{\scriptsize\itshape PossibleZeroQ::\allowbreak{}ztest1} & yes & 2$\to$1 & 0.026\\
A19 & Strad[\allowbreak{}(\allowbreak{}3 +\allowbreak{} 2*\allowbreak{}Sqrt[\allowbreak{}2]\allowbreak{})\allowbreak{}\textasciicircum{}\allowbreak{}(\allowbreak{}-2\textasciicircum{}\allowbreak{}(\allowbreak{}-1)\allowbreak{})\allowbreak{}]\allowbreak{} & -1 +\allowbreak{} Sqrt[\allowbreak{}2]\allowbreak{}\newline{\scriptsize\itshape PossibleZeroQ::\allowbreak{}ztest1} & yes & 2$\to$1 & 0.018\\
A20 & Strad[\allowbreak{}1/\allowbreak{}Sqrt[\allowbreak{}3 +\allowbreak{} 2*\allowbreak{}Sqrt[\allowbreak{}2]\allowbreak{}]\allowbreak{}]\allowbreak{} & -1 +\allowbreak{} Sqrt[\allowbreak{}2]\allowbreak{} & yes & 2$\to$1 & 0.009\\
A21 & Strad[\allowbreak{}Sqrt[\allowbreak{}3 +\allowbreak{} 2*\allowbreak{}Sqrt[\allowbreak{}2]\allowbreak{}]\allowbreak{} +\allowbreak{} Sqrt[\allowbreak{}5 +\allowbreak{} 2*\allowbreak{}Sqrt[\allowbreak{}6]\allowbreak{}]\allowbreak{}]\allowbreak{} & 1 +\allowbreak{} 2*\allowbreak{}Sqrt[\allowbreak{}2]\allowbreak{} +\allowbreak{} Sqrt[\allowbreak{}3]\allowbreak{} & yes & 2$\to$1 & 0.030\\
A22 & Strad[\allowbreak{}Sqrt[\allowbreak{}3 +\allowbreak{} 2*\allowbreak{}Sqrt[\allowbreak{}2]\allowbreak{}]\allowbreak{}*\allowbreak{}Sqrt[\allowbreak{}5 +\allowbreak{} 2*\allowbreak{}Sqrt[\allowbreak{}6]\allowbreak{}]\allowbreak{}]\allowbreak{} & 2 +\allowbreak{} Sqrt[\allowbreak{}2]\allowbreak{} +\allowbreak{} Sqrt[\allowbreak{}3]\allowbreak{} +\allowbreak{} Sqrt[\allowbreak{}6]\allowbreak{}\newline{\scriptsize\itshape PossibleZeroQ::\allowbreak{}ztest1} & yes & 2$\to$1 & 0.046\\
A23 & Strad[\allowbreak{}Sqrt[\allowbreak{}2 +\allowbreak{} Sqrt[\allowbreak{}3]\allowbreak{}]\allowbreak{} +\allowbreak{} Sqrt[\allowbreak{}2 -\allowbreak{} Sqrt[\allowbreak{}3]\allowbreak{}]\allowbreak{}]\allowbreak{} & Sqrt[\allowbreak{}6]\allowbreak{} & yes & 2$\to$1 & 0.048\\
A24 & Strad[\allowbreak{}Sqrt[\allowbreak{}-1 +\allowbreak{} 2*\allowbreak{}I*\allowbreak{}Sqrt[\allowbreak{}2]\allowbreak{}]\allowbreak{}]\allowbreak{} & 1 +\allowbreak{} I*\allowbreak{}Sqrt[\allowbreak{}2]\allowbreak{}\newline{\scriptsize\itshape PossibleZeroQ::\allowbreak{}ztest1} & yes & 2$\to$1 & 0.026\\
A25 & Strad[\allowbreak{}Sqrt[\allowbreak{}2 +\allowbreak{} 2*\allowbreak{}I*\allowbreak{}Sqrt[\allowbreak{}3]\allowbreak{}]\allowbreak{}]\allowbreak{} & I +\allowbreak{} Sqrt[\allowbreak{}3]\allowbreak{}\newline{\scriptsize\itshape PossibleZeroQ::\allowbreak{}ztest1} & yes & 2$\to$1 & 0.046\\
A26 & Strad[\allowbreak{}Sqrt[\allowbreak{}-3 -\allowbreak{} 2*\allowbreak{}Sqrt[\allowbreak{}2]\allowbreak{}]\allowbreak{}]\allowbreak{} & I*\allowbreak{}(\allowbreak{}1 +\allowbreak{} Sqrt[\allowbreak{}2]\allowbreak{})\allowbreak{} & yes & 2$\to$1 & 0.009\\
A27 & Strad[\allowbreak{}(\allowbreak{}3 +\allowbreak{} 2*\allowbreak{}Sqrt[\allowbreak{}2]\allowbreak{})\allowbreak{}\textasciicircum{}\allowbreak{}(\allowbreak{}1/\allowbreak{}4)\allowbreak{}]\allowbreak{} & Sqrt[\allowbreak{}1 +\allowbreak{} Sqrt[\allowbreak{}2]\allowbreak{}]\allowbreak{}\newline{\scriptsize\itshape PossibleZeroQ::\allowbreak{}ztest1} & yes & 2$\to$2 & 0.030\\
A28 & Strad[\allowbreak{}Sqrt[\allowbreak{}3 +\allowbreak{} 2*\allowbreak{}Sqrt[\allowbreak{}3 +\allowbreak{} 2*\allowbreak{}Sqrt[\allowbreak{}2]\allowbreak{}]\allowbreak{}]\allowbreak{},\allowbreak{} True]\allowbreak{} & Sqrt[\allowbreak{}5 +\allowbreak{} 2*\allowbreak{}Sqrt[\allowbreak{}2]\allowbreak{}]\allowbreak{} & yes & 3$\to$2 & 0.090\\
A29 & Strad[\allowbreak{}Sqrt[\allowbreak{}3 +\allowbreak{} 2*\allowbreak{}Sqrt[\allowbreak{}3 +\allowbreak{} 2*\allowbreak{}Sqrt[\allowbreak{}2]\allowbreak{}]\allowbreak{}]\allowbreak{}]\allowbreak{} & Sqrt[\allowbreak{}3 +\allowbreak{} 2*\allowbreak{}Sqrt[\allowbreak{}3 +\allowbreak{} 2*\allowbreak{}Sqrt[\allowbreak{}2]\allowbreak{}]\allowbreak{}]\allowbreak{} & yes & 3$\to$3 & 0.099\\
A30 & Strad[\allowbreak{}Sqrt[\allowbreak{}16 -\allowbreak{} 2*\allowbreak{}Sqrt[\allowbreak{}29]\allowbreak{} +\allowbreak{} 2*\allowbreak{}Sqrt[\allowbreak{}55 -\allowbreak{} 10*\allowbreak{}Sqrt[\allowbreak{}29]\allowbreak{}]\allowbreak{}]\allowbreak{},\allowbreak{} True]\allowbreak{} & (\allowbreak{}Sqrt[\allowbreak{}319 -\allowbreak{} 58*\allowbreak{}Sqrt[\allowbreak{}29]\allowbreak{}]\allowbreak{} -\allowbreak{} 10*\allowbreak{}Sqrt[\allowbreak{}11 -\allowbreak{} 2*\allowbreak{}Sqrt[\allowbreak{}29]\allowbreak{}]\allowbreak{} +\allowbreak{} Sqrt[\allowbreak{}5]\allowbreak{}*\allowbreak{}(\allowbreak{}-16 +\allowbreak{} 3*\allowbreak{}Sqrt[\allowbreak{}29]\allowbreak{})\allowbreak{})\allowbreak{}/\allowbreak{}(\allowbreak{}-5 +\allowbreak{} Sqrt[\allowbreak{}29]\allowbreak{} -\allowbreak{} Sqrt[\allowbreak{}55 -\allowbreak{} 10*\allowbreak{}Sqrt[\allowbreak{}29]\allowbreak{}]\allowbreak{})\allowbreak{}\newline{\scriptsize\itshape PossibleZeroQ::\allowbreak{}ztest1} & yes & 3$\to$2 & 0.598\\
A31 & Strad[\allowbreak{}Sqrt[\allowbreak{}16 -\allowbreak{} 2*\allowbreak{}Sqrt[\allowbreak{}29]\allowbreak{} +\allowbreak{} 2*\allowbreak{}Sqrt[\allowbreak{}55 -\allowbreak{} 10*\allowbreak{}Sqrt[\allowbreak{}29]\allowbreak{}]\allowbreak{}]\allowbreak{}]\allowbreak{} & (\allowbreak{}Sqrt[\allowbreak{}319 -\allowbreak{} 58*\allowbreak{}Sqrt[\allowbreak{}29]\allowbreak{}]\allowbreak{} -\allowbreak{} 10*\allowbreak{}Sqrt[\allowbreak{}11 -\allowbreak{} 2*\allowbreak{}Sqrt[\allowbreak{}29]\allowbreak{}]\allowbreak{} +\allowbreak{} Sqrt[\allowbreak{}5]\allowbreak{}*\allowbreak{}(\allowbreak{}-16 +\allowbreak{} 3*\allowbreak{}Sqrt[\allowbreak{}29]\allowbreak{})\allowbreak{})\allowbreak{}/\allowbreak{}(\allowbreak{}-5 +\allowbreak{} Sqrt[\allowbreak{}29]\allowbreak{} -\allowbreak{} Sqrt[\allowbreak{}55 -\allowbreak{} 10*\allowbreak{}Sqrt[\allowbreak{}29]\allowbreak{}]\allowbreak{})\allowbreak{}\newline{\scriptsize\itshape PossibleZeroQ::\allowbreak{}ztest1} & yes & 3$\to$2 & 0.118\\
\end{longtable}

\paragraph{Non-denestable and trivial inputs (Table~\ref{tab:origB}).}
Values are preserved. B01 shows the final \code{Simplify} rewriting a
non-denestable input into a different, not simpler, form; B07 shows a
\code{Root} object turned into a \emph{nested} radical (\U{20}).
\begin{longtable}{@{}p{0.055\textwidth}>{\ttfamily\footnotesize\raggedright\arraybackslash}p{0.33\textwidth}>{\ttfamily\footnotesize\raggedright\arraybackslash}p{0.33\textwidth}p{0.05\textwidth}p{0.05\textwidth}p{0.05\textwidth}@{}}
\caption{Original program on non-denestable and trivial inputs.}\label{tab:origB}\\
\toprule ID & Input & Output & Equal & Depth & Time (s)\\\midrule\endfirsthead
\toprule ID & Input & Output & Equal & Depth & Time (s)\\\midrule\endhead
\bottomrule\endfoot
B01 & Strad[\allowbreak{}Sqrt[\allowbreak{}2 +\allowbreak{} Sqrt[\allowbreak{}2]\allowbreak{}]\allowbreak{}]\allowbreak{} & 2\textasciicircum{}\allowbreak{}(\allowbreak{}1/\allowbreak{}4)\allowbreak{}*\allowbreak{}Sqrt[\allowbreak{}1 +\allowbreak{} Sqrt[\allowbreak{}2]\allowbreak{}]\allowbreak{} & yes & 2$\to$2 & 0.031\\
B02 & Strad[\allowbreak{}Sqrt[\allowbreak{}1 +\allowbreak{} Sqrt[\allowbreak{}2]\allowbreak{}]\allowbreak{}]\allowbreak{} & Sqrt[\allowbreak{}1 +\allowbreak{} Sqrt[\allowbreak{}2]\allowbreak{}]\allowbreak{} & yes & 2$\to$2 & 0.022\\
B03 & Strad[\allowbreak{}(\allowbreak{}1 +\allowbreak{} Sqrt[\allowbreak{}2]\allowbreak{})\allowbreak{}\textasciicircum{}\allowbreak{}(\allowbreak{}1/\allowbreak{}3)\allowbreak{}]\allowbreak{} & (\allowbreak{}1 +\allowbreak{} Sqrt[\allowbreak{}2]\allowbreak{})\allowbreak{}\textasciicircum{}\allowbreak{}(\allowbreak{}1/\allowbreak{}3)\allowbreak{} & yes & 2$\to$2 & 0.085\\
B04 & Strad[\allowbreak{}Sqrt[\allowbreak{}2]\allowbreak{}]\allowbreak{} & Sqrt[\allowbreak{}2]\allowbreak{} & yes & 1$\to$1 & 0.002\\
B05 & Strad[\allowbreak{}2]\allowbreak{} & 2 & yes & 0$\to$0 & 0.000\\
B06 & Strad[\allowbreak{}Sqrt[\allowbreak{}1 +\allowbreak{} I]\allowbreak{}]\allowbreak{} & Sqrt[\allowbreak{}1 +\allowbreak{} I]\allowbreak{} & yes & 1$\to$1 & 0.002\\
B07 & Strad[\allowbreak{}Root[\allowbreak{}\#1\textasciicircum{}\allowbreak{}4 -\allowbreak{} 10*\allowbreak{}\#1\textasciicircum{}\allowbreak{}2 +\allowbreak{} 1 \& ,\allowbreak{} 4]\allowbreak{}]\allowbreak{} & Sqrt[\allowbreak{}5 +\allowbreak{} 2*\allowbreak{}Sqrt[\allowbreak{}6]\allowbreak{}]\allowbreak{} & yes & 0$\to$2 & 0.006\\
B08 & Strad[\allowbreak{}Sqrt[\allowbreak{}Sqrt[\allowbreak{}2]\allowbreak{} +\allowbreak{} Sqrt[\allowbreak{}3]\allowbreak{}]\allowbreak{}]\allowbreak{} & Sqrt[\allowbreak{}Sqrt[\allowbreak{}2]\allowbreak{} +\allowbreak{} Sqrt[\allowbreak{}3]\allowbreak{}]\allowbreak{} & yes & 2$\to$2 & 0.074\\
B09 & Strad[\allowbreak{}(\allowbreak{}2 +\allowbreak{} 2*\allowbreak{}I*\allowbreak{}Sqrt[\allowbreak{}3]\allowbreak{})\allowbreak{}\textasciicircum{}\allowbreak{}(\allowbreak{}1/\allowbreak{}3)\allowbreak{}]\allowbreak{} & (\allowbreak{}2 +\allowbreak{} (\allowbreak{}2*\allowbreak{}I)\allowbreak{}*\allowbreak{}Sqrt[\allowbreak{}3]\allowbreak{})\allowbreak{}\textasciicircum{}\allowbreak{}(\allowbreak{}1/\allowbreak{}3)\allowbreak{} & yes & 2$\to$2 & 0.261\\
\end{longtable}

\paragraph{Branch counterexamples through the wrapper (Table~\ref{tab:origC}).}
C02 and C03 are the wrong values predicted by all three reviews: the program
returns $5-i\sqrt2$ for $(-1+2i\sqrt2)^{3/2}=-5+i\sqrt2$, and
$(\sqrt2-i)(\sqrt3-i)$ for $\sqrt{-1+2i\sqrt2}\sqrt{-2+2i\sqrt3}=(1+i\sqrt2)(1+i\sqrt3)$.
C01, C04, C07, C08 are right because their values happen to lie in the
principal sector of Proposition~\ref{prop:principal}
($\Arg=-\pi/2$ is still inside for $q=2$).
\begin{longtable}{@{}p{0.055\textwidth}>{\ttfamily\footnotesize\raggedright\arraybackslash}p{0.33\textwidth}>{\ttfamily\footnotesize\raggedright\arraybackslash}p{0.33\textwidth}p{0.05\textwidth}p{0.05\textwidth}p{0.05\textwidth}@{}}
\caption{Original program on branch counterexamples (wrapper level). ``Equal'' compares with the true value of the input.}\label{tab:origC}\\
\toprule ID & Input & Output & Equal & Depth & Time (s)\\\midrule\endfirsthead
\toprule ID & Input & Output & Equal & Depth & Time (s)\\\midrule\endhead
\bottomrule\endfoot
C01 & Strad[\allowbreak{}-Sqrt[\allowbreak{}3 +\allowbreak{} 2*\allowbreak{}Sqrt[\allowbreak{}2]\allowbreak{}]\allowbreak{}]\allowbreak{} & -1 -\allowbreak{} Sqrt[\allowbreak{}2]\allowbreak{} & yes & 2$\to$1 & 0.016\\
C02 & Strad[\allowbreak{}(\allowbreak{}-1 +\allowbreak{} 2*\allowbreak{}I*\allowbreak{}Sqrt[\allowbreak{}2]\allowbreak{})\allowbreak{}\textasciicircum{}\allowbreak{}(\allowbreak{}3/\allowbreak{}2)\allowbreak{}]\allowbreak{} & 5 -\allowbreak{} I*\allowbreak{}Sqrt[\allowbreak{}2]\allowbreak{}\newline{\scriptsize\itshape PossibleZeroQ::\allowbreak{}ztest1} & \textbf{NO} & 2$\to$1 & 0.396\\
C03 & Strad[\allowbreak{}Sqrt[\allowbreak{}-1 +\allowbreak{} 2*\allowbreak{}I*\allowbreak{}Sqrt[\allowbreak{}2]\allowbreak{}]\allowbreak{}*\allowbreak{}Sqrt[\allowbreak{}-2 +\allowbreak{} 2*\allowbreak{}I*\allowbreak{}Sqrt[\allowbreak{}3]\allowbreak{}]\allowbreak{}]\allowbreak{} & (\allowbreak{}-I +\allowbreak{} Sqrt[\allowbreak{}2]\allowbreak{})\allowbreak{}*\allowbreak{}(\allowbreak{}-I +\allowbreak{} Sqrt[\allowbreak{}3]\allowbreak{})\allowbreak{}\newline{\scriptsize\itshape PossibleZeroQ::\allowbreak{}ztest1} & \textbf{NO} & 2$\to$1 & 0.128\\
C04 & Strad[\allowbreak{}(\allowbreak{}-3 -\allowbreak{} 2*\allowbreak{}Sqrt[\allowbreak{}2]\allowbreak{})\allowbreak{}\textasciicircum{}\allowbreak{}(\allowbreak{}3/\allowbreak{}2)\allowbreak{}]\allowbreak{} & (\allowbreak{}-I)\allowbreak{}*\allowbreak{}(\allowbreak{}7 +\allowbreak{} 5*\allowbreak{}Sqrt[\allowbreak{}2]\allowbreak{})\allowbreak{} & yes & 2$\to$1 & 0.013\\
C05 & Strad[\allowbreak{}Sqrt[\allowbreak{}-2 -\allowbreak{} 2*\allowbreak{}Sqrt[\allowbreak{}2]\allowbreak{}]\allowbreak{}*\allowbreak{}Sqrt[\allowbreak{}-3 -\allowbreak{} 2*\allowbreak{}Sqrt[\allowbreak{}2]\allowbreak{}]\allowbreak{}]\allowbreak{} & -Sqrt[\allowbreak{}2*\allowbreak{}(\allowbreak{}7 +\allowbreak{} 5*\allowbreak{}Sqrt[\allowbreak{}2]\allowbreak{})\allowbreak{}]\allowbreak{} & yes & 2$\to$2 & 0.065\\
C06 & Strad[\allowbreak{}Sqrt[\allowbreak{}-3 -\allowbreak{} 2*\allowbreak{}Sqrt[\allowbreak{}2]\allowbreak{}]\allowbreak{}*\allowbreak{}Sqrt[\allowbreak{}-3 -\allowbreak{} 2*\allowbreak{}Sqrt[\allowbreak{}2]\allowbreak{}]\allowbreak{}]\allowbreak{} & -3 -\allowbreak{} 2*\allowbreak{}Sqrt[\allowbreak{}2]\allowbreak{} & yes & 1$\to$1 & 0.009\\
C07 & Strad[\allowbreak{}(\allowbreak{}-1 +\allowbreak{} 2*\allowbreak{}I*\allowbreak{}Sqrt[\allowbreak{}2]\allowbreak{})\allowbreak{}\textasciicircum{}\allowbreak{}(\allowbreak{}5/\allowbreak{}2)\allowbreak{}]\allowbreak{} & 1 -\allowbreak{} (\allowbreak{}11*\allowbreak{}I)\allowbreak{}*\allowbreak{}Sqrt[\allowbreak{}2]\allowbreak{}\newline{\scriptsize\itshape PossibleZeroQ::\allowbreak{}ztest1} & yes & 2$\to$1 & 0.034\\
C08 & Strad[\allowbreak{}(\allowbreak{}-5 -\allowbreak{} 2*\allowbreak{}Sqrt[\allowbreak{}6]\allowbreak{})\allowbreak{}\textasciicircum{}\allowbreak{}(\allowbreak{}3/\allowbreak{}2)\allowbreak{}]\allowbreak{} & (\allowbreak{}-I)\allowbreak{}*\allowbreak{}(\allowbreak{}11*\allowbreak{}Sqrt[\allowbreak{}2]\allowbreak{} +\allowbreak{} 9*\allowbreak{}Sqrt[\allowbreak{}3]\allowbreak{})\allowbreak{}\newline{\scriptsize\itshape PossibleZeroQ::\allowbreak{}ztest1} & yes & 2$\to$1 & 0.045\\
\end{longtable}

\paragraph{The engine called directly (Table~\ref{tab:origD}).}
D01, D04, D05 reproduce the branch loss at the engine level. D02 (multiplier
$0$) prints ``success'' and returns the unevaluated \code{{}[[1]]}; D10 (an
inexact multiplier) returns \code{Outer[][[1]]}; D11 ($\pi$) is harmless. D06
versus D07 is \RII's equal-degree gate: with $\{1,m\}$ the trial $m$ is skipped
because its degree equals the best degree, with $\{m\}$ alone it is tried and
``succeeds'' with an uglier expression of the same depth (\U{9}, \U{12}). D08
shows that \code{Extension -> Automatic} copes with a nested radicand but that
single-level mode cannot find $\sqrt{3+2\sqrt{3+2\sqrt2}}=\sqrt{5+2\sqrt2}$; it
tries 12 multipliers and gives up (compare A28). D12 shows that a ``forced''
list does not prevent generated batches. D14 shows that the imaginary
multiplier $i$ works.
\begin{longtable}{@{}p{0.055\textwidth}>{\ttfamily\footnotesize\raggedright\arraybackslash}p{0.33\textwidth}>{\ttfamily\footnotesize\raggedright\arraybackslash}p{0.33\textwidth}p{0.05\textwidth}p{0.05\textwidth}p{0.05\textwidth}@{}}
\caption{Original core \texttt{denest11} called directly with forced multipliers.}\label{tab:origD}\\
\toprule ID & Input & Output & Equal & Depth & Time (s)\\\midrule\endfirsthead
\toprule ID & Input & Output & Equal & Depth & Time (s)\\\midrule\endhead
\bottomrule\endfoot
D01 & denest11[\allowbreak{}-Sqrt[\allowbreak{}3 +\allowbreak{} 2*\allowbreak{}Sqrt[\allowbreak{}2]\allowbreak{}]\allowbreak{},\allowbreak{} \{1\}]\allowbreak{} & 1 +\allowbreak{} Sqrt[\allowbreak{}2]\allowbreak{} & \textbf{NO} & 2$\to$1 & 0.005\\
D02 & denest11[\allowbreak{}Sqrt[\allowbreak{}3 +\allowbreak{} 2*\allowbreak{}Sqrt[\allowbreak{}2]\allowbreak{}]\allowbreak{},\allowbreak{} \{0\}]\allowbreak{} & \{\}[\allowbreak{}[\allowbreak{}1]\allowbreak{}]\allowbreak{}\newline{\scriptsize\itshape Power::\allowbreak{}infy, Part::\allowbreak{}partw} & \textbf{NO} & 2$\to$0 & 0.089\\
D03 & denest11[\allowbreak{}Sqrt[\allowbreak{}3 +\allowbreak{} 2*\allowbreak{}Sqrt[\allowbreak{}2]\allowbreak{}]\allowbreak{},\allowbreak{} \{\}]\allowbreak{} & Sqrt[\allowbreak{}3 +\allowbreak{} 2*\allowbreak{}Sqrt[\allowbreak{}2]\allowbreak{}]\allowbreak{} & yes & 2$\to$2 & 0.001\\
D04 & denest11[\allowbreak{}(\allowbreak{}-1 +\allowbreak{} 2*\allowbreak{}I*\allowbreak{}Sqrt[\allowbreak{}2]\allowbreak{})\allowbreak{}\textasciicircum{}\allowbreak{}(\allowbreak{}3/\allowbreak{}2)\allowbreak{},\allowbreak{} \{1\}]\allowbreak{} & 5 -\allowbreak{} I*\allowbreak{}Sqrt[\allowbreak{}2]\allowbreak{} & \textbf{NO} & 2$\to$1 & 0.010\\
D05 & denest11[\allowbreak{}Sqrt[\allowbreak{}-1 +\allowbreak{} 2*\allowbreak{}I*\allowbreak{}Sqrt[\allowbreak{}2]\allowbreak{}]\allowbreak{}*\allowbreak{}Sqrt[\allowbreak{}-2 +\allowbreak{} 2*\allowbreak{}I*\allowbreak{}Sqrt[\allowbreak{}3]\allowbreak{}]\allowbreak{},\allowbreak{} \{1\}]\allowbreak{} & (\allowbreak{}-I +\allowbreak{} Sqrt[\allowbreak{}2]\allowbreak{})\allowbreak{}*\allowbreak{}(\allowbreak{}-I +\allowbreak{} Sqrt[\allowbreak{}3]\allowbreak{})\allowbreak{}\newline{\scriptsize\itshape PossibleZeroQ::\allowbreak{}ztest1} & \textbf{NO} & 2$\to$1 & 0.026\\
D06 & denest11[\allowbreak{}Sqrt[\allowbreak{}5 +\allowbreak{} 2*\allowbreak{}Sqrt[\allowbreak{}6]\allowbreak{}]\allowbreak{},\allowbreak{} \{1,\allowbreak{} 6 -\allowbreak{} 2*\allowbreak{}Sqrt[\allowbreak{}6]\allowbreak{} -\allowbreak{} 2*\allowbreak{}Sqrt[\allowbreak{}2]\allowbreak{} +\allowbreak{} 2*\allowbreak{}Sqrt[\allowbreak{}3]\allowbreak{}\}]\allowbreak{} & Sqrt[\allowbreak{}5 +\allowbreak{} 2*\allowbreak{}Sqrt[\allowbreak{}6]\allowbreak{}]\allowbreak{} & yes & 2$\to$2 & 0.012\\
D07 & denest11[\allowbreak{}Sqrt[\allowbreak{}5 +\allowbreak{} 2*\allowbreak{}Sqrt[\allowbreak{}6]\allowbreak{}]\allowbreak{},\allowbreak{} \{6 -\allowbreak{} 2*\allowbreak{}Sqrt[\allowbreak{}6]\allowbreak{} -\allowbreak{} 2*\allowbreak{}Sqrt[\allowbreak{}2]\allowbreak{} +\allowbreak{} 2*\allowbreak{}Sqrt[\allowbreak{}3]\allowbreak{}\}]\allowbreak{} & (\allowbreak{}1 +\allowbreak{} Sqrt[\allowbreak{}2]\allowbreak{} +\allowbreak{} Sqrt[\allowbreak{}3]\allowbreak{})\allowbreak{}/\allowbreak{}Sqrt[\allowbreak{}2*\allowbreak{}(\allowbreak{}3 -\allowbreak{} Sqrt[\allowbreak{}2]\allowbreak{} +\allowbreak{} Sqrt[\allowbreak{}3]\allowbreak{} -\allowbreak{} Sqrt[\allowbreak{}6]\allowbreak{})\allowbreak{}]\allowbreak{}\newline{\scriptsize\itshape PossibleZeroQ::\allowbreak{}ztest1} & yes & 2$\to$2 & 0.056\\
D08 & denest11[\allowbreak{}Sqrt[\allowbreak{}3 +\allowbreak{} 2*\allowbreak{}Sqrt[\allowbreak{}3 +\allowbreak{} 2*\allowbreak{}Sqrt[\allowbreak{}2]\allowbreak{}]\allowbreak{}]\allowbreak{}]\allowbreak{} & Sqrt[\allowbreak{}3 +\allowbreak{} 2*\allowbreak{}Sqrt[\allowbreak{}3 +\allowbreak{} 2*\allowbreak{}Sqrt[\allowbreak{}2]\allowbreak{}]\allowbreak{}]\allowbreak{} & yes & 3$\to$3 & 0.100\\
D09 & denest11[\allowbreak{}2 +\allowbreak{} Sqrt[\allowbreak{}3]\allowbreak{}]\allowbreak{} & 2 +\allowbreak{} Sqrt[\allowbreak{}3]\allowbreak{} & yes & 1$\to$1 & 0.000\\
D10 & denest11[\allowbreak{}Sqrt[\allowbreak{}3 +\allowbreak{} 2*\allowbreak{}Sqrt[\allowbreak{}2]\allowbreak{}]\allowbreak{},\allowbreak{} \{2.\}]\allowbreak{} & Outer[\allowbreak{}]\allowbreak{}[\allowbreak{}[\allowbreak{}1]\allowbreak{}]\allowbreak{}\newline{\scriptsize\itshape MinimalPolynomial::\allowbreak{}nalg, Outer::\allowbreak{}heads, Outer::\allowbreak{}argm, Part::\allowbreak{}partw} & \textbf{NO} & 2$\to$0 & 0.238\\
D11 & denest11[\allowbreak{}Sqrt[\allowbreak{}3 +\allowbreak{} 2*\allowbreak{}Sqrt[\allowbreak{}2]\allowbreak{}]\allowbreak{},\allowbreak{} \{Pi\}]\allowbreak{} & Sqrt[\allowbreak{}3 +\allowbreak{} 2*\allowbreak{}Sqrt[\allowbreak{}2]\allowbreak{}]\allowbreak{}\newline{\scriptsize\itshape MinimalPolynomial::\allowbreak{}nalg} & yes & 2$\to$2 & 0.004\\
D12 & denest11[\allowbreak{}Sqrt[\allowbreak{}5 +\allowbreak{} 2*\allowbreak{}Sqrt[\allowbreak{}6]\allowbreak{}]\allowbreak{},\allowbreak{} \{1,\allowbreak{} Sqrt[\allowbreak{}5]\allowbreak{}\}]\allowbreak{} & Sqrt[\allowbreak{}2]\allowbreak{} +\allowbreak{} Sqrt[\allowbreak{}3]\allowbreak{} & yes & 2$\to$1 & 0.014\\
D13 & denest11[\allowbreak{}Sqrt[\allowbreak{}3 +\allowbreak{} 2*\allowbreak{}Sqrt[\allowbreak{}2]\allowbreak{}]\allowbreak{},\allowbreak{} \{-1\}]\allowbreak{} & Sqrt[\allowbreak{}3 +\allowbreak{} 2*\allowbreak{}Sqrt[\allowbreak{}2]\allowbreak{}]\allowbreak{} & yes & 2$\to$2 & 0.005\\
D14 & denest11[\allowbreak{}Sqrt[\allowbreak{}3 +\allowbreak{} 2*\allowbreak{}Sqrt[\allowbreak{}2]\allowbreak{}]\allowbreak{},\allowbreak{} \{I\}]\allowbreak{} & 1 +\allowbreak{} Sqrt[\allowbreak{}2]\allowbreak{}\newline{\scriptsize\itshape PossibleZeroQ::\allowbreak{}ztest1} & yes & 2$\to$1 & 0.032\\
\end{longtable}

\paragraph{Symbolic input (Table~\ref{tab:origE}).}
$\sqrt{z^2}\mapsto z$ and $\sqrt{-x-y}\mapsto i\sqrt{x+y}$ are the branch-unsafe
rewrites of \U{2} reaching the public interface. E07 is \RII's marker
collision: \code{Strad[f[1], False, Identity]} returns \code{1}. E08 shows
what a user function named \code{f} wrapped around a radical provokes: a
cascade of seven different messages.
\begin{longtable}{@{}p{0.055\textwidth}>{\ttfamily\footnotesize\raggedright\arraybackslash}p{0.33\textwidth}>{\ttfamily\footnotesize\raggedright\arraybackslash}p{0.33\textwidth}p{0.05\textwidth}p{0.05\textwidth}p{0.05\textwidth}@{}}
\caption{Original program on symbolic input.}\label{tab:origE}\\
\toprule ID & Input & Output & Equal & Depth & Time (s)\\\midrule\endfirsthead
\toprule ID & Input & Output & Equal & Depth & Time (s)\\\midrule\endhead
\bottomrule\endfoot
E01 & Strad[\allowbreak{}Sqrt[\allowbreak{}z\textasciicircum{}\allowbreak{}2]\allowbreak{}]\allowbreak{} & z & -- & 1$\to$0 & 0.001\\
E02 & Strad[\allowbreak{}Sqrt[\allowbreak{}t +\allowbreak{} t\textasciicircum{}\allowbreak{}2]\allowbreak{}]\allowbreak{} & Sqrt[\allowbreak{}t*\allowbreak{}(\allowbreak{}1 +\allowbreak{} t)\allowbreak{}]\allowbreak{} & -- & 1$\to$1 & 0.002\\
E03 & Strad[\allowbreak{}Sqrt[\allowbreak{}-x -\allowbreak{} y]\allowbreak{}]\allowbreak{} & I*\allowbreak{}Sqrt[\allowbreak{}x +\allowbreak{} y]\allowbreak{} & -- & 1$\to$1 & 0.004\\
E04 & Strad[\allowbreak{}Sqrt[\allowbreak{}a +\allowbreak{} b]\allowbreak{}]\allowbreak{} & Sqrt[\allowbreak{}a +\allowbreak{} b]\allowbreak{} & -- & 1$\to$1 & 0.002\\
E05 & Strad[\allowbreak{}x]\allowbreak{} & x & -- & 0$\to$0 & 0.000\\
E06 & Strad[\allowbreak{}Sqrt[\allowbreak{}x]\allowbreak{}*\allowbreak{}Sqrt[\allowbreak{}3 +\allowbreak{} 2*\allowbreak{}Sqrt[\allowbreak{}2]\allowbreak{}]\allowbreak{}]\allowbreak{} & (\allowbreak{}1 +\allowbreak{} Sqrt[\allowbreak{}2]\allowbreak{})\allowbreak{}*\allowbreak{}Sqrt[\allowbreak{}x]\allowbreak{} & -- & 2$\to$1 & 0.010\\
E07 & Strad[\allowbreak{}f[\allowbreak{}1]\allowbreak{},\allowbreak{} False,\allowbreak{} Identity]\allowbreak{} & 1 & -- & 0$\to$0 & 0.001\\
E08 & Strad[\allowbreak{}f[\allowbreak{}Sqrt[\allowbreak{}3 +\allowbreak{} 2*\allowbreak{}Sqrt[\allowbreak{}2]\allowbreak{}]\allowbreak{}]\allowbreak{}]\allowbreak{} & f[\allowbreak{}Sqrt[\allowbreak{}3 +\allowbreak{} 2*\allowbreak{}Sqrt[\allowbreak{}2]\allowbreak{}]\allowbreak{}]\allowbreak{}\newline{\scriptsize\itshape MinimalPolynomial::\allowbreak{}nalg, General::\allowbreak{}stop, Discriminant::\allowbreak{}poly2, FactorInteger::\allowbreak{}exact, Part::\allowbreak{}partd, Union::\allowbreak{}heads, Part::\allowbreak{}pkspec1} & -- & 2$\to$2 & 0.363\\
\end{longtable}

\paragraph{Helpers (Tables~\ref{tab:helpH} and~\ref{tab:helpR}).}
\code{Factorc} preserves the value on 18 of 20 numeric probes; H05 is the
numeric counterexample of Proposition~\ref{prop:power}(2): the result
$i\sqrt{(-1)^{1/3}+\sqrt2}$ is the negative of the input. Note that
\code{Factorc[Sqrt[t + t^2]]} (H02) returns \code{Sqrt[t (1 + t)]}, so \RI's
probe does not trigger the defect. \code{RationalizeDenominator} is correct on
every numeric denominator (R01--R10); R05 confirms \RII's prediction that a
numerator named \code{x} is erased (result 0), and R11 shows the helper
emitting three messages on a symbolic denominator. S01 shows the comparator
reversing a list; S02 the out-of-range slice; S04--S06 that
\code{PossibleZeroQ} correctly rejects three very close non-zero differences,
so \U{3} is a matter of unreliable \emph{acceptance}, not of easily
demonstrable false positives.
\begin{longtable}{@{}p{0.055\textwidth}>{\ttfamily\footnotesize\raggedright\arraybackslash}p{0.33\textwidth}>{\ttfamily\footnotesize\raggedright\arraybackslash}p{0.33\textwidth}p{0.05\textwidth}p{0.05\textwidth}p{0.05\textwidth}@{}}
\caption{\texttt{Factorc} on symbolic and numeric inputs. ``Equal'' is exact equality with the input.}\label{tab:helpH}\\
\toprule ID & Input & Output & Equal & Depth & Time (s)\\\midrule\endfirsthead
\toprule ID & Input & Output & Equal & Depth & Time (s)\\\midrule\endhead
\bottomrule\endfoot
H01 & Factorc[\allowbreak{}Sqrt[\allowbreak{}z\textasciicircum{}\allowbreak{}2]\allowbreak{}]\allowbreak{} & z & -- & 0$\to$0 & 0.000\\
H02 & Factorc[\allowbreak{}Sqrt[\allowbreak{}t +\allowbreak{} t\textasciicircum{}\allowbreak{}2]\allowbreak{}]\allowbreak{} & Sqrt[\allowbreak{}t*\allowbreak{}(\allowbreak{}1 +\allowbreak{} t)\allowbreak{}]\allowbreak{} & -- & 1$\to$1 & 0.000\\
H03 & Factorc[\allowbreak{}Sqrt[\allowbreak{}-x -\allowbreak{} y]\allowbreak{}]\allowbreak{} & I*\allowbreak{}Sqrt[\allowbreak{}x +\allowbreak{} y]\allowbreak{} & -- & 1$\to$1 & 0.001\\
H04 & Factorc[\allowbreak{}Sqrt[\allowbreak{}-2 -\allowbreak{} 2*\allowbreak{}Sqrt[\allowbreak{}2]\allowbreak{}]\allowbreak{}]\allowbreak{} & I*\allowbreak{}Sqrt[\allowbreak{}2*\allowbreak{}(\allowbreak{}1 +\allowbreak{} Sqrt[\allowbreak{}2]\allowbreak{})\allowbreak{}]\allowbreak{} & yes & 2$\to$2 & 0.001\\
H05 & Factorc[\allowbreak{}Sqrt[\allowbreak{}-(\allowbreak{}Sqrt[\allowbreak{}2]\allowbreak{} +\allowbreak{} (\allowbreak{}-1)\allowbreak{}\textasciicircum{}\allowbreak{}(\allowbreak{}1/\allowbreak{}3)\allowbreak{})\allowbreak{}]\allowbreak{}]\allowbreak{} & I*\allowbreak{}Sqrt[\allowbreak{}(\allowbreak{}-1)\allowbreak{}\textasciicircum{}\allowbreak{}(\allowbreak{}1/\allowbreak{}3)\allowbreak{} +\allowbreak{} Sqrt[\allowbreak{}2]\allowbreak{}]\allowbreak{} & \textbf{NO} & 2$\to$2 & 0.001\\
H06 & Factorc[\allowbreak{}Sqrt[\allowbreak{}(\allowbreak{}1 +\allowbreak{} I)\allowbreak{}*\allowbreak{}(\allowbreak{}3 +\allowbreak{} 2*\allowbreak{}Sqrt[\allowbreak{}2]\allowbreak{})\allowbreak{}]\allowbreak{}]\allowbreak{} & Sqrt[\allowbreak{}(\allowbreak{}1 +\allowbreak{} I)\allowbreak{}*\allowbreak{}(\allowbreak{}3 +\allowbreak{} 2*\allowbreak{}Sqrt[\allowbreak{}2]\allowbreak{})\allowbreak{}]\allowbreak{} & yes & 2$\to$2 & 0.001\\
H07 & Factorc[\allowbreak{}Sqrt[\allowbreak{}-1 -\allowbreak{} I]\allowbreak{}]\allowbreak{} & Sqrt[\allowbreak{}-1 -\allowbreak{} I]\allowbreak{} & yes & 1$\to$1 & 0.000\\
H08 & Factorc[\allowbreak{}Sqrt[\allowbreak{}-2 -\allowbreak{} 2*\allowbreak{}I*\allowbreak{}Sqrt[\allowbreak{}3]\allowbreak{}]\allowbreak{}]\allowbreak{} & Sqrt[\allowbreak{}2*\allowbreak{}(\allowbreak{}-1 -\allowbreak{} I*\allowbreak{}Sqrt[\allowbreak{}3]\allowbreak{})\allowbreak{}]\allowbreak{} & yes & 2$\to$2 & 0.001\\
H09 & Factorc[\allowbreak{}Sqrt[\allowbreak{}(\allowbreak{}-1 -\allowbreak{} I*\allowbreak{}Sqrt[\allowbreak{}3]\allowbreak{})\allowbreak{}*\allowbreak{}(\allowbreak{}1 +\allowbreak{} Sqrt[\allowbreak{}2]\allowbreak{})\allowbreak{}]\allowbreak{}]\allowbreak{} & Sqrt[\allowbreak{}(\allowbreak{}-I)\allowbreak{}*\allowbreak{}(\allowbreak{}1 +\allowbreak{} Sqrt[\allowbreak{}2]\allowbreak{})\allowbreak{}*\allowbreak{}(\allowbreak{}-I +\allowbreak{} Sqrt[\allowbreak{}3]\allowbreak{})\allowbreak{}]\allowbreak{} & yes & 2$\to$2 & 0.001\\
H10 & Factorc[\allowbreak{}(\allowbreak{}-3 -\allowbreak{} 2*\allowbreak{}Sqrt[\allowbreak{}2]\allowbreak{})\allowbreak{}\textasciicircum{}\allowbreak{}(\allowbreak{}1/\allowbreak{}3)\allowbreak{}]\allowbreak{} & (\allowbreak{}-3 -\allowbreak{} 2*\allowbreak{}Sqrt[\allowbreak{}2]\allowbreak{})\allowbreak{}\textasciicircum{}\allowbreak{}(\allowbreak{}1/\allowbreak{}3)\allowbreak{} & yes & 2$\to$2 & 0.000\\
H11 & Factorc[\allowbreak{}(\allowbreak{}(\allowbreak{}-1 +\allowbreak{} I)\allowbreak{}*\allowbreak{}(\allowbreak{}3 +\allowbreak{} 2*\allowbreak{}Sqrt[\allowbreak{}2]\allowbreak{})\allowbreak{})\allowbreak{}\textasciicircum{}\allowbreak{}(\allowbreak{}1/\allowbreak{}3)\allowbreak{}]\allowbreak{} & (\allowbreak{}(\allowbreak{}-1 +\allowbreak{} I)\allowbreak{}*\allowbreak{}(\allowbreak{}3 +\allowbreak{} 2*\allowbreak{}Sqrt[\allowbreak{}2]\allowbreak{})\allowbreak{})\allowbreak{}\textasciicircum{}\allowbreak{}(\allowbreak{}1/\allowbreak{}3)\allowbreak{} & yes & 2$\to$2 & 0.001\\
H12 & Factorc[\allowbreak{}Sqrt[\allowbreak{}2]\allowbreak{} +\allowbreak{} Sqrt[\allowbreak{}3]\allowbreak{}]\allowbreak{} & Sqrt[\allowbreak{}2]\allowbreak{} +\allowbreak{} Sqrt[\allowbreak{}3]\allowbreak{} & yes & 1$\to$1 & 0.001\\
H13 & Factorc[\allowbreak{}(\allowbreak{}1 +\allowbreak{} Sqrt[\allowbreak{}2]\allowbreak{})\allowbreak{}\textasciicircum{}\allowbreak{}2]\allowbreak{} & (\allowbreak{}1 +\allowbreak{} Sqrt[\allowbreak{}2]\allowbreak{})\allowbreak{}\textasciicircum{}\allowbreak{}2 & yes & 1$\to$1 & 0.000\\
H14 & Factorc[\allowbreak{}Sqrt[\allowbreak{}8 +\allowbreak{} 4*\allowbreak{}Sqrt[\allowbreak{}3]\allowbreak{}]\allowbreak{}]\allowbreak{} & 2*\allowbreak{}Sqrt[\allowbreak{}2 +\allowbreak{} Sqrt[\allowbreak{}3]\allowbreak{}]\allowbreak{} & yes & 2$\to$2 & 0.001\\
H15 & Factorc[\allowbreak{}Sqrt[\allowbreak{}-8 -\allowbreak{} 4*\allowbreak{}Sqrt[\allowbreak{}3]\allowbreak{}]\allowbreak{}]\allowbreak{} & (\allowbreak{}2*\allowbreak{}I)\allowbreak{}*\allowbreak{}Sqrt[\allowbreak{}2 +\allowbreak{} Sqrt[\allowbreak{}3]\allowbreak{}]\allowbreak{} & yes & 2$\to$2 & 0.001\\
H16 & Factorc[\allowbreak{}Sqrt[\allowbreak{}(\allowbreak{}-1)\allowbreak{}\textasciicircum{}\allowbreak{}(\allowbreak{}1/\allowbreak{}3)\allowbreak{}*\allowbreak{}(\allowbreak{}3 +\allowbreak{} 2*\allowbreak{}Sqrt[\allowbreak{}2]\allowbreak{})\allowbreak{}]\allowbreak{}]\allowbreak{} & (\allowbreak{}-1)\allowbreak{}\textasciicircum{}\allowbreak{}(\allowbreak{}1/\allowbreak{}6)\allowbreak{}*\allowbreak{}Sqrt[\allowbreak{}3 +\allowbreak{} 2*\allowbreak{}Sqrt[\allowbreak{}2]\allowbreak{}]\allowbreak{} & yes & 2$\to$2 & 0.001\\
H17 & Factorc[\allowbreak{}Sqrt[\allowbreak{}(\allowbreak{}-1)\allowbreak{}\textasciicircum{}\allowbreak{}(\allowbreak{}2/\allowbreak{}3)\allowbreak{}*\allowbreak{}(\allowbreak{}3 +\allowbreak{} 2*\allowbreak{}Sqrt[\allowbreak{}2]\allowbreak{})\allowbreak{}]\allowbreak{}]\allowbreak{} & (\allowbreak{}-1)\allowbreak{}\textasciicircum{}\allowbreak{}(\allowbreak{}1/\allowbreak{}3)\allowbreak{}*\allowbreak{}Sqrt[\allowbreak{}3 +\allowbreak{} 2*\allowbreak{}Sqrt[\allowbreak{}2]\allowbreak{}]\allowbreak{} & yes & 2$\to$2 & 0.001\\
H18 & Factorc[\allowbreak{}Sqrt[\allowbreak{}x\textasciicircum{}\allowbreak{}4]\allowbreak{}]\allowbreak{} & x\textasciicircum{}\allowbreak{}2 & -- & 0$\to$0 & 0.000\\
H19 & Factorc[\allowbreak{}(\allowbreak{}x\textasciicircum{}\allowbreak{}2)\allowbreak{}\textasciicircum{}\allowbreak{}(\allowbreak{}1/\allowbreak{}3)\allowbreak{}]\allowbreak{} & x\textasciicircum{}\allowbreak{}(\allowbreak{}2/\allowbreak{}3)\allowbreak{} & -- & 1$\to$1 & 0.000\\
H20 & Factorc[\allowbreak{}Sqrt[\allowbreak{}4*\allowbreak{}x\textasciicircum{}\allowbreak{}2 +\allowbreak{} 4*\allowbreak{}x\textasciicircum{}\allowbreak{}3]\allowbreak{}]\allowbreak{} & 2*\allowbreak{}Sqrt[\allowbreak{}x\textasciicircum{}\allowbreak{}2*\allowbreak{}(\allowbreak{}1 +\allowbreak{} x)\allowbreak{}]\allowbreak{} & -- & 1$\to$1 & 0.001\\
H21 & Factorc[\allowbreak{}Sqrt[\allowbreak{}-3 -\allowbreak{} 2*\allowbreak{}Sqrt[\allowbreak{}2]\allowbreak{}]\allowbreak{}]\allowbreak{} & I*\allowbreak{}Sqrt[\allowbreak{}3 +\allowbreak{} 2*\allowbreak{}Sqrt[\allowbreak{}2]\allowbreak{}]\allowbreak{} & yes & 2$\to$2 & 0.001\\
H22 & Factorc[\allowbreak{}Sqrt[\allowbreak{}(\allowbreak{}-1 +\allowbreak{} I*\allowbreak{}Sqrt[\allowbreak{}3]\allowbreak{})\allowbreak{}*\allowbreak{}(\allowbreak{}-1 +\allowbreak{} I*\allowbreak{}Sqrt[\allowbreak{}3]\allowbreak{})\allowbreak{}]\allowbreak{}]\allowbreak{} & 1 -\allowbreak{} I*\allowbreak{}Sqrt[\allowbreak{}3]\allowbreak{} & yes & 1$\to$1 & 0.000\\
H23 & Factorc[\allowbreak{}Sqrt[\allowbreak{}-(\allowbreak{}1 +\allowbreak{} Sqrt[\allowbreak{}2]\allowbreak{})\allowbreak{} -\allowbreak{} (\allowbreak{}1 +\allowbreak{} Sqrt[\allowbreak{}2]\allowbreak{})\allowbreak{}*\allowbreak{}I*\allowbreak{}Sqrt[\allowbreak{}3]\allowbreak{}]\allowbreak{}]\allowbreak{} & Sqrt[\allowbreak{}(\allowbreak{}-I)\allowbreak{}*\allowbreak{}(\allowbreak{}1 +\allowbreak{} Sqrt[\allowbreak{}2]\allowbreak{})\allowbreak{}*\allowbreak{}(\allowbreak{}-I +\allowbreak{} Sqrt[\allowbreak{}3]\allowbreak{})\allowbreak{}]\allowbreak{} & yes & 2$\to$2 & 0.001\\
\end{longtable}
\begin{longtable}{@{}p{0.055\textwidth}>{\ttfamily\footnotesize\raggedright\arraybackslash}p{0.33\textwidth}>{\ttfamily\footnotesize\raggedright\arraybackslash}p{0.33\textwidth}p{0.05\textwidth}p{0.05\textwidth}p{0.05\textwidth}@{}}
\caption{Helper-level probes: \texttt{RationalizeDenominator}, marker collisions, nesting helpers, comparator, slice, \texttt{PossibleZeroQ}.}\label{tab:helpR}\\
\toprule ID & Input & Output & Equal & Depth & Time (s)\\\midrule\endfirsthead
\toprule ID & Input & Output & Equal & Depth & Time (s)\\\midrule\endhead
\bottomrule\endfoot
R01 & RationalizeDenominator[\allowbreak{}1/\allowbreak{}(\allowbreak{}1 +\allowbreak{} Sqrt[\allowbreak{}2]\allowbreak{})\allowbreak{}]\allowbreak{} & -1 +\allowbreak{} Sqrt[\allowbreak{}2]\allowbreak{} & yes & 1$\to$1 & 0.000\\
R02 & RationalizeDenominator[\allowbreak{}1/\allowbreak{}(\allowbreak{}Sqrt[\allowbreak{}2]\allowbreak{} +\allowbreak{} Sqrt[\allowbreak{}3]\allowbreak{})\allowbreak{}]\allowbreak{} & -Sqrt[\allowbreak{}2]\allowbreak{} +\allowbreak{} Sqrt[\allowbreak{}3]\allowbreak{} & yes & 1$\to$1 & 0.002\\
R03 & RationalizeDenominator[\allowbreak{}1/\allowbreak{}(\allowbreak{}1 +\allowbreak{} I*\allowbreak{}Sqrt[\allowbreak{}2]\allowbreak{})\allowbreak{}]\allowbreak{} & (\allowbreak{}1 -\allowbreak{} I*\allowbreak{}Sqrt[\allowbreak{}2]\allowbreak{})\allowbreak{}/\allowbreak{}3 & yes & 1$\to$1 & 0.001\\
R04 & RationalizeDenominator[\allowbreak{}(\allowbreak{}1 +\allowbreak{} Sqrt[\allowbreak{}2]\allowbreak{})\allowbreak{}/\allowbreak{}(\allowbreak{}Sqrt[\allowbreak{}2]\allowbreak{} -\allowbreak{} Sqrt[\allowbreak{}3]\allowbreak{})\allowbreak{}]\allowbreak{} & -2 -\allowbreak{} Sqrt[\allowbreak{}2]\allowbreak{} -\allowbreak{} Sqrt[\allowbreak{}3]\allowbreak{} -\allowbreak{} Sqrt[\allowbreak{}6]\allowbreak{} & yes & 1$\to$1 & 0.002\\
R05 & RationalizeDenominator[\allowbreak{}x/\allowbreak{}(\allowbreak{}1 +\allowbreak{} Sqrt[\allowbreak{}2]\allowbreak{})\allowbreak{}]\allowbreak{} & 0 & yes & 0$\to$0 & 0.000\\
R06 & RationalizeDenominator[\allowbreak{}1/\allowbreak{}2 +\allowbreak{} 1/\allowbreak{}Sqrt[\allowbreak{}2]\allowbreak{}]\allowbreak{} & 1/\allowbreak{}2 +\allowbreak{} 1/\allowbreak{}Sqrt[\allowbreak{}2]\allowbreak{} & yes & 1$\to$1 & 0.000\\
R07 & RationalizeDenominator[\allowbreak{}1/\allowbreak{}Sqrt[\allowbreak{}2]\allowbreak{}]\allowbreak{} & 1/\allowbreak{}Sqrt[\allowbreak{}2]\allowbreak{} & yes & 1$\to$1 & 0.000\\
R08 & RationalizeDenominator[\allowbreak{}1/\allowbreak{}(\allowbreak{}1 +\allowbreak{} 2\textasciicircum{}\allowbreak{}(\allowbreak{}1/\allowbreak{}3)\allowbreak{})\allowbreak{}]\allowbreak{} & (\allowbreak{}1 -\allowbreak{} 2\textasciicircum{}\allowbreak{}(\allowbreak{}1/\allowbreak{}3)\allowbreak{} +\allowbreak{} 2\textasciicircum{}\allowbreak{}(\allowbreak{}2/\allowbreak{}3)\allowbreak{})\allowbreak{}/\allowbreak{}3 & yes & 1$\to$1 & 0.001\\
R09 & RationalizeDenominator[\allowbreak{}1/\allowbreak{}Sqrt[\allowbreak{}1 +\allowbreak{} Sqrt[\allowbreak{}2]\allowbreak{}]\allowbreak{}]\allowbreak{} & -Sqrt[\allowbreak{}1 +\allowbreak{} Sqrt[\allowbreak{}2]\allowbreak{}]\allowbreak{} +\allowbreak{} Sqrt[\allowbreak{}2*\allowbreak{}(\allowbreak{}1 +\allowbreak{} Sqrt[\allowbreak{}2]\allowbreak{})\allowbreak{}]\allowbreak{} & yes & 2$\to$2 & 0.001\\
R10 & RationalizeDenominator[\allowbreak{}0]\allowbreak{} & 0 & yes & 0$\to$0 & 0.000\\
R11 & RationalizeDenominator[\allowbreak{}1/\allowbreak{}x]\allowbreak{} & -(\allowbreak{}PolynomialQuotient[\allowbreak{}0,\allowbreak{} 0,\allowbreak{} 0]\allowbreak{}/\allowbreak{}MinimalPolynomial[\allowbreak{}x]\allowbreak{}[\allowbreak{}0]\allowbreak{})\allowbreak{}\newline{\scriptsize\itshape MinimalPolynomial::\allowbreak{}nalg, General::\allowbreak{}stop, PolynomialQuotient::\allowbreak{}ivar} & -- & 0$\to$0 & 0.008\\
M01 & DenestRadicals3[\allowbreak{}f[\allowbreak{}1]\allowbreak{},\allowbreak{} False,\allowbreak{} Identity]\allowbreak{} & 1 & -- & 0$\to$0 & 0.000\\
M02 & DenestRadicals3[\allowbreak{}g[\allowbreak{}1]\allowbreak{},\allowbreak{} False,\allowbreak{} Identity]\allowbreak{} & g[\allowbreak{}1]\allowbreak{} & -- & 0$\to$0 & 0.000\\
M03 & DenestRadicals3[\allowbreak{}f[\allowbreak{}Sqrt[\allowbreak{}3 +\allowbreak{} 2*\allowbreak{}Sqrt[\allowbreak{}2]\allowbreak{}]\allowbreak{}]\allowbreak{},\allowbreak{} False,\allowbreak{} Identity]\allowbreak{} & f[\allowbreak{}Sqrt[\allowbreak{}3 +\allowbreak{} 2*\allowbreak{}Sqrt[\allowbreak{}2]\allowbreak{}]\allowbreak{}]\allowbreak{}\newline{\scriptsize\itshape General::\allowbreak{}newsym} & -- & 2$\to$2 & 0.007\\
N01 & maxnestdepth[\allowbreak{}Sqrt[\allowbreak{}3 +\allowbreak{} 2*\allowbreak{}Sqrt[\allowbreak{}2]\allowbreak{}]\allowbreak{},\allowbreak{} MatchQ[\allowbreak{}\#1,\allowbreak{} \_\textasciicircum{}\allowbreak{}\_Rational]\allowbreak{} \& ]\allowbreak{} & \{\{Sqrt[\allowbreak{}3 +\allowbreak{} 2*\allowbreak{}Sqrt[\allowbreak{}2]\allowbreak{}]\allowbreak{},\allowbreak{} 2\}\} & -- & 2$\to$2 & 0.000\\
N02 & maxnestdepth[\allowbreak{}(\allowbreak{}1 +\allowbreak{} 2\textasciicircum{}\allowbreak{}(\allowbreak{}1/\allowbreak{}3)\allowbreak{})\allowbreak{}\textasciicircum{}\allowbreak{}3,\allowbreak{} MatchQ[\allowbreak{}\#1,\allowbreak{} \_\textasciicircum{}\allowbreak{}\_Rational]\allowbreak{} \& ]\allowbreak{} & \{\{2\textasciicircum{}\allowbreak{}(\allowbreak{}1/\allowbreak{}3)\allowbreak{},\allowbreak{} 1\}\} & -- & 1$\to$1 & 0.000\\
N03 & LCM @@ (\allowbreak{}maxnestdepth[\allowbreak{}(\allowbreak{}1 +\allowbreak{} 2\textasciicircum{}\allowbreak{}(\allowbreak{}1/\allowbreak{}3)\allowbreak{})\allowbreak{}\textasciicircum{}\allowbreak{}3,\allowbreak{} MatchQ[\allowbreak{}\#1,\allowbreak{} \_\textasciicircum{}\allowbreak{}\_Rational]\allowbreak{} \& ]\allowbreak{} /\allowbreak{}. \{(\allowbreak{}base\_)\allowbreak{}\textasciicircum{}\allowbreak{}(\allowbreak{}exp\_)\allowbreak{},\allowbreak{} \_\} :> Denominator[\allowbreak{}exp]\allowbreak{})\allowbreak{} & 3 & yes & 0$\to$0 & 0.000\\
N04 & maxnestdepth[\allowbreak{}Sqrt[\allowbreak{}1 +\allowbreak{} 2\textasciicircum{}\allowbreak{}(\allowbreak{}1/\allowbreak{}3)\allowbreak{}]\allowbreak{}*\allowbreak{}(\allowbreak{}1 +\allowbreak{} Sqrt[\allowbreak{}2]\allowbreak{})\allowbreak{}\textasciicircum{}\allowbreak{}(\allowbreak{}1/\allowbreak{}3)\allowbreak{},\allowbreak{} MatchQ[\allowbreak{}\#1,\allowbreak{} \_\textasciicircum{}\allowbreak{}\_Rational]\allowbreak{} \& ]\allowbreak{} & \{\{Sqrt[\allowbreak{}1 +\allowbreak{} 2\textasciicircum{}\allowbreak{}(\allowbreak{}1/\allowbreak{}3)\allowbreak{}]\allowbreak{},\allowbreak{} 2\},\allowbreak{} \{(\allowbreak{}1 +\allowbreak{} Sqrt[\allowbreak{}2]\allowbreak{})\allowbreak{}\textasciicircum{}\allowbreak{}(\allowbreak{}1/\allowbreak{}3)\allowbreak{},\allowbreak{} 2\}\} & -- & 2$\to$2 & 0.000\\
N05 & nestdepthsort[\allowbreak{}Sqrt[\allowbreak{}3 +\allowbreak{} 2*\allowbreak{}Sqrt[\allowbreak{}3 +\allowbreak{} 2*\allowbreak{}Sqrt[\allowbreak{}2]\allowbreak{}]\allowbreak{}]\allowbreak{},\allowbreak{} MatchQ[\allowbreak{}\#1,\allowbreak{} \_\textasciicircum{}\allowbreak{}\_Rational]\allowbreak{} \& ]\allowbreak{} & \{\{\{Sqrt[\allowbreak{}3 +\allowbreak{} 2*\allowbreak{}Sqrt[\allowbreak{}3 +\allowbreak{} 2*\allowbreak{}Sqrt[\allowbreak{}2]\allowbreak{}]\allowbreak{}]\allowbreak{},\allowbreak{} 3\}\},\allowbreak{} \{\{3 +\allowbreak{} 2*\allowbreak{}Sqrt[\allowbreak{}3 +\allowbreak{} 2*\allowbreak{}Sqrt[\allowbreak{}2]\allowbreak{}]\allowbreak{},\allowbreak{} 2\},\allowbreak{} \{2*\allowbreak{}Sqrt[\allowbreak{}3 +\allowbreak{} 2*\allowbreak{}Sqrt[\allowbreak{}2]\allowbreak{}]\allowbreak{},\allowbreak{} 2\},\allowbreak{} \{Sqrt[\allowbreak{}3 +\allowbreak{} 2*\allowbreak{}Sqrt[\allowbreak{}2]\allowbreak{}]\allowbreak{},\allowbreak{} 2\}\},\allowbreak{} \{\{3 +\allowbreak{} 2*\allowbreak{}Sqrt[\allowbreak{}2]\allowbreak{},\allowbreak{} 1\},\allowbreak{} \{2*\allowbreak{}Sqrt[\allowbreak{}2]\allowbreak{},\allowbreak{} 1\},\allowbreak{} \{Sqrt[\allowbreak{}2]\allowbreak{},\allowbreak{} 1\}\},\allowbreak{} \{\{3,\allowbreak{} 0\},\allowbreak{} \{2,\allowbreak{} 0\},\allowbreak{} \{3,\allowbreak{} 0\},\allowbreak{} \{2,\allowbreak{} 0\},\allowbreak{} \{2,\allowbreak{} 0\},\allowbreak{} \{1/\allowbreak{}2,\allowbreak{} 0\},\allowbreak{} \{1/\allowbreak{}2,\allowbreak{} 0\},\allowbreak{} \{1/\allowbreak{}2,\allowbreak{} 0\}\}\} & -- & 3$\to$3 & 0.000\\
N06 & maxnestdepth[\allowbreak{}Sqrt[\allowbreak{}3 +\allowbreak{} 2*\allowbreak{}Sqrt[\allowbreak{}2]\allowbreak{}]\allowbreak{} +\allowbreak{} (\allowbreak{}5 +\allowbreak{} 2*\allowbreak{}Sqrt[\allowbreak{}6]\allowbreak{})\allowbreak{}\textasciicircum{}\allowbreak{}(\allowbreak{}1/\allowbreak{}3)\allowbreak{},\allowbreak{} MatchQ[\allowbreak{}\#1,\allowbreak{} \_\textasciicircum{}\allowbreak{}\_Rational]\allowbreak{} \& ]\allowbreak{} & \{\{Sqrt[\allowbreak{}3 +\allowbreak{} 2*\allowbreak{}Sqrt[\allowbreak{}2]\allowbreak{}]\allowbreak{},\allowbreak{} 2\},\allowbreak{} \{(\allowbreak{}5 +\allowbreak{} 2*\allowbreak{}Sqrt[\allowbreak{}6]\allowbreak{})\allowbreak{}\textasciicircum{}\allowbreak{}(\allowbreak{}1/\allowbreak{}3)\allowbreak{},\allowbreak{} 2\}\} & -- & 2$\to$2 & 0.000\\
S01 & Sort[\allowbreak{}\{\{2,\allowbreak{} 1\},\allowbreak{} \{3,\allowbreak{} 1\},\allowbreak{} \{5,\allowbreak{} 1\},\allowbreak{} \{7,\allowbreak{} 1\},\allowbreak{} \{11,\allowbreak{} 1\},\allowbreak{} \{13,\allowbreak{} 1\},\allowbreak{} \{2003,\allowbreak{} 1\}\},\allowbreak{} \#1[\allowbreak{}[\allowbreak{}1]\allowbreak{}]\allowbreak{} <= \#1[\allowbreak{}[\allowbreak{}2]\allowbreak{}]\allowbreak{} \& ]\allowbreak{} & \{\{2003,\allowbreak{} 1\},\allowbreak{} \{13,\allowbreak{} 1\},\allowbreak{} \{11,\allowbreak{} 1\},\allowbreak{} \{7,\allowbreak{} 1\},\allowbreak{} \{5,\allowbreak{} 1\},\allowbreak{} \{3,\allowbreak{} 1\},\allowbreak{} \{2,\allowbreak{} 1\}\} & -- & 0$\to$0 & 0.000\\
S02 & Module[\allowbreak{}\{p = \{2,\allowbreak{} 3,\allowbreak{} 5,\allowbreak{} 7,\allowbreak{} 11,\allowbreak{} 0\}\},\allowbreak{} DeleteCases[\allowbreak{}p,\allowbreak{} 0]\allowbreak{}[\allowbreak{}[\allowbreak{}1 ;; Min[\allowbreak{}10,\allowbreak{} Length[\allowbreak{}p]\allowbreak{}]\allowbreak{}]\allowbreak{}]\allowbreak{}]\allowbreak{} & \{2,\allowbreak{} 3,\allowbreak{} 5,\allowbreak{} 7,\allowbreak{} 11\}[\allowbreak{}[\allowbreak{}1 ;; 6]\allowbreak{}]\allowbreak{}\newline{\scriptsize\itshape Part::\allowbreak{}take} & -- & 0$\to$0 & 0.000\\
S03 & ComplexitySort[\allowbreak{}\{6,\allowbreak{} Sqrt[\allowbreak{}2]\allowbreak{},\allowbreak{} 2,\allowbreak{} 3,\allowbreak{} 1 +\allowbreak{} Sqrt[\allowbreak{}2]\allowbreak{},\allowbreak{} 10,\allowbreak{} -3\}]\allowbreak{} & \{2,\allowbreak{} 3,\allowbreak{} -3,\allowbreak{} 6,\allowbreak{} 10,\allowbreak{} Sqrt[\allowbreak{}2]\allowbreak{},\allowbreak{} 1 +\allowbreak{} Sqrt[\allowbreak{}2]\allowbreak{}\} & -- & 1$\to$1 & 0.000\\
S04 & PossibleZeroQ[\allowbreak{}Sqrt[\allowbreak{}2]\allowbreak{} -\allowbreak{} FromContinuedFraction[\allowbreak{}ContinuedFraction[\allowbreak{}Sqrt[\allowbreak{}2]\allowbreak{},\allowbreak{} 400]\allowbreak{}]\allowbreak{}]\allowbreak{} & False & -- & 0$\to$0 & 0.000\\
S05 & PossibleZeroQ[\allowbreak{}Sqrt[\allowbreak{}2]\allowbreak{} -\allowbreak{} FromContinuedFraction[\allowbreak{}ContinuedFraction[\allowbreak{}Sqrt[\allowbreak{}2]\allowbreak{},\allowbreak{} 400]\allowbreak{}]\allowbreak{},\allowbreak{} Method -> "ExactAlgebraics"]\allowbreak{} & False\newline{\scriptsize\itshape N::\allowbreak{}meprec} & -- & 0$\to$0 & 0.001\\
S06 & PossibleZeroQ[\allowbreak{}Sqrt[\allowbreak{}2]\allowbreak{} -\allowbreak{} Root[\allowbreak{}\#1\textasciicircum{}\allowbreak{}2 -\allowbreak{} 2 +\allowbreak{} 10\textasciicircum{}\allowbreak{}(\allowbreak{}-100)\allowbreak{} \& ,\allowbreak{} 2]\allowbreak{}]\allowbreak{} & False & -- & 0$\to$0 & 0.001\\
S07 & DeleteDuplicates[\allowbreak{}\{1 +\allowbreak{} Sqrt[\allowbreak{}2]\allowbreak{},\allowbreak{} Sqrt[\allowbreak{}3 +\allowbreak{} 2*\allowbreak{}Sqrt[\allowbreak{}2]\allowbreak{}]\allowbreak{},\allowbreak{} 1 +\allowbreak{} Sqrt[\allowbreak{}2]\allowbreak{}\}]\allowbreak{} & \{1 +\allowbreak{} Sqrt[\allowbreak{}2]\allowbreak{},\allowbreak{} Sqrt[\allowbreak{}3 +\allowbreak{} 2*\allowbreak{}Sqrt[\allowbreak{}2]\allowbreak{}]\allowbreak{}\} & -- & 2$\to$2 & 0.000\\
\end{longtable}

\paragraph{Wrapper-level numeric branch errors and odd inputs (Table~\ref{tab:r2JK}).}
J01 and J02 are the most important new findings: a purely numeric input,
$\sqrt{-(\sqrt2+(-1)^{1/3})}$, alone or inside a sum, is returned with the
wrong sign, silently. J06 is the cube-root analogue. K14 shows that a
\emph{list} of two radicals is turned into
$(3+2\sqrt2)^{(\sqrt2+\sqrt3)/2}$. K04 shows that the ugly BFHT output is a
fixed point of the program. K12: inexact input is returned as a machine number.
\begin{longtable}{@{}p{0.055\textwidth}>{\ttfamily\footnotesize\raggedright\arraybackslash}p{0.33\textwidth}>{\ttfamily\footnotesize\raggedright\arraybackslash}p{0.33\textwidth}p{0.05\textwidth}p{0.05\textwidth}p{0.05\textwidth}@{}}
\caption{Original program: wrapper-level numeric branch errors caused by \texttt{Factorc} (J) and miscellaneous inputs (K).}\label{tab:r2JK}\\
\toprule ID & Input & Output & Equal & Depth & Time (s)\\\midrule\endfirsthead
\toprule ID & Input & Output & Equal & Depth & Time (s)\\\midrule\endhead
\bottomrule\endfoot
J01 & Strad[\allowbreak{}Sqrt[\allowbreak{}-(\allowbreak{}Sqrt[\allowbreak{}2]\allowbreak{} +\allowbreak{} (\allowbreak{}-1)\allowbreak{}\textasciicircum{}\allowbreak{}(\allowbreak{}1/\allowbreak{}3)\allowbreak{})\allowbreak{}]\allowbreak{}]\allowbreak{} & I*\allowbreak{}Sqrt[\allowbreak{}(\allowbreak{}-1)\allowbreak{}\textasciicircum{}\allowbreak{}(\allowbreak{}1/\allowbreak{}3)\allowbreak{} +\allowbreak{} Sqrt[\allowbreak{}2]\allowbreak{}]\allowbreak{} & \textbf{NO} & 2$\to$2 & 0.097\\
J02 & Strad[\allowbreak{}Sqrt[\allowbreak{}-(\allowbreak{}Sqrt[\allowbreak{}2]\allowbreak{} +\allowbreak{} (\allowbreak{}-1)\allowbreak{}\textasciicircum{}\allowbreak{}(\allowbreak{}1/\allowbreak{}3)\allowbreak{})\allowbreak{}]\allowbreak{} +\allowbreak{} 1]\allowbreak{} & 1 +\allowbreak{} I*\allowbreak{}Sqrt[\allowbreak{}(\allowbreak{}-1)\allowbreak{}\textasciicircum{}\allowbreak{}(\allowbreak{}1/\allowbreak{}3)\allowbreak{} +\allowbreak{} Sqrt[\allowbreak{}2]\allowbreak{}]\allowbreak{} & \textbf{NO} & 2$\to$2 & 0.101\\
J03 & Factorc[\allowbreak{}Sqrt[\allowbreak{}(\allowbreak{}-(\allowbreak{}1 +\allowbreak{} Sqrt[\allowbreak{}2]\allowbreak{})\allowbreak{})\allowbreak{}*\allowbreak{}(\allowbreak{}(\allowbreak{}1 +\allowbreak{} I*\allowbreak{}Sqrt[\allowbreak{}3]\allowbreak{})\allowbreak{}/\allowbreak{}2)\allowbreak{}]\allowbreak{}]\allowbreak{} & Sqrt[\allowbreak{}(\allowbreak{}-1/\allowbreak{}2*\allowbreak{}I)\allowbreak{}*\allowbreak{}(\allowbreak{}1 +\allowbreak{} Sqrt[\allowbreak{}2]\allowbreak{})\allowbreak{}*\allowbreak{}(\allowbreak{}-I +\allowbreak{} Sqrt[\allowbreak{}3]\allowbreak{})\allowbreak{}]\allowbreak{} & yes & 2$\to$2 & 0.004\\
J04 & Strad[\allowbreak{}Sqrt[\allowbreak{}(\allowbreak{}-(\allowbreak{}1 +\allowbreak{} Sqrt[\allowbreak{}2]\allowbreak{})\allowbreak{})\allowbreak{}*\allowbreak{}(\allowbreak{}(\allowbreak{}1 +\allowbreak{} I*\allowbreak{}Sqrt[\allowbreak{}3]\allowbreak{})\allowbreak{}/\allowbreak{}2)\allowbreak{}]\allowbreak{}]\allowbreak{} & Sqrt[\allowbreak{}(\allowbreak{}-1/\allowbreak{}2*\allowbreak{}I)\allowbreak{}*\allowbreak{}(\allowbreak{}1 +\allowbreak{} Sqrt[\allowbreak{}2]\allowbreak{})\allowbreak{}*\allowbreak{}(\allowbreak{}-I +\allowbreak{} Sqrt[\allowbreak{}3]\allowbreak{})\allowbreak{}]\allowbreak{} & yes & 2$\to$2 & 0.313\\
J05 & Strad[\allowbreak{}Sqrt[\allowbreak{}-3 -\allowbreak{} 2*\allowbreak{}Sqrt[\allowbreak{}2]\allowbreak{} -\allowbreak{} I*\allowbreak{}(\allowbreak{}3 +\allowbreak{} 2*\allowbreak{}Sqrt[\allowbreak{}2]\allowbreak{})\allowbreak{}*\allowbreak{}Sqrt[\allowbreak{}3]\allowbreak{}]\allowbreak{}]\allowbreak{} & (\allowbreak{}(\allowbreak{}2 +\allowbreak{} Sqrt[\allowbreak{}2]\allowbreak{})\allowbreak{}*\allowbreak{}(\allowbreak{}1 -\allowbreak{} I*\allowbreak{}Sqrt[\allowbreak{}3]\allowbreak{})\allowbreak{})\allowbreak{}/\allowbreak{}2 & yes & 1$\to$1 & 0.026\\
J06 & Strad[\allowbreak{}(\allowbreak{}-(\allowbreak{}Sqrt[\allowbreak{}2]\allowbreak{} +\allowbreak{} (\allowbreak{}-1)\allowbreak{}\textasciicircum{}\allowbreak{}(\allowbreak{}1/\allowbreak{}3)\allowbreak{})\allowbreak{})\allowbreak{}\textasciicircum{}\allowbreak{}(\allowbreak{}1/\allowbreak{}3)\allowbreak{}]\allowbreak{} & (\allowbreak{}-1)\allowbreak{}\textasciicircum{}\allowbreak{}(\allowbreak{}1/\allowbreak{}3)\allowbreak{}*\allowbreak{}(\allowbreak{}(\allowbreak{}-1)\allowbreak{}\textasciicircum{}\allowbreak{}(\allowbreak{}1/\allowbreak{}3)\allowbreak{} +\allowbreak{} Sqrt[\allowbreak{}2]\allowbreak{})\allowbreak{}\textasciicircum{}\allowbreak{}(\allowbreak{}1/\allowbreak{}3)\allowbreak{} & \textbf{NO} & 2$\to$2 & 0.792\\
J07 & Strad[\allowbreak{}Sqrt[\allowbreak{}-Sqrt[\allowbreak{}2]\allowbreak{} -\allowbreak{} I*\allowbreak{}Sqrt[\allowbreak{}3]\allowbreak{}]\allowbreak{}]\allowbreak{} & Sqrt[\allowbreak{}-Sqrt[\allowbreak{}2]\allowbreak{} -\allowbreak{} I*\allowbreak{}Sqrt[\allowbreak{}3]\allowbreak{}]\allowbreak{} & yes & 2$\to$2 & 0.109\\
J08 & Strad[\allowbreak{}Sqrt[\allowbreak{}-1 -\allowbreak{} I*\allowbreak{}Sqrt[\allowbreak{}3]\allowbreak{} -\allowbreak{} Sqrt[\allowbreak{}2]\allowbreak{} -\allowbreak{} I*\allowbreak{}Sqrt[\allowbreak{}6]\allowbreak{}]\allowbreak{}]\allowbreak{} & Sqrt[\allowbreak{}(\allowbreak{}-I)\allowbreak{}*\allowbreak{}(\allowbreak{}1 +\allowbreak{} Sqrt[\allowbreak{}2]\allowbreak{})\allowbreak{}*\allowbreak{}(\allowbreak{}-I +\allowbreak{} Sqrt[\allowbreak{}3]\allowbreak{})\allowbreak{}]\allowbreak{} & yes & 2$\to$2 & 0.450\\
K01 & Strad[\allowbreak{}Sqrt[\allowbreak{}28\textasciicircum{}\allowbreak{}(\allowbreak{}1/\allowbreak{}3)\allowbreak{} -\allowbreak{} 3]\allowbreak{}]\allowbreak{} & (\allowbreak{}-1 -\allowbreak{} 2\textasciicircum{}\allowbreak{}(\allowbreak{}2/\allowbreak{}3)\allowbreak{}*\allowbreak{}7\textasciicircum{}\allowbreak{}(\allowbreak{}1/\allowbreak{}3)\allowbreak{} +\allowbreak{} 2\textasciicircum{}\allowbreak{}(\allowbreak{}1/\allowbreak{}3)\allowbreak{}*\allowbreak{}7\textasciicircum{}\allowbreak{}(\allowbreak{}2/\allowbreak{}3)\allowbreak{})\allowbreak{}/\allowbreak{}3 & yes & 1$\to$1 & 0.018\\
K02 & Strad[\allowbreak{}(\allowbreak{}1 +\allowbreak{} 2\textasciicircum{}\allowbreak{}(\allowbreak{}1/\allowbreak{}3)\allowbreak{})\allowbreak{}\textasciicircum{}\allowbreak{}3]\allowbreak{} & (\allowbreak{}1 +\allowbreak{} 2\textasciicircum{}\allowbreak{}(\allowbreak{}1/\allowbreak{}3)\allowbreak{})\allowbreak{}\textasciicircum{}\allowbreak{}3 & yes & 1$\to$1 & 0.002\\
K03 & denest11[\allowbreak{}(\allowbreak{}1 +\allowbreak{} 2\textasciicircum{}\allowbreak{}(\allowbreak{}1/\allowbreak{}3)\allowbreak{})\allowbreak{}\textasciicircum{}\allowbreak{}3]\allowbreak{} & 3*\allowbreak{}(\allowbreak{}1 +\allowbreak{} 2\textasciicircum{}\allowbreak{}(\allowbreak{}1/\allowbreak{}3)\allowbreak{} +\allowbreak{} 2\textasciicircum{}\allowbreak{}(\allowbreak{}2/\allowbreak{}3)\allowbreak{})\allowbreak{} & yes & 1$\to$1 & 0.006\\
K04 & Strad[\allowbreak{}Strad[\allowbreak{}Sqrt[\allowbreak{}16 -\allowbreak{} 2*\allowbreak{}Sqrt[\allowbreak{}29]\allowbreak{} +\allowbreak{} 2*\allowbreak{}Sqrt[\allowbreak{}55 -\allowbreak{} 10*\allowbreak{}Sqrt[\allowbreak{}29]\allowbreak{}]\allowbreak{}]\allowbreak{},\allowbreak{} True]\allowbreak{}]\allowbreak{} & (\allowbreak{}Sqrt[\allowbreak{}319 -\allowbreak{} 58*\allowbreak{}Sqrt[\allowbreak{}29]\allowbreak{}]\allowbreak{} -\allowbreak{} 10*\allowbreak{}Sqrt[\allowbreak{}11 -\allowbreak{} 2*\allowbreak{}Sqrt[\allowbreak{}29]\allowbreak{}]\allowbreak{} +\allowbreak{} Sqrt[\allowbreak{}5]\allowbreak{}*\allowbreak{}(\allowbreak{}-16 +\allowbreak{} 3*\allowbreak{}Sqrt[\allowbreak{}29]\allowbreak{})\allowbreak{})\allowbreak{}/\allowbreak{}(\allowbreak{}-5 +\allowbreak{} Sqrt[\allowbreak{}29]\allowbreak{} -\allowbreak{} Sqrt[\allowbreak{}55 -\allowbreak{} 10*\allowbreak{}Sqrt[\allowbreak{}29]\allowbreak{}]\allowbreak{})\allowbreak{}\newline{\scriptsize\itshape PossibleZeroQ::\allowbreak{}ztest1} & yes & 2$\to$2 & 0.370\\
K05 & Strad[\allowbreak{}Sqrt[\allowbreak{}2 +\allowbreak{} Sqrt[\allowbreak{}3]\allowbreak{}]\allowbreak{}*\allowbreak{}Sqrt[\allowbreak{}2 -\allowbreak{} Sqrt[\allowbreak{}3]\allowbreak{}]\allowbreak{}]\allowbreak{} & 1 & yes & 0$\to$0 & 0.001\\
K06 & Strad[\allowbreak{}Sqrt[\allowbreak{}Sqrt[\allowbreak{}2]\allowbreak{} -\allowbreak{} 1]\allowbreak{}*\allowbreak{}Sqrt[\allowbreak{}Sqrt[\allowbreak{}2]\allowbreak{} +\allowbreak{} 1]\allowbreak{}]\allowbreak{} & 1 & yes & 0$\to$0 & 0.000\\
K07 & Strad[\allowbreak{}Sqrt[\allowbreak{}3 +\allowbreak{} 2*\allowbreak{}Sqrt[\allowbreak{}2]\allowbreak{}]\allowbreak{}\textasciicircum{}\allowbreak{}(\allowbreak{}1/\allowbreak{}3)\allowbreak{}]\allowbreak{} & (\allowbreak{}1 +\allowbreak{} Sqrt[\allowbreak{}2]\allowbreak{})\allowbreak{}\textasciicircum{}\allowbreak{}(\allowbreak{}1/\allowbreak{}3)\allowbreak{} & yes & 2$\to$2 & 0.010\\
K08 & Strad[\allowbreak{}(\allowbreak{}3 +\allowbreak{} 2*\allowbreak{}Sqrt[\allowbreak{}2]\allowbreak{})\allowbreak{}\textasciicircum{}\allowbreak{}(\allowbreak{}1/\allowbreak{}6)\allowbreak{}]\allowbreak{} & (\allowbreak{}1 +\allowbreak{} Sqrt[\allowbreak{}2]\allowbreak{})\allowbreak{}\textasciicircum{}\allowbreak{}(\allowbreak{}1/\allowbreak{}3)\allowbreak{} & yes & 2$\to$2 & 0.012\\
K09 & Strad[\allowbreak{}Sqrt[\allowbreak{}5 +\allowbreak{} 2*\allowbreak{}Sqrt[\allowbreak{}6]\allowbreak{}]\allowbreak{}/\allowbreak{}Sqrt[\allowbreak{}3 +\allowbreak{} 2*\allowbreak{}Sqrt[\allowbreak{}2]\allowbreak{}]\allowbreak{}]\allowbreak{} & 2 -\allowbreak{} Sqrt[\allowbreak{}2]\allowbreak{} -\allowbreak{} Sqrt[\allowbreak{}3]\allowbreak{} +\allowbreak{} Sqrt[\allowbreak{}6]\allowbreak{} & yes & 1$\to$1 & 0.024\\
K10 & Strad[\allowbreak{}Exp[\allowbreak{}I*\allowbreak{}(\allowbreak{}Pi/\allowbreak{}7)\allowbreak{}]\allowbreak{}*\allowbreak{}Sqrt[\allowbreak{}3 +\allowbreak{} 2*\allowbreak{}Sqrt[\allowbreak{}2]\allowbreak{}]\allowbreak{}]\allowbreak{} & (\allowbreak{}-1)\allowbreak{}\textasciicircum{}\allowbreak{}(\allowbreak{}1/\allowbreak{}7)\allowbreak{}*\allowbreak{}(\allowbreak{}1 +\allowbreak{} Sqrt[\allowbreak{}2]\allowbreak{})\allowbreak{} & yes & 1$\to$1 & 0.010\\
K11 & Strad[\allowbreak{}Sqrt[\allowbreak{}3 +\allowbreak{} 2*\allowbreak{}Sqrt[\allowbreak{}2]\allowbreak{}]\allowbreak{} +\allowbreak{} Pi]\allowbreak{} & 1 +\allowbreak{} Sqrt[\allowbreak{}2]\allowbreak{} +\allowbreak{} Pi & yes & 1$\to$1 & 0.010\\
K12 & Strad[\allowbreak{}Sqrt[\allowbreak{}3. +\allowbreak{} 2*\allowbreak{}Sqrt[\allowbreak{}2]\allowbreak{}]\allowbreak{}]\allowbreak{} & 2.414213562373095 & yes & 0$\to$0 & 0.000\\
K13 & Strad[\allowbreak{}Sqrt[\allowbreak{}2 +\allowbreak{} Sqrt[\allowbreak{}3]\allowbreak{}]\allowbreak{},\allowbreak{} "yes"]\allowbreak{} & (\allowbreak{}1 +\allowbreak{} Sqrt[\allowbreak{}3]\allowbreak{})\allowbreak{}/\allowbreak{}Sqrt[\allowbreak{}2]\allowbreak{} & yes & 1$\to$1 & 0.013\\
K14 & Strad[\allowbreak{}\{Sqrt[\allowbreak{}3 +\allowbreak{} 2*\allowbreak{}Sqrt[\allowbreak{}2]\allowbreak{}]\allowbreak{},\allowbreak{} Sqrt[\allowbreak{}5 +\allowbreak{} 2*\allowbreak{}Sqrt[\allowbreak{}6]\allowbreak{}]\allowbreak{}\}]\allowbreak{} & (\allowbreak{}3 +\allowbreak{} 2*\allowbreak{}Sqrt[\allowbreak{}2]\allowbreak{})\allowbreak{}\textasciicircum{}\allowbreak{}(\allowbreak{}(\allowbreak{}Sqrt[\allowbreak{}2]\allowbreak{} +\allowbreak{} Sqrt[\allowbreak{}3]\allowbreak{})\allowbreak{}/\allowbreak{}2)\allowbreak{} & yes & 1$\to$1 & 0.019\\
K15 & Strad[\allowbreak{}Sqrt[\allowbreak{}3 +\allowbreak{} 2*\allowbreak{}Sqrt[\allowbreak{}2]\allowbreak{}]\allowbreak{} == 1 +\allowbreak{} Sqrt[\allowbreak{}2]\allowbreak{}]\allowbreak{} & True & -- & 0$\to$0 & 0.002\\
K16 & Strad[\allowbreak{}Sqrt[\allowbreak{}9 +\allowbreak{} 4*\allowbreak{}Sqrt[\allowbreak{}5]\allowbreak{}]\allowbreak{}]\allowbreak{} & 2 +\allowbreak{} Sqrt[\allowbreak{}5]\allowbreak{} & yes & 1$\to$1 & 0.009\\
K17 & Strad[\allowbreak{}Sqrt[\allowbreak{}(\allowbreak{}3 +\allowbreak{} 2*\allowbreak{}Sqrt[\allowbreak{}2]\allowbreak{})\allowbreak{}/\allowbreak{}(\allowbreak{}5 +\allowbreak{} 2*\allowbreak{}Sqrt[\allowbreak{}6]\allowbreak{})\allowbreak{}]\allowbreak{}]\allowbreak{} & -2 -\allowbreak{} Sqrt[\allowbreak{}2]\allowbreak{} +\allowbreak{} Sqrt[\allowbreak{}3]\allowbreak{} +\allowbreak{} Sqrt[\allowbreak{}6]\allowbreak{} & yes & 1$\to$1 & 0.024\\
K18 & Strad[\allowbreak{}Sqrt[\allowbreak{}3 +\allowbreak{} 2*\allowbreak{}Sqrt[\allowbreak{}2]\allowbreak{}]\allowbreak{}*\allowbreak{}Sqrt[\allowbreak{}x]\allowbreak{}]\allowbreak{} & (\allowbreak{}1 +\allowbreak{} Sqrt[\allowbreak{}2]\allowbreak{})\allowbreak{}*\allowbreak{}Sqrt[\allowbreak{}x]\allowbreak{} & yes & 1$\to$1 & 0.008\\
K19 & Strad[\allowbreak{}Sqrt[\allowbreak{}Sqrt[\allowbreak{}2]\allowbreak{} +\allowbreak{} 1]\allowbreak{}\textasciicircum{}\allowbreak{}3]\allowbreak{} & (\allowbreak{}1 +\allowbreak{} Sqrt[\allowbreak{}2]\allowbreak{})\allowbreak{}\textasciicircum{}\allowbreak{}(\allowbreak{}3/\allowbreak{}2)\allowbreak{} & yes & 2$\to$2 & 0.028\\
K20 & Strad[\allowbreak{}Sqrt[\allowbreak{}2 +\allowbreak{} Sqrt[\allowbreak{}3]\allowbreak{}]\allowbreak{}*\allowbreak{}Sqrt[\allowbreak{}3 +\allowbreak{} 2*\allowbreak{}Sqrt[\allowbreak{}2]\allowbreak{}]\allowbreak{}]\allowbreak{} & (\allowbreak{}(\allowbreak{}2 +\allowbreak{} Sqrt[\allowbreak{}2]\allowbreak{})\allowbreak{}*\allowbreak{}(\allowbreak{}1 +\allowbreak{} Sqrt[\allowbreak{}3]\allowbreak{})\allowbreak{})\allowbreak{}/\allowbreak{}2 & yes & 1$\to$1 & 0.017\\
\end{longtable}

\paragraph{Higher-order roots, \code{PossibleZeroQ}, prime processing (Table~\ref{tab:r2LMN}).}
L04, a four-term square root, is denested to a correct but enormous quotient
(rationalizing its denominator and simplifying gives
$\sqrt2+\sqrt3+\sqrt5+\sqrt7$; the corrected version does this). L14, a tenth
root, takes 20\,s and emits \code{Ceiling::meprec} because
\code{Ceiling[Log[1000]/Log[10]]} does not evaluate (R01--R04 in
Table~\ref{tab:r3QRS}). N01--N04 compare the prime processing with the original
and the corrected comparator on the same factor lists: the original reverses
the list and never prunes; the corrected one sorts, prunes, and then hits the
stale-length \code{Part::take} (\U{6}). N05/N06 materialize the $2^{10}$ and
$10^5$ divisor lists that the ``cap'' permits.
\begin{longtable}{@{}p{0.055\textwidth}>{\ttfamily\footnotesize\raggedright\arraybackslash}p{0.33\textwidth}>{\ttfamily\footnotesize\raggedright\arraybackslash}p{0.33\textwidth}p{0.05\textwidth}p{0.05\textwidth}p{0.05\textwidth}@{}}
\caption{Original program: higher-order roots and timing (L), \texttt{PossibleZeroQ} on actual candidate differences (M), prime processing with the original and the corrected comparator (N).}\label{tab:r2LMN}\\
\toprule ID & Input & Output & Equal & Depth & Time (s)\\\midrule\endfirsthead
\toprule ID & Input & Output & Equal & Depth & Time (s)\\\midrule\endhead
\bottomrule\endfoot
L01 & Strad[\allowbreak{}(\allowbreak{}239 +\allowbreak{} 169*\allowbreak{}Sqrt[\allowbreak{}2]\allowbreak{})\allowbreak{}\textasciicircum{}\allowbreak{}(\allowbreak{}1/\allowbreak{}7)\allowbreak{}]\allowbreak{} & 1 +\allowbreak{} Sqrt[\allowbreak{}2]\allowbreak{} & yes & 1$\to$1 & 0.008\\
L02 & Strad[\allowbreak{}(\allowbreak{}2 +\allowbreak{} Sqrt[\allowbreak{}2]\allowbreak{})\allowbreak{}\textasciicircum{}\allowbreak{}(\allowbreak{}1/\allowbreak{}5)\allowbreak{}]\allowbreak{} & 2\textasciicircum{}\allowbreak{}(\allowbreak{}1/\allowbreak{}10)\allowbreak{}*\allowbreak{}(\allowbreak{}1 +\allowbreak{} Sqrt[\allowbreak{}2]\allowbreak{})\allowbreak{}\textasciicircum{}\allowbreak{}(\allowbreak{}1/\allowbreak{}5)\allowbreak{} & yes & 2$\to$2 & 0.198\\
L03 & Strad[\allowbreak{}(\allowbreak{}1 +\allowbreak{} Sqrt[\allowbreak{}2]\allowbreak{})\allowbreak{}\textasciicircum{}\allowbreak{}(\allowbreak{}1/\allowbreak{}7)\allowbreak{}]\allowbreak{} & (\allowbreak{}1 +\allowbreak{} Sqrt[\allowbreak{}2]\allowbreak{})\allowbreak{}\textasciicircum{}\allowbreak{}(\allowbreak{}1/\allowbreak{}7)\allowbreak{} & yes & 2$\to$2 & 0.494\\
L04 & Strad[\allowbreak{}Sqrt[\allowbreak{}17 +\allowbreak{} 2*\allowbreak{}Sqrt[\allowbreak{}6]\allowbreak{} +\allowbreak{} 2*\allowbreak{}Sqrt[\allowbreak{}10]\allowbreak{} +\allowbreak{} 2*\allowbreak{}Sqrt[\allowbreak{}14]\allowbreak{} +\allowbreak{} 2*\allowbreak{}Sqrt[\allowbreak{}15]\allowbreak{} +\allowbreak{} 2*\allowbreak{}Sqrt[\allowbreak{}21]\allowbreak{} +\allowbreak{} 2*\allowbreak{}Sqrt[\allowbreak{}35]\allowbreak{}]\allowbreak{}]\allowbreak{} & (\allowbreak{}6397*\allowbreak{}Sqrt[\allowbreak{}2]\allowbreak{} +\allowbreak{} 5316*\allowbreak{}Sqrt[\allowbreak{}3]\allowbreak{} +\allowbreak{} 4090*\allowbreak{}Sqrt[\allowbreak{}5]\allowbreak{} +\allowbreak{} 3432*\allowbreak{}Sqrt[\allowbreak{}7]\allowbreak{} +\allowbreak{} 1695*\allowbreak{}Sqrt[\allowbreak{}30]\allowbreak{} +\allowbreak{} 1423*\allowbreak{}Sqrt[\allowbreak{}42]\allowbreak{} +\allowbreak{} 1095*\allowbreak{}Sqrt[\allowbreak{}70]\allowbreak{} +\allowbreak{} 910*\allowbreak{}Sqrt[\allowbreak{}105]\allowbreak{})\allowbreak{}/\allowbreak{}(\allowbreak{}1108 +\allowbreak{} 467*\allowbreak{}Sqrt[\allowbreak{}6]\allowbreak{} +\allowbreak{} 359*\allowbreak{}Sqrt[\allowbreak{}10]\allowbreak{} +\allowbreak{} 299*\allowbreak{}Sqrt[\allowbreak{}14]\allowbreak{} +\allowbreak{} 302*\allowbreak{}Sqrt[\allowbreak{}15]\allowbreak{} +\allowbreak{} 252*\allowbreak{}Sqrt[\allowbreak{}21]\allowbreak{} +\allowbreak{} 194*\allowbreak{}Sqrt[\allowbreak{}35]\allowbreak{} +\allowbreak{} 81*\allowbreak{}Sqrt[\allowbreak{}210]\allowbreak{})\allowbreak{} & yes & 1$\to$1 & 0.956\\
L05 & Strad[\allowbreak{}Expand[\allowbreak{}(\allowbreak{}1 +\allowbreak{} Sqrt[\allowbreak{}2]\allowbreak{} +\allowbreak{} Sqrt[\allowbreak{}3]\allowbreak{})\allowbreak{}\textasciicircum{}\allowbreak{}3]\allowbreak{}\textasciicircum{}\allowbreak{}(\allowbreak{}1/\allowbreak{}3)\allowbreak{}]\allowbreak{} & 1 +\allowbreak{} Sqrt[\allowbreak{}2]\allowbreak{} +\allowbreak{} Sqrt[\allowbreak{}3]\allowbreak{} & yes & 1$\to$1 & 0.020\\
L06 & Strad[\allowbreak{}Expand[\allowbreak{}(\allowbreak{}Sqrt[\allowbreak{}2]\allowbreak{} +\allowbreak{} Sqrt[\allowbreak{}3]\allowbreak{})\allowbreak{}\textasciicircum{}\allowbreak{}4]\allowbreak{}\textasciicircum{}\allowbreak{}(\allowbreak{}1/\allowbreak{}4)\allowbreak{}]\allowbreak{} & Sqrt[\allowbreak{}5 +\allowbreak{} 2*\allowbreak{}Sqrt[\allowbreak{}6]\allowbreak{}]\allowbreak{} & yes & 2$\to$2 & 0.013\\
L07 & Strad[\allowbreak{}Sqrt[\allowbreak{}1000003 +\allowbreak{} 2*\allowbreak{}Sqrt[\allowbreak{}999983]\allowbreak{}]\allowbreak{}]\allowbreak{} & Sqrt[\allowbreak{}1000003 +\allowbreak{} 2*\allowbreak{}Sqrt[\allowbreak{}999983]\allowbreak{}]\allowbreak{} & yes & 2$\to$2 & 0.215\\
L08 & Strad[\allowbreak{}Sqrt[\allowbreak{}2 +\allowbreak{} Sqrt[\allowbreak{}3]\allowbreak{} +\allowbreak{} Sqrt[\allowbreak{}5]\allowbreak{}]\allowbreak{}]\allowbreak{} & Sqrt[\allowbreak{}2 +\allowbreak{} Sqrt[\allowbreak{}3]\allowbreak{} +\allowbreak{} Sqrt[\allowbreak{}5]\allowbreak{}]\allowbreak{} & yes & 2$\to$2 & 0.216\\
L09 & Strad[\allowbreak{}(\allowbreak{}5 +\allowbreak{} 2*\allowbreak{}Sqrt[\allowbreak{}6]\allowbreak{})\allowbreak{}\textasciicircum{}\allowbreak{}(\allowbreak{}1/\allowbreak{}6)\allowbreak{}]\allowbreak{} & (\allowbreak{}2 +\allowbreak{} Sqrt[\allowbreak{}6]\allowbreak{})\allowbreak{}\textasciicircum{}\allowbreak{}(\allowbreak{}1/\allowbreak{}3)\allowbreak{}/\allowbreak{}2\textasciicircum{}\allowbreak{}(\allowbreak{}1/\allowbreak{}6)\allowbreak{} & yes & 2$\to$2 & 2.256\\
L10 & Strad[\allowbreak{}(\allowbreak{}11*\allowbreak{}Sqrt[\allowbreak{}2]\allowbreak{} +\allowbreak{} 9*\allowbreak{}Sqrt[\allowbreak{}3]\allowbreak{})\allowbreak{}\textasciicircum{}\allowbreak{}(\allowbreak{}1/\allowbreak{}3)\allowbreak{}]\allowbreak{} & Sqrt[\allowbreak{}2]\allowbreak{} +\allowbreak{} Sqrt[\allowbreak{}3]\allowbreak{} & yes & 1$\to$1 & 0.015\\
L11 & Strad[\allowbreak{}Expand[\allowbreak{}(\allowbreak{}2\textasciicircum{}\allowbreak{}(\allowbreak{}1/\allowbreak{}3)\allowbreak{} +\allowbreak{} 3\textasciicircum{}\allowbreak{}(\allowbreak{}1/\allowbreak{}3)\allowbreak{})\allowbreak{}\textasciicircum{}\allowbreak{}3]\allowbreak{}\textasciicircum{}\allowbreak{}(\allowbreak{}1/\allowbreak{}3)\allowbreak{}]\allowbreak{} & 2\textasciicircum{}\allowbreak{}(\allowbreak{}1/\allowbreak{}3)\allowbreak{} +\allowbreak{} 3\textasciicircum{}\allowbreak{}(\allowbreak{}1/\allowbreak{}3)\allowbreak{} & yes & 1$\to$1 & 0.133\\
L12 & Strad[\allowbreak{}Sqrt[\allowbreak{}Expand[\allowbreak{}(\allowbreak{}1 +\allowbreak{} 2\textasciicircum{}\allowbreak{}(\allowbreak{}1/\allowbreak{}3)\allowbreak{} +\allowbreak{} 3\textasciicircum{}\allowbreak{}(\allowbreak{}1/\allowbreak{}3)\allowbreak{})\allowbreak{}\textasciicircum{}\allowbreak{}2]\allowbreak{}]\allowbreak{}]\allowbreak{} & 1 +\allowbreak{} 2\textasciicircum{}\allowbreak{}(\allowbreak{}1/\allowbreak{}3)\allowbreak{} +\allowbreak{} 3\textasciicircum{}\allowbreak{}(\allowbreak{}1/\allowbreak{}3)\allowbreak{} & yes & 1$\to$1 & 0.006\\
L13 & Strad[\allowbreak{}Expand[\allowbreak{}(\allowbreak{}1 +\allowbreak{} 2\textasciicircum{}\allowbreak{}(\allowbreak{}1/\allowbreak{}5)\allowbreak{})\allowbreak{}\textasciicircum{}\allowbreak{}5]\allowbreak{}\textasciicircum{}\allowbreak{}(\allowbreak{}1/\allowbreak{}5)\allowbreak{}]\allowbreak{} & 1 +\allowbreak{} 2\textasciicircum{}\allowbreak{}(\allowbreak{}1/\allowbreak{}5)\allowbreak{} & yes & 1$\to$1 & 0.023\\
L14 & Strad[\allowbreak{}Sqrt[\allowbreak{}5 +\allowbreak{} 2*\allowbreak{}Sqrt[\allowbreak{}6]\allowbreak{}]\allowbreak{}\textasciicircum{}\allowbreak{}(\allowbreak{}1/\allowbreak{}5)\allowbreak{}]\allowbreak{} & -(\allowbreak{}(\allowbreak{}(\allowbreak{}-1)\allowbreak{}\textasciicircum{}\allowbreak{}(\allowbreak{}4/\allowbreak{}5)\allowbreak{}*\allowbreak{}(\allowbreak{}-2 -\allowbreak{} Sqrt[\allowbreak{}6]\allowbreak{})\allowbreak{}\textasciicircum{}\allowbreak{}(\allowbreak{}1/\allowbreak{}5)\allowbreak{})\allowbreak{}/\allowbreak{}2\textasciicircum{}\allowbreak{}(\allowbreak{}1/\allowbreak{}10)\allowbreak{})\allowbreak{}\newline{\scriptsize\itshape Ceiling::\allowbreak{}meprec} & yes & 2$\to$2 & 20.350\\
M01 & PossibleZeroQ[\allowbreak{}Sqrt[\allowbreak{}3/\allowbreak{}2]\allowbreak{} +\allowbreak{} 1/\allowbreak{}Sqrt[\allowbreak{}2]\allowbreak{} -\allowbreak{} Sqrt[\allowbreak{}2 +\allowbreak{} Sqrt[\allowbreak{}3]\allowbreak{}]\allowbreak{}]\allowbreak{} & True & -- & 0$\to$0 & 0.000\\
M02 & PossibleZeroQ[\allowbreak{}Sqrt[\allowbreak{}3/\allowbreak{}2]\allowbreak{} +\allowbreak{} 1/\allowbreak{}Sqrt[\allowbreak{}2]\allowbreak{} -\allowbreak{} Sqrt[\allowbreak{}2 +\allowbreak{} Sqrt[\allowbreak{}3]\allowbreak{}]\allowbreak{},\allowbreak{} Method -> "ExactAlgebraics"]\allowbreak{} & True & -- & 0$\to$0 & 0.000\\
M03 & RootReduce[\allowbreak{}Sqrt[\allowbreak{}3/\allowbreak{}2]\allowbreak{} +\allowbreak{} 1/\allowbreak{}Sqrt[\allowbreak{}2]\allowbreak{} -\allowbreak{} Sqrt[\allowbreak{}2 +\allowbreak{} Sqrt[\allowbreak{}3]\allowbreak{}]\allowbreak{}]\allowbreak{} & 0 & yes & 0$\to$0 & 0.014\\
M04 & PossibleZeroQ[\allowbreak{}-Sqrt[\allowbreak{}3/\allowbreak{}2]\allowbreak{} -\allowbreak{} 1/\allowbreak{}Sqrt[\allowbreak{}2]\allowbreak{} -\allowbreak{} Sqrt[\allowbreak{}2 +\allowbreak{} Sqrt[\allowbreak{}3]\allowbreak{}]\allowbreak{}]\allowbreak{} & False & -- & 0$\to$0 & 0.000\\
M05 & PossibleZeroQ[\allowbreak{}(\allowbreak{}1 +\allowbreak{} I*\allowbreak{}Sqrt[\allowbreak{}2]\allowbreak{})\allowbreak{}*\allowbreak{}(\allowbreak{}1 +\allowbreak{} I*\allowbreak{}Sqrt[\allowbreak{}3]\allowbreak{})\allowbreak{} -\allowbreak{} Sqrt[\allowbreak{}Expand[\allowbreak{}(\allowbreak{}1 +\allowbreak{} I*\allowbreak{}Sqrt[\allowbreak{}2]\allowbreak{})\allowbreak{}\textasciicircum{}\allowbreak{}2*\allowbreak{}(\allowbreak{}1 +\allowbreak{} I*\allowbreak{}Sqrt[\allowbreak{}3]\allowbreak{})\allowbreak{}\textasciicircum{}\allowbreak{}2]\allowbreak{}]\allowbreak{}]\allowbreak{} & False & -- & 0$\to$0 & 0.000\\
M06 & PossibleZeroQ[\allowbreak{}(\allowbreak{}-(\allowbreak{}1 +\allowbreak{} I*\allowbreak{}Sqrt[\allowbreak{}2]\allowbreak{})\allowbreak{})\allowbreak{}*\allowbreak{}(\allowbreak{}1 +\allowbreak{} I*\allowbreak{}Sqrt[\allowbreak{}3]\allowbreak{})\allowbreak{} -\allowbreak{} Sqrt[\allowbreak{}Expand[\allowbreak{}(\allowbreak{}1 +\allowbreak{} I*\allowbreak{}Sqrt[\allowbreak{}2]\allowbreak{})\allowbreak{}\textasciicircum{}\allowbreak{}2*\allowbreak{}(\allowbreak{}1 +\allowbreak{} I*\allowbreak{}Sqrt[\allowbreak{}3]\allowbreak{})\allowbreak{}\textasciicircum{}\allowbreak{}2]\allowbreak{}]\allowbreak{}]\allowbreak{} & True & -- & 0$\to$0 & 0.000\\
N01 & primeProc[\allowbreak{}facs,\allowbreak{} \#1[\allowbreak{}[\allowbreak{}1]\allowbreak{}]\allowbreak{} <= \#1[\allowbreak{}[\allowbreak{}2]\allowbreak{}]\allowbreak{} \& ,\allowbreak{} 2]\allowbreak{} & \{\{2,\allowbreak{} 2003,\allowbreak{} 11,\allowbreak{} 7,\allowbreak{} 5,\allowbreak{} 3\},\allowbreak{} \{2,\allowbreak{} 2003,\allowbreak{} 11,\allowbreak{} 7,\allowbreak{} 5,\allowbreak{} 3\}\} & -- & 0$\to$0 & 0.000\\
N02 & primeProc[\allowbreak{}facs,\allowbreak{} \#1[\allowbreak{}[\allowbreak{}1]\allowbreak{}]\allowbreak{} <= \#2[\allowbreak{}[\allowbreak{}1]\allowbreak{}]\allowbreak{} \& ,\allowbreak{} 2]\allowbreak{} & \{\{2,\allowbreak{} 3,\allowbreak{} 5,\allowbreak{} 7,\allowbreak{} 11,\allowbreak{} 2003\},\allowbreak{} \{2,\allowbreak{} 3,\allowbreak{} 5,\allowbreak{} 7,\allowbreak{} 11\}[\allowbreak{}[\allowbreak{}1 ;; 6]\allowbreak{}]\allowbreak{}\}\newline{\scriptsize\itshape Part::\allowbreak{}take} & -- & 0$\to$0 & 0.000\\
N03 & primeProc[\allowbreak{}\{\{2,\allowbreak{} 14\},\allowbreak{} \{3,\allowbreak{} 2\},\allowbreak{} \{5,\allowbreak{} 1\},\allowbreak{} \{7,\allowbreak{} 1\},\allowbreak{} \{11,\allowbreak{} 1\},\allowbreak{} \{13,\allowbreak{} 1\},\allowbreak{} \{17,\allowbreak{} 1\},\allowbreak{} \{19,\allowbreak{} 1\},\allowbreak{} \{23,\allowbreak{} 1\},\allowbreak{} \{29,\allowbreak{} 1\},\allowbreak{} \{31,\allowbreak{} 1\}\},\allowbreak{} \#1[\allowbreak{}[\allowbreak{}1]\allowbreak{}]\allowbreak{} <= \#1[\allowbreak{}[\allowbreak{}2]\allowbreak{}]\allowbreak{} \& ,\allowbreak{} 2]\allowbreak{} & \{\{2,\allowbreak{} 31,\allowbreak{} 29,\allowbreak{} 23,\allowbreak{} 19,\allowbreak{} 17,\allowbreak{} 13,\allowbreak{} 11,\allowbreak{} 7,\allowbreak{} 5,\allowbreak{} 3\},\allowbreak{} \{2,\allowbreak{} 31,\allowbreak{} 29,\allowbreak{} 23,\allowbreak{} 19,\allowbreak{} 17,\allowbreak{} 13,\allowbreak{} 11,\allowbreak{} 7,\allowbreak{} 5\}\} & -- & 0$\to$0 & 0.000\\
N04 & primeProc[\allowbreak{}\{\{2,\allowbreak{} 14\},\allowbreak{} \{3,\allowbreak{} 2\},\allowbreak{} \{5,\allowbreak{} 1\},\allowbreak{} \{7,\allowbreak{} 1\},\allowbreak{} \{11,\allowbreak{} 1\},\allowbreak{} \{13,\allowbreak{} 1\},\allowbreak{} \{17,\allowbreak{} 1\},\allowbreak{} \{19,\allowbreak{} 1\},\allowbreak{} \{23,\allowbreak{} 1\},\allowbreak{} \{29,\allowbreak{} 1\},\allowbreak{} \{31,\allowbreak{} 1\}\},\allowbreak{} \#1[\allowbreak{}[\allowbreak{}1]\allowbreak{}]\allowbreak{} <= \#2[\allowbreak{}[\allowbreak{}1]\allowbreak{}]\allowbreak{} \& ,\allowbreak{} 2]\allowbreak{} & \{\{2,\allowbreak{} 3,\allowbreak{} 5,\allowbreak{} 7,\allowbreak{} 11,\allowbreak{} 13,\allowbreak{} 17,\allowbreak{} 19,\allowbreak{} 23,\allowbreak{} 29,\allowbreak{} 31\},\allowbreak{} \{2,\allowbreak{} 3,\allowbreak{} 5,\allowbreak{} 7,\allowbreak{} 11,\allowbreak{} 13,\allowbreak{} 17,\allowbreak{} 19,\allowbreak{} 23,\allowbreak{} 29\}\} & -- & 0$\to$0 & 0.000\\
N05 & Length[\allowbreak{}Divisors[\allowbreak{}(\allowbreak{}2*\allowbreak{}3*\allowbreak{}5*\allowbreak{}7*\allowbreak{}11*\allowbreak{}13*\allowbreak{}17*\allowbreak{}19*\allowbreak{}23*\allowbreak{}29)\allowbreak{}\textasciicircum{}\allowbreak{}(\allowbreak{}2 -\allowbreak{} 1)\allowbreak{}]\allowbreak{}]\allowbreak{} & 1024 & yes & 0$\to$0 & 0.000\\
N06 & Length[\allowbreak{}Divisors[\allowbreak{}(\allowbreak{}2*\allowbreak{}3*\allowbreak{}5*\allowbreak{}7*\allowbreak{}11)\allowbreak{}\textasciicircum{}\allowbreak{}(\allowbreak{}10 -\allowbreak{} 1)\allowbreak{}]\allowbreak{}]\allowbreak{} & 100000 & yes & 0$\to$0 & 0.039\\
\end{longtable}

\paragraph{Negative bases and output quality (Table~\ref{tab:r3QRS}).}
Q01, Q02, Q07, Q08: for $(p+q\sqrt c)^3<0$ the principal cube root is
$e^{i\pi/3}|p+q\sqrt c|$; the program's ``solutions'' are equal to the input but
contain $\sqrt{9\pm6\sqrt2}$ and similar denestable surds, so the depth does not
decrease (\U{23}). Q03--Q06 (fourth roots and positive bases) are fine.
\begin{longtable}{@{}p{0.055\textwidth}>{\ttfamily\footnotesize\raggedright\arraybackslash}p{0.33\textwidth}>{\ttfamily\footnotesize\raggedright\arraybackslash}p{0.33\textwidth}p{0.05\textwidth}p{0.05\textwidth}p{0.05\textwidth}@{}}
\caption{Original program: roots of negative bases (Q), the \texttt{Ceiling[Log]} issue (R), output quality for a four-term square root (S).}\label{tab:r3QRS}\\
\toprule ID & Input & Output & Equal & Depth & Time (s)\\\midrule\endfirsthead
\toprule ID & Input & Output & Equal & Depth & Time (s)\\\midrule\endhead
\bottomrule\endfoot
Q01 & Strad[\allowbreak{}Expand[\allowbreak{}(\allowbreak{}-1 -\allowbreak{} Sqrt[\allowbreak{}2]\allowbreak{})\allowbreak{}\textasciicircum{}\allowbreak{}3]\allowbreak{}\textasciicircum{}\allowbreak{}(\allowbreak{}1/\allowbreak{}3)\allowbreak{}]\allowbreak{} & (\allowbreak{}(\allowbreak{}-1)\allowbreak{}\textasciicircum{}\allowbreak{}(\allowbreak{}2/\allowbreak{}3)\allowbreak{}*\allowbreak{}(\allowbreak{}1 +\allowbreak{} Sqrt[\allowbreak{}2]\allowbreak{} -\allowbreak{} I*\allowbreak{}Sqrt[\allowbreak{}9 +\allowbreak{} 6*\allowbreak{}Sqrt[\allowbreak{}2]\allowbreak{}]\allowbreak{})\allowbreak{})\allowbreak{}/\allowbreak{}2 & yes & 2$\to$2 & 0.021\\
Q02 & Strad[\allowbreak{}Expand[\allowbreak{}(\allowbreak{}1 -\allowbreak{} Sqrt[\allowbreak{}2]\allowbreak{})\allowbreak{}\textasciicircum{}\allowbreak{}3]\allowbreak{}\textasciicircum{}\allowbreak{}(\allowbreak{}1/\allowbreak{}3)\allowbreak{}]\allowbreak{} & (\allowbreak{}(\allowbreak{}-1)\allowbreak{}\textasciicircum{}\allowbreak{}(\allowbreak{}2/\allowbreak{}3)\allowbreak{}*\allowbreak{}(\allowbreak{}-1 +\allowbreak{} Sqrt[\allowbreak{}2]\allowbreak{} -\allowbreak{} I*\allowbreak{}Sqrt[\allowbreak{}9 -\allowbreak{} 6*\allowbreak{}Sqrt[\allowbreak{}2]\allowbreak{}]\allowbreak{})\allowbreak{})\allowbreak{}/\allowbreak{}2 & yes & 2$\to$2 & 0.019\\
Q03 & Strad[\allowbreak{}Expand[\allowbreak{}(\allowbreak{}-1 -\allowbreak{} Sqrt[\allowbreak{}2]\allowbreak{})\allowbreak{}\textasciicircum{}\allowbreak{}4]\allowbreak{}\textasciicircum{}\allowbreak{}(\allowbreak{}1/\allowbreak{}4)\allowbreak{}]\allowbreak{} & 1 +\allowbreak{} Sqrt[\allowbreak{}2]\allowbreak{} & yes & 1$\to$1 & 0.008\\
Q04 & Strad[\allowbreak{}Expand[\allowbreak{}(\allowbreak{}1 -\allowbreak{} Sqrt[\allowbreak{}2]\allowbreak{})\allowbreak{}\textasciicircum{}\allowbreak{}4]\allowbreak{}\textasciicircum{}\allowbreak{}(\allowbreak{}1/\allowbreak{}4)\allowbreak{}]\allowbreak{} & -1 +\allowbreak{} Sqrt[\allowbreak{}2]\allowbreak{} & yes & 1$\to$1 & 0.013\\
Q05 & Strad[\allowbreak{}Expand[\allowbreak{}(\allowbreak{}2 -\allowbreak{} 3*\allowbreak{}Sqrt[\allowbreak{}3]\allowbreak{})\allowbreak{}\textasciicircum{}\allowbreak{}4]\allowbreak{}\textasciicircum{}\allowbreak{}(\allowbreak{}1/\allowbreak{}4)\allowbreak{}]\allowbreak{} & -2 +\allowbreak{} 3*\allowbreak{}Sqrt[\allowbreak{}3]\allowbreak{} & yes & 1$\to$1 & 0.015\\
Q06 & Strad[\allowbreak{}Expand[\allowbreak{}(\allowbreak{}3 -\allowbreak{} Sqrt[\allowbreak{}2]\allowbreak{})\allowbreak{}\textasciicircum{}\allowbreak{}3]\allowbreak{}\textasciicircum{}\allowbreak{}(\allowbreak{}1/\allowbreak{}3)\allowbreak{}]\allowbreak{} & 3 -\allowbreak{} Sqrt[\allowbreak{}2]\allowbreak{} & yes & 1$\to$1 & 0.014\\
Q07 & Strad[\allowbreak{}Expand[\allowbreak{}(\allowbreak{}2 -\allowbreak{} 3*\allowbreak{}Sqrt[\allowbreak{}2]\allowbreak{})\allowbreak{}\textasciicircum{}\allowbreak{}3]\allowbreak{}\textasciicircum{}\allowbreak{}(\allowbreak{}1/\allowbreak{}3)\allowbreak{}]\allowbreak{} & -(\allowbreak{}(\allowbreak{}(\allowbreak{}-1)\allowbreak{}\textasciicircum{}\allowbreak{}(\allowbreak{}2/\allowbreak{}3)\allowbreak{}*\allowbreak{}(\allowbreak{}-3 +\allowbreak{} Sqrt[\allowbreak{}2]\allowbreak{} +\allowbreak{} I*\allowbreak{}Sqrt[\allowbreak{}33 -\allowbreak{} 18*\allowbreak{}Sqrt[\allowbreak{}2]\allowbreak{}]\allowbreak{})\allowbreak{})\allowbreak{}/\allowbreak{}Sqrt[\allowbreak{}2]\allowbreak{})\allowbreak{} & yes & 2$\to$2 & 0.023\\
Q08 & Strad[\allowbreak{}(\allowbreak{}-7 -\allowbreak{} 5*\allowbreak{}Sqrt[\allowbreak{}2]\allowbreak{})\allowbreak{}\textasciicircum{}\allowbreak{}(\allowbreak{}1/\allowbreak{}3)\allowbreak{}]\allowbreak{} & (\allowbreak{}(\allowbreak{}-1)\allowbreak{}\textasciicircum{}\allowbreak{}(\allowbreak{}2/\allowbreak{}3)\allowbreak{}*\allowbreak{}(\allowbreak{}1 +\allowbreak{} Sqrt[\allowbreak{}2]\allowbreak{} -\allowbreak{} I*\allowbreak{}Sqrt[\allowbreak{}9 +\allowbreak{} 6*\allowbreak{}Sqrt[\allowbreak{}2]\allowbreak{}]\allowbreak{})\allowbreak{})\allowbreak{}/\allowbreak{}2 & yes & 2$\to$2 & 0.017\\
Q09 & denest11[\allowbreak{}(\allowbreak{}7 -\allowbreak{} 5*\allowbreak{}Sqrt[\allowbreak{}2]\allowbreak{})\allowbreak{}\textasciicircum{}\allowbreak{}(\allowbreak{}1/\allowbreak{}3)\allowbreak{}]\allowbreak{} & (\allowbreak{}(\allowbreak{}-1)\allowbreak{}\textasciicircum{}\allowbreak{}(\allowbreak{}2/\allowbreak{}3)\allowbreak{}*\allowbreak{}(\allowbreak{}-1 +\allowbreak{} Sqrt[\allowbreak{}2]\allowbreak{} -\allowbreak{} I*\allowbreak{}Sqrt[\allowbreak{}9 -\allowbreak{} 6*\allowbreak{}Sqrt[\allowbreak{}2]\allowbreak{}]\allowbreak{})\allowbreak{})\allowbreak{}/\allowbreak{}2 & yes & 2$\to$2 & 0.013\\
R01 & Ceiling[\allowbreak{}Log[\allowbreak{}1000]\allowbreak{}/\allowbreak{}Log[\allowbreak{}10]\allowbreak{}]\allowbreak{} & Ceiling[\allowbreak{}Log[\allowbreak{}1000]\allowbreak{}/\allowbreak{}Log[\allowbreak{}10]\allowbreak{}]\allowbreak{}\newline{\scriptsize\itshape Ceiling::\allowbreak{}meprec} & yes & 0$\to$0 & 0.001\\
R02 & Max[\allowbreak{}5,\allowbreak{} Ceiling[\allowbreak{}Log[\allowbreak{}1000]\allowbreak{}/\allowbreak{}Log[\allowbreak{}10]\allowbreak{}]\allowbreak{}]\allowbreak{} & 5\newline{\scriptsize\itshape Ceiling::\allowbreak{}meprec} & yes & 0$\to$0 & 0.001\\
R03 & Min[\allowbreak{}Max[\allowbreak{}5,\allowbreak{} Ceiling[\allowbreak{}Log[\allowbreak{}1000]\allowbreak{}/\allowbreak{}Log[\allowbreak{}10]\allowbreak{}]\allowbreak{}]\allowbreak{},\allowbreak{} 3]\allowbreak{} & 3\newline{\scriptsize\itshape Ceiling::\allowbreak{}meprec} & yes & 0$\to$0 & 0.000\\
R04 & \{2,\allowbreak{} 3,\allowbreak{} 5\}[\allowbreak{}[\allowbreak{}1 ;; Min[\allowbreak{}Max[\allowbreak{}5,\allowbreak{} Ceiling[\allowbreak{}Log[\allowbreak{}1000]\allowbreak{}/\allowbreak{}Log[\allowbreak{}10]\allowbreak{}]\allowbreak{}]\allowbreak{},\allowbreak{} 3]\allowbreak{}]\allowbreak{}]\allowbreak{} & \{2,\allowbreak{} 3,\allowbreak{} 5\}\newline{\scriptsize\itshape Ceiling::\allowbreak{}meprec} & -- & 0$\to$0 & 0.000\\
R05 & Ceiling[\allowbreak{}Log[\allowbreak{}1000]\allowbreak{}/\allowbreak{}Log[\allowbreak{}100]\allowbreak{}]\allowbreak{} & 2 & yes & 0$\to$0 & 0.000\\
R06 & Ceiling[\allowbreak{}Log[\allowbreak{}1000]\allowbreak{}/\allowbreak{}Log[\allowbreak{}1000]\allowbreak{}]\allowbreak{} & 1 & yes & 0$\to$0 & 0.000\\
R07 & Ceiling[\allowbreak{}Log[\allowbreak{}1000]\allowbreak{}/\allowbreak{}Log[\allowbreak{}6]\allowbreak{}]\allowbreak{} & 4 & yes & 0$\to$0 & 0.000\\
S01 & Strad[\allowbreak{}Sqrt[\allowbreak{}17 +\allowbreak{} 2*\allowbreak{}Sqrt[\allowbreak{}6]\allowbreak{} +\allowbreak{} 2*\allowbreak{}Sqrt[\allowbreak{}10]\allowbreak{} +\allowbreak{} 2*\allowbreak{}Sqrt[\allowbreak{}14]\allowbreak{} +\allowbreak{} 2*\allowbreak{}Sqrt[\allowbreak{}15]\allowbreak{} +\allowbreak{} 2*\allowbreak{}Sqrt[\allowbreak{}21]\allowbreak{} +\allowbreak{} 2*\allowbreak{}Sqrt[\allowbreak{}35]\allowbreak{}]\allowbreak{}]\allowbreak{} & (\allowbreak{}6397*\allowbreak{}Sqrt[\allowbreak{}2]\allowbreak{} +\allowbreak{} 5316*\allowbreak{}Sqrt[\allowbreak{}3]\allowbreak{} +\allowbreak{} 4090*\allowbreak{}Sqrt[\allowbreak{}5]\allowbreak{} +\allowbreak{} 3432*\allowbreak{}Sqrt[\allowbreak{}7]\allowbreak{} +\allowbreak{} 1695*\allowbreak{}Sqrt[\allowbreak{}30]\allowbreak{} +\allowbreak{} 1423*\allowbreak{}Sqrt[\allowbreak{}42]\allowbreak{} +\allowbreak{} 1095*\allowbreak{}Sqrt[\allowbreak{}70]\allowbreak{} +\allowbreak{} 910*\allowbreak{}Sqrt[\allowbreak{}105]\allowbreak{})\allowbreak{}/\allowbreak{}(\allowbreak{}1108 +\allowbreak{} 467*\allowbreak{}Sqrt[\allowbreak{}6]\allowbreak{} +\allowbreak{} 359*\allowbreak{}Sqrt[\allowbreak{}10]\allowbreak{} +\allowbreak{} 299*\allowbreak{}Sqrt[\allowbreak{}14]\allowbreak{} +\allowbreak{} 302*\allowbreak{}Sqrt[\allowbreak{}15]\allowbreak{} +\allowbreak{} 252*\allowbreak{}Sqrt[\allowbreak{}21]\allowbreak{} +\allowbreak{} 194*\allowbreak{}Sqrt[\allowbreak{}35]\allowbreak{} +\allowbreak{} 81*\allowbreak{}Sqrt[\allowbreak{}210]\allowbreak{})\allowbreak{} & yes & 1$\to$1 & 2.355\\
\end{longtable}

\subsection{Randomized families}
\label{sec:fuzz}
Ten structured families with a fixed seed (\code{SeedRandom[20260905]}),
$p,q,r\in\{\pm1,\ldots,\pm7\}$ and $c,d$ square-free in $\{2,\ldots,17\}$, were
run through \code{Strad} of both versions with a 60\,s limit per input. Each
result was classified by exact comparison with the input. Table~\ref{tab:fuzz}
summarizes; ``wrong'' means the returned value differs from the input.

\begin{table}[h]\centering\scriptsize\setlength{\tabcolsep}{3.5pt}
\begin{tabular}{@{}l r rrrr rrr@{}}
\toprule
& & \multicolumn{4}{c}{Original} & \multicolumn{3}{c}{Corrected}\\
\cmidrule(lr){3-6}\cmidrule(lr){7-9}
Family & $n$ & denested & equal, rewritten & unchanged & \textbf{wrong} & denested & unchanged & \textbf{wrong}\\\midrule
F1 $\sqrt{(p+q\sqrt c)^2}$ & 40 & 40 & 0 & 0 & 0 & 40 & 0 & 0\\
F2 $\big((p+q\sqrt c)^3\big)^{1/3}$ & 30 & 13 & 17 & 0 & 0 & 30 & 0 & 0\\
F3 $\sqrt{a+b\sqrt c}$, random & 40 & 6 & 8 & 26 & 0 & 6 & 34 & 0\\
F4 $\sqrt{(p+qi\sqrt c)^2}$ & 40 & 40 & 0 & 0 & 0 & 40 & 0 & 0\\
F5 $\big((p+qi\sqrt c)^2\big)^{3/2}$ & 40 & 1 & 0 & 0 & \textbf{39} & 40 & 0 & 0\\
F6 $\sqrt{u^2}\sqrt{v^2}$, complex & 40 & 23 & 0 & 0 & \textbf{17} & 40 & 0 & 0\\
F7 $\big((p+q\sqrt c)^4\big)^{1/4}$ & 25 & 17 & 8 & 0 & 0 & 25 & 0 & 0\\
F8 $\sqrt{(p+q\sqrt c+r\sqrt d)^2}$ & 25 & 25 & 0 & 0 & 0 & 25 & 0 & 0\\
F9 $\big((p+q\sqrt c)^2\big)^{3/2}$, real & 30 & 30 & 0 & 0 & 0 & 30 & 0 & 0\\
F10 $\sqrt{-u^2}\sqrt{-v^2}$, real & 30 & 30 & 0 & 0 & 0 & 30 & 0 & 0\\
\midrule
Total & 340 & 225 & 33 & 26 & \textbf{56} & 306 & 34 & \textbf{0}\\
\bottomrule
\end{tabular}
\caption{Randomized families. The original program emitted
\code{PossibleZeroQ::ztest1} on 272 of the 340 inputs; the corrected version
emitted no message. Wall-clock totals: original 27\,s, corrected 39\,s;
the slowest single input took 0.4\,s (original) and 1.7\,s (corrected).}
\label{tab:fuzz}
\end{table}

The F3 family is a control: with $a\in[-30,30]$ random, only 6 of 40
radicands satisfy the square criterion of Proposition~\ref{prop:surd}, and
both versions denest exactly those 6. The 8 ``equal, rewritten'' outputs of
the original in F3 are the final \code{Simplify} at work (\U{20}); the
corrected version returns such inputs unchanged. The 17 rewritten outputs in
F2 and 8 in F7 are the negative-base cases of Table~\ref{tab:r3QRS}.

\subsection{The reviews' regression suites}
\label{sec:expsuites}
Table~\ref{tab:suites} lists the outcome of each review's own Wolfram tests on
the original program. Every ``expected to fail'' contract fails, every control
passes, and no test aborted.

\begin{table}[h]\centering\small
\begin{tabularx}{\textwidth}{@{}p{0.1\textwidth} X X@{}}
\toprule Suite & Failing tests (all predicted by the review) & Passing\\\midrule
\RI\ (14) & Print protection; comparator; negative branch; grouped complex branch; wrapper complex product; \code{Factorc[Sqrt[t^2]]}; zero multiplier & 7: rationalizer (2), direct and multiplier success, \code{Sqrt[t + t^2]} (\emph{not} predicted to pass), empty seed list, all-levels\\
\RII\ (15) & comparator; negative branch; complex $3/2$ power; merged product; \code{Factorc[Sqrt[z^2]]}; caller's \code{f}; numerator \code{x} & 4 on the original (Print protection\footnotemark, worked example, rationalizer, all-level control); the 4 tests of its guarded component pass when the component is loaded (Section~\ref{sec:expkernels})\\
\RIII\ (20) & negative core target; complex $3/2$; merged product (core and wrapper); \code{Factorc[Sqrt[-x-y]]}; zero multiplier; comparator; stale slice; cap; Print protection & 10: classical success, rationalizer (2), non-nested input, wrapper modes with \code{Identity}, and all 5 tests of its proposed kernel\\
\bottomrule
\end{tabularx}
\caption{The reviews' regression suites executed against the original program.}
\label{tab:suites}
\end{table}
\footnotetext{\RII's suite restores the attributes of \code{Print} itself after
loading and then tests the recorded values, so this test measures the suite's
own bookkeeping rather than the defect; \RI's and \RIII's versions of the same
test fail as they should.}

\subsection{The reviews' proposed kernels}
\label{sec:expkernels}
The three one-trial components (\RI\ \code{VerifiedMultiplierTrial},
\RII\ \code{CertifiedMultiplierStep}, \RIII\ \code{TryDenestCandidate}) were
loaded and called on eleven common cases. All three certify the correct
branch for the negative target, the complex $3/2$ power, the complex product
and the cube-root case with $m=9$; all three reject a zero multiplier, an
inexact multiplier, root order 1 and a symbolic target with a structured
\code{Failure}; all three report ``no lower-degree GCD'' for
$\sqrt{5+2\sqrt6}$ with $m=1$ and find $\sqrt2+\sqrt3$ with $m=2$, and, because
none of them has a degree gate, all three also find it with \RII's
equal-degree multiplier (returning the same ugly quotient as the original,
D07). \RIII's \code{BestImprovement} policy correctly rejects that quotient
as ``not an improvement'', but would accept the original's BFHT output
(Table~
ef{tab:origA}, A31), because its cost tuple puts radical depth before
size: $(0,2,9,69)$ beats $(0,3,4,28)$. In other words the components work as
their authors said, and they are single steps, not denesters: none contains a
multiplier generator, a wrapper, or a search.

\subsection{The corrected version}
\label{sec:expfixed}
Tables~\ref{tab:fixedA}--\ref{tab:fixedLQ} repeat the battery with
\code{StradFixed.wl}. Every row is exactly equal to the true value; no row
emits a message; \code{Print} keeps its \code{Protected} attribute. Beyond
fixing every wrong value, the corrected version finds denestings the original
does not: the BFHT example (A30/A31), the four-term square root as
$\sqrt2+\sqrt3+\sqrt5+\sqrt7$ (L04), $(49+20\sqrt6)^{1/4}=\sqrt2+\sqrt3$ (L06),
the cube roots of negative bases (Q01--Q08) and the equal-degree multiplier
(D06). It also returns inputs unchanged when it cannot improve them (B01, B08,
B09, L02, L03, L09) instead of rewriting them, converts a \code{Root} input to
$\sqrt2+\sqrt3$ (B07), maps over lists (K14), leaves symbolic hosts alone
(E01--E05, E07) while still denesting numeric parts inside them (E06, E08), and
rejects invalid multipliers without garbage (D02, D10, D11). The price is a
factor of roughly two to three in time on the small inputs, almost entirely
spent in \code{RootReduce} certification.
\begin{longtable}{@{}p{0.055\textwidth}>{\ttfamily\footnotesize\raggedright\arraybackslash}p{0.33\textwidth}>{\ttfamily\footnotesize\raggedright\arraybackslash}p{0.33\textwidth}p{0.05\textwidth}p{0.05\textwidth}p{0.05\textwidth}@{}}
\caption{Corrected version on the classical identities (compare Table~\ref{tab:origA}).}\label{tab:fixedA}\\
\toprule ID & Input & Output & Equal & Depth & Time (s)\\\midrule\endfirsthead
\toprule ID & Input & Output & Equal & Depth & Time (s)\\\midrule\endhead
\bottomrule\endfoot
A01 & Strad[\allowbreak{}Sqrt[\allowbreak{}3 +\allowbreak{} 2*\allowbreak{}Sqrt[\allowbreak{}2]\allowbreak{}]\allowbreak{}]\allowbreak{} & 1 +\allowbreak{} Sqrt[\allowbreak{}2]\allowbreak{} & yes & 2$\to$1 & 0.103\\
A02 & Strad[\allowbreak{}Sqrt[\allowbreak{}5 +\allowbreak{} 2*\allowbreak{}Sqrt[\allowbreak{}6]\allowbreak{}]\allowbreak{}]\allowbreak{} & Sqrt[\allowbreak{}2]\allowbreak{} +\allowbreak{} Sqrt[\allowbreak{}3]\allowbreak{} & yes & 2$\to$1 & 0.055\\
A03 & Strad[\allowbreak{}Sqrt[\allowbreak{}2 +\allowbreak{} Sqrt[\allowbreak{}3]\allowbreak{}]\allowbreak{}]\allowbreak{} & Sqrt[\allowbreak{}3/\allowbreak{}2]\allowbreak{} +\allowbreak{} 1/\allowbreak{}Sqrt[\allowbreak{}2]\allowbreak{} & yes & 2$\to$1 & 0.057\\
A04 & Strad[\allowbreak{}Sqrt[\allowbreak{}11 +\allowbreak{} 6*\allowbreak{}Sqrt[\allowbreak{}2]\allowbreak{}]\allowbreak{}]\allowbreak{} & 3 +\allowbreak{} Sqrt[\allowbreak{}2]\allowbreak{} & yes & 2$\to$1 & 0.030\\
A05 & Strad[\allowbreak{}Sqrt[\allowbreak{}3 +\allowbreak{} Sqrt[\allowbreak{}5]\allowbreak{}]\allowbreak{}]\allowbreak{} & 1/\allowbreak{}Sqrt[\allowbreak{}2]\allowbreak{} +\allowbreak{} Sqrt[\allowbreak{}5/\allowbreak{}2]\allowbreak{} & yes & 2$\to$1 & 0.060\\
A06 & Strad[\allowbreak{}Sqrt[\allowbreak{}7 +\allowbreak{} 2*\allowbreak{}Sqrt[\allowbreak{}10]\allowbreak{}]\allowbreak{}]\allowbreak{} & Sqrt[\allowbreak{}2]\allowbreak{} +\allowbreak{} Sqrt[\allowbreak{}5]\allowbreak{} & yes & 2$\to$1 & 0.040\\
A07 & Strad[\allowbreak{}Sqrt[\allowbreak{}6 +\allowbreak{} 2*\allowbreak{}Sqrt[\allowbreak{}2]\allowbreak{} +\allowbreak{} 2*\allowbreak{}Sqrt[\allowbreak{}3]\allowbreak{} +\allowbreak{} 2*\allowbreak{}Sqrt[\allowbreak{}6]\allowbreak{}]\allowbreak{}]\allowbreak{} & 1 +\allowbreak{} Sqrt[\allowbreak{}2]\allowbreak{} +\allowbreak{} Sqrt[\allowbreak{}3]\allowbreak{} & yes & 2$\to$1 & 0.093\\
A08 & Strad[\allowbreak{}Sqrt[\allowbreak{}12 +\allowbreak{} 2*\allowbreak{}Sqrt[\allowbreak{}6]\allowbreak{} +\allowbreak{} 2*\allowbreak{}Sqrt[\allowbreak{}14]\allowbreak{} +\allowbreak{} 2*\allowbreak{}Sqrt[\allowbreak{}21]\allowbreak{}]\allowbreak{}]\allowbreak{} & Sqrt[\allowbreak{}2]\allowbreak{} +\allowbreak{} Sqrt[\allowbreak{}3]\allowbreak{} +\allowbreak{} Sqrt[\allowbreak{}7]\allowbreak{} & yes & 2$\to$1 & 0.231\\
A09 & Strad[\allowbreak{}(\allowbreak{}Sqrt[\allowbreak{}5]\allowbreak{} -\allowbreak{} 2)\allowbreak{}\textasciicircum{}\allowbreak{}(\allowbreak{}1/\allowbreak{}3)\allowbreak{}]\allowbreak{} & (\allowbreak{}-1 +\allowbreak{} Sqrt[\allowbreak{}5]\allowbreak{})\allowbreak{}/\allowbreak{}2 & yes & 2$\to$1 & 0.045\\
A10 & Strad[\allowbreak{}(\allowbreak{}2 +\allowbreak{} Sqrt[\allowbreak{}5]\allowbreak{})\allowbreak{}\textasciicircum{}\allowbreak{}(\allowbreak{}1/\allowbreak{}3)\allowbreak{}]\allowbreak{} & (\allowbreak{}1 +\allowbreak{} Sqrt[\allowbreak{}5]\allowbreak{})\allowbreak{}/\allowbreak{}2 & yes & 2$\to$1 & 0.046\\
A11 & Strad[\allowbreak{}(\allowbreak{}10 +\allowbreak{} 6*\allowbreak{}Sqrt[\allowbreak{}3]\allowbreak{})\allowbreak{}\textasciicircum{}\allowbreak{}(\allowbreak{}1/\allowbreak{}3)\allowbreak{}]\allowbreak{} & 1 +\allowbreak{} Sqrt[\allowbreak{}3]\allowbreak{} & yes & 2$\to$1 & 0.066\\
A12 & Strad[\allowbreak{}(\allowbreak{}2\textasciicircum{}\allowbreak{}(\allowbreak{}1/\allowbreak{}3)\allowbreak{} -\allowbreak{} 1)\allowbreak{}\textasciicircum{}\allowbreak{}(\allowbreak{}1/\allowbreak{}3)\allowbreak{}]\allowbreak{} & (\allowbreak{}1 -\allowbreak{} 2\textasciicircum{}\allowbreak{}(\allowbreak{}1/\allowbreak{}3)\allowbreak{} +\allowbreak{} 2\textasciicircum{}\allowbreak{}(\allowbreak{}2/\allowbreak{}3)\allowbreak{})\allowbreak{}/\allowbreak{}3\textasciicircum{}\allowbreak{}(\allowbreak{}2/\allowbreak{}3)\allowbreak{} & yes & 2$\to$1 & 0.337\\
A13 & Strad[\allowbreak{}(\allowbreak{}5\textasciicircum{}\allowbreak{}(\allowbreak{}1/\allowbreak{}3)\allowbreak{} -\allowbreak{} 4\textasciicircum{}\allowbreak{}(\allowbreak{}1/\allowbreak{}3)\allowbreak{})\allowbreak{}\textasciicircum{}\allowbreak{}(\allowbreak{}1/\allowbreak{}3)\allowbreak{}]\allowbreak{} & (\allowbreak{}-2\textasciicircum{}\allowbreak{}(\allowbreak{}2/\allowbreak{}3)\allowbreak{} +\allowbreak{} 5\textasciicircum{}\allowbreak{}(\allowbreak{}1/\allowbreak{}3)\allowbreak{})\allowbreak{}\textasciicircum{}\allowbreak{}(\allowbreak{}1/\allowbreak{}3)\allowbreak{} & yes & 2$\to$2 & 1.503\\
A14 & Strad[\allowbreak{}(\allowbreak{}(\allowbreak{}11 +\allowbreak{} 5*\allowbreak{}Sqrt[\allowbreak{}5]\allowbreak{})\allowbreak{}/\allowbreak{}2)\allowbreak{}\textasciicircum{}\allowbreak{}(\allowbreak{}1/\allowbreak{}5)\allowbreak{}]\allowbreak{} & (\allowbreak{}1 +\allowbreak{} Sqrt[\allowbreak{}5]\allowbreak{})\allowbreak{}/\allowbreak{}2 & yes & 2$\to$1 & 0.045\\
A15 & Strad[\allowbreak{}Sqrt[\allowbreak{}1 +\allowbreak{} 2*\allowbreak{}2\textasciicircum{}\allowbreak{}(\allowbreak{}1/\allowbreak{}3)\allowbreak{} +\allowbreak{} 2\textasciicircum{}\allowbreak{}(\allowbreak{}2/\allowbreak{}3)\allowbreak{}]\allowbreak{}]\allowbreak{} & 1 +\allowbreak{} 2\textasciicircum{}\allowbreak{}(\allowbreak{}1/\allowbreak{}3)\allowbreak{} & yes & 2$\to$1 & 0.019\\
A16 & Strad[\allowbreak{}(\allowbreak{}17 +\allowbreak{} 12*\allowbreak{}Sqrt[\allowbreak{}2]\allowbreak{})\allowbreak{}\textasciicircum{}\allowbreak{}(\allowbreak{}1/\allowbreak{}4)\allowbreak{}]\allowbreak{} & 1 +\allowbreak{} Sqrt[\allowbreak{}2]\allowbreak{} & yes & 2$\to$1 & 0.042\\
A17 & Strad[\allowbreak{}(\allowbreak{}99 +\allowbreak{} 70*\allowbreak{}Sqrt[\allowbreak{}2]\allowbreak{})\allowbreak{}\textasciicircum{}\allowbreak{}(\allowbreak{}1/\allowbreak{}6)\allowbreak{}]\allowbreak{} & 1 +\allowbreak{} Sqrt[\allowbreak{}2]\allowbreak{} & yes & 2$\to$1 & 0.044\\
A18 & Strad[\allowbreak{}(\allowbreak{}3 +\allowbreak{} 2*\allowbreak{}Sqrt[\allowbreak{}2]\allowbreak{})\allowbreak{}\textasciicircum{}\allowbreak{}(\allowbreak{}3/\allowbreak{}2)\allowbreak{}]\allowbreak{} & 7 +\allowbreak{} 5*\allowbreak{}Sqrt[\allowbreak{}2]\allowbreak{} & yes & 2$\to$1 & 0.027\\
A19 & Strad[\allowbreak{}(\allowbreak{}3 +\allowbreak{} 2*\allowbreak{}Sqrt[\allowbreak{}2]\allowbreak{})\allowbreak{}\textasciicircum{}\allowbreak{}(\allowbreak{}-2\textasciicircum{}\allowbreak{}(\allowbreak{}-1)\allowbreak{})\allowbreak{}]\allowbreak{} & -1 +\allowbreak{} Sqrt[\allowbreak{}2]\allowbreak{} & yes & 2$\to$1 & 0.049\\
A20 & Strad[\allowbreak{}1/\allowbreak{}Sqrt[\allowbreak{}3 +\allowbreak{} 2*\allowbreak{}Sqrt[\allowbreak{}2]\allowbreak{}]\allowbreak{}]\allowbreak{} & -1 +\allowbreak{} Sqrt[\allowbreak{}2]\allowbreak{} & yes & 2$\to$1 & 0.042\\
A21 & Strad[\allowbreak{}Sqrt[\allowbreak{}3 +\allowbreak{} 2*\allowbreak{}Sqrt[\allowbreak{}2]\allowbreak{}]\allowbreak{} +\allowbreak{} Sqrt[\allowbreak{}5 +\allowbreak{} 2*\allowbreak{}Sqrt[\allowbreak{}6]\allowbreak{}]\allowbreak{}]\allowbreak{} & 1 +\allowbreak{} 2*\allowbreak{}Sqrt[\allowbreak{}2]\allowbreak{} +\allowbreak{} Sqrt[\allowbreak{}3]\allowbreak{} & yes & 2$\to$1 & 0.043\\
A22 & Strad[\allowbreak{}Sqrt[\allowbreak{}3 +\allowbreak{} 2*\allowbreak{}Sqrt[\allowbreak{}2]\allowbreak{}]\allowbreak{}*\allowbreak{}Sqrt[\allowbreak{}5 +\allowbreak{} 2*\allowbreak{}Sqrt[\allowbreak{}6]\allowbreak{}]\allowbreak{}]\allowbreak{} & 2 +\allowbreak{} Sqrt[\allowbreak{}2]\allowbreak{} +\allowbreak{} Sqrt[\allowbreak{}3]\allowbreak{} +\allowbreak{} Sqrt[\allowbreak{}6]\allowbreak{} & yes & 2$\to$1 & 0.094\\
A23 & Strad[\allowbreak{}Sqrt[\allowbreak{}2 +\allowbreak{} Sqrt[\allowbreak{}3]\allowbreak{}]\allowbreak{} +\allowbreak{} Sqrt[\allowbreak{}2 -\allowbreak{} Sqrt[\allowbreak{}3]\allowbreak{}]\allowbreak{}]\allowbreak{} & Sqrt[\allowbreak{}6]\allowbreak{} & yes & 2$\to$1 & 0.058\\
A24 & Strad[\allowbreak{}Sqrt[\allowbreak{}-1 +\allowbreak{} 2*\allowbreak{}I*\allowbreak{}Sqrt[\allowbreak{}2]\allowbreak{}]\allowbreak{}]\allowbreak{} & 1 +\allowbreak{} I*\allowbreak{}Sqrt[\allowbreak{}2]\allowbreak{} & yes & 2$\to$1 & 0.049\\
A25 & Strad[\allowbreak{}Sqrt[\allowbreak{}2 +\allowbreak{} 2*\allowbreak{}I*\allowbreak{}Sqrt[\allowbreak{}3]\allowbreak{}]\allowbreak{}]\allowbreak{} & I +\allowbreak{} Sqrt[\allowbreak{}3]\allowbreak{} & yes & 2$\to$1 & 0.071\\
A26 & Strad[\allowbreak{}Sqrt[\allowbreak{}-3 -\allowbreak{} 2*\allowbreak{}Sqrt[\allowbreak{}2]\allowbreak{}]\allowbreak{}]\allowbreak{} & I*\allowbreak{}(\allowbreak{}1 +\allowbreak{} Sqrt[\allowbreak{}2]\allowbreak{})\allowbreak{} & yes & 2$\to$1 & 0.028\\
A27 & Strad[\allowbreak{}(\allowbreak{}3 +\allowbreak{} 2*\allowbreak{}Sqrt[\allowbreak{}2]\allowbreak{})\allowbreak{}\textasciicircum{}\allowbreak{}(\allowbreak{}1/\allowbreak{}4)\allowbreak{}]\allowbreak{} & Sqrt[\allowbreak{}1 +\allowbreak{} Sqrt[\allowbreak{}2]\allowbreak{}]\allowbreak{} & yes & 2$\to$2 & 0.070\\
A28 & Strad[\allowbreak{}Sqrt[\allowbreak{}3 +\allowbreak{} 2*\allowbreak{}Sqrt[\allowbreak{}3 +\allowbreak{} 2*\allowbreak{}Sqrt[\allowbreak{}2]\allowbreak{}]\allowbreak{}]\allowbreak{},\allowbreak{} True]\allowbreak{} & Sqrt[\allowbreak{}5 +\allowbreak{} 2*\allowbreak{}Sqrt[\allowbreak{}2]\allowbreak{}]\allowbreak{} & yes & 3$\to$2 & 0.090\\
A29 & Strad[\allowbreak{}Sqrt[\allowbreak{}3 +\allowbreak{} 2*\allowbreak{}Sqrt[\allowbreak{}3 +\allowbreak{} 2*\allowbreak{}Sqrt[\allowbreak{}2]\allowbreak{}]\allowbreak{}]\allowbreak{}]\allowbreak{} & Sqrt[\allowbreak{}3 +\allowbreak{} 2*\allowbreak{}(\allowbreak{}1 +\allowbreak{} Sqrt[\allowbreak{}2]\allowbreak{})\allowbreak{}]\allowbreak{} & yes & 3$\to$2 & 0.093\\
A30 & Strad[\allowbreak{}Sqrt[\allowbreak{}16 -\allowbreak{} 2*\allowbreak{}Sqrt[\allowbreak{}29]\allowbreak{} +\allowbreak{} 2*\allowbreak{}Sqrt[\allowbreak{}55 -\allowbreak{} 10*\allowbreak{}Sqrt[\allowbreak{}29]\allowbreak{}]\allowbreak{}]\allowbreak{},\allowbreak{} True]\allowbreak{} & Sqrt[\allowbreak{}5]\allowbreak{} +\allowbreak{} Sqrt[\allowbreak{}11 -\allowbreak{} 2*\allowbreak{}Sqrt[\allowbreak{}29]\allowbreak{}]\allowbreak{} & yes & 3$\to$2 & 1.165\\
A31 & Strad[\allowbreak{}Sqrt[\allowbreak{}16 -\allowbreak{} 2*\allowbreak{}Sqrt[\allowbreak{}29]\allowbreak{} +\allowbreak{} 2*\allowbreak{}Sqrt[\allowbreak{}55 -\allowbreak{} 10*\allowbreak{}Sqrt[\allowbreak{}29]\allowbreak{}]\allowbreak{}]\allowbreak{}]\allowbreak{} & Sqrt[\allowbreak{}5]\allowbreak{} +\allowbreak{} Sqrt[\allowbreak{}11 -\allowbreak{} 2*\allowbreak{}Sqrt[\allowbreak{}29]\allowbreak{}]\allowbreak{} & yes & 3$\to$2 & 0.531\\
\end{longtable}
\begin{longtable}{@{}p{0.055\textwidth}>{\ttfamily\footnotesize\raggedright\arraybackslash}p{0.33\textwidth}>{\ttfamily\footnotesize\raggedright\arraybackslash}p{0.33\textwidth}p{0.05\textwidth}p{0.05\textwidth}p{0.05\textwidth}@{}}
\caption{Corrected version on non-denestable inputs (B) and on the branch counterexamples (C).}\label{tab:fixedBC}\\
\toprule ID & Input & Output & Equal & Depth & Time (s)\\\midrule\endfirsthead
\toprule ID & Input & Output & Equal & Depth & Time (s)\\\midrule\endhead
\bottomrule\endfoot
B01 & Strad[\allowbreak{}Sqrt[\allowbreak{}2 +\allowbreak{} Sqrt[\allowbreak{}2]\allowbreak{}]\allowbreak{}]\allowbreak{} & Sqrt[\allowbreak{}2 +\allowbreak{} Sqrt[\allowbreak{}2]\allowbreak{}]\allowbreak{} & yes & 2$\to$2 & 0.041\\
B02 & Strad[\allowbreak{}Sqrt[\allowbreak{}1 +\allowbreak{} Sqrt[\allowbreak{}2]\allowbreak{}]\allowbreak{}]\allowbreak{} & Sqrt[\allowbreak{}1 +\allowbreak{} Sqrt[\allowbreak{}2]\allowbreak{}]\allowbreak{} & yes & 2$\to$2 & 0.023\\
B03 & Strad[\allowbreak{}(\allowbreak{}1 +\allowbreak{} Sqrt[\allowbreak{}2]\allowbreak{})\allowbreak{}\textasciicircum{}\allowbreak{}(\allowbreak{}1/\allowbreak{}3)\allowbreak{}]\allowbreak{} & (\allowbreak{}1 +\allowbreak{} Sqrt[\allowbreak{}2]\allowbreak{})\allowbreak{}\textasciicircum{}\allowbreak{}(\allowbreak{}1/\allowbreak{}3)\allowbreak{} & yes & 2$\to$2 & 0.069\\
B04 & Strad[\allowbreak{}Sqrt[\allowbreak{}2]\allowbreak{}]\allowbreak{} & Sqrt[\allowbreak{}2]\allowbreak{} & yes & 1$\to$1 & 0.003\\
B05 & Strad[\allowbreak{}2]\allowbreak{} & 2 & yes & 0$\to$0 & 0.001\\
B06 & Strad[\allowbreak{}Sqrt[\allowbreak{}1 +\allowbreak{} I]\allowbreak{}]\allowbreak{} & Sqrt[\allowbreak{}1 +\allowbreak{} I]\allowbreak{} & yes & 1$\to$1 & 0.003\\
B07 & Strad[\allowbreak{}Root[\allowbreak{}\#1\textasciicircum{}\allowbreak{}4 -\allowbreak{} 10*\allowbreak{}\#1\textasciicircum{}\allowbreak{}2 +\allowbreak{} 1 \& ,\allowbreak{} 4]\allowbreak{}]\allowbreak{} & Sqrt[\allowbreak{}2]\allowbreak{} +\allowbreak{} Sqrt[\allowbreak{}3]\allowbreak{} & yes & 0$\to$1 & 0.017\\
B08 & Strad[\allowbreak{}Sqrt[\allowbreak{}Sqrt[\allowbreak{}2]\allowbreak{} +\allowbreak{} Sqrt[\allowbreak{}3]\allowbreak{}]\allowbreak{}]\allowbreak{} & Sqrt[\allowbreak{}Sqrt[\allowbreak{}2]\allowbreak{} +\allowbreak{} Sqrt[\allowbreak{}3]\allowbreak{}]\allowbreak{} & yes & 2$\to$2 & 0.080\\
B09 & Strad[\allowbreak{}(\allowbreak{}2 +\allowbreak{} 2*\allowbreak{}I*\allowbreak{}Sqrt[\allowbreak{}3]\allowbreak{})\allowbreak{}\textasciicircum{}\allowbreak{}(\allowbreak{}1/\allowbreak{}3)\allowbreak{}]\allowbreak{} & (\allowbreak{}2 +\allowbreak{} (\allowbreak{}2*\allowbreak{}I)\allowbreak{}*\allowbreak{}Sqrt[\allowbreak{}3]\allowbreak{})\allowbreak{}\textasciicircum{}\allowbreak{}(\allowbreak{}1/\allowbreak{}3)\allowbreak{} & yes & 2$\to$2 & 0.269\\
C01 & Strad[\allowbreak{}-Sqrt[\allowbreak{}3 +\allowbreak{} 2*\allowbreak{}Sqrt[\allowbreak{}2]\allowbreak{}]\allowbreak{}]\allowbreak{} & -1 -\allowbreak{} Sqrt[\allowbreak{}2]\allowbreak{} & yes & 2$\to$1 & 0.027\\
C02 & Strad[\allowbreak{}(\allowbreak{}-1 +\allowbreak{} 2*\allowbreak{}I*\allowbreak{}Sqrt[\allowbreak{}2]\allowbreak{})\allowbreak{}\textasciicircum{}\allowbreak{}(\allowbreak{}3/\allowbreak{}2)\allowbreak{}]\allowbreak{} & -5 +\allowbreak{} I*\allowbreak{}Sqrt[\allowbreak{}2]\allowbreak{} & yes & 2$\to$1 & 0.043\\
C03 & Strad[\allowbreak{}Sqrt[\allowbreak{}-1 +\allowbreak{} 2*\allowbreak{}I*\allowbreak{}Sqrt[\allowbreak{}2]\allowbreak{}]\allowbreak{}*\allowbreak{}Sqrt[\allowbreak{}-2 +\allowbreak{} 2*\allowbreak{}I*\allowbreak{}Sqrt[\allowbreak{}3]\allowbreak{}]\allowbreak{}]\allowbreak{} & -(\allowbreak{}(\allowbreak{}-I +\allowbreak{} Sqrt[\allowbreak{}2]\allowbreak{})\allowbreak{}*\allowbreak{}(\allowbreak{}-I +\allowbreak{} Sqrt[\allowbreak{}3]\allowbreak{})\allowbreak{})\allowbreak{} & yes & 2$\to$1 & 0.138\\
C04 & Strad[\allowbreak{}(\allowbreak{}-3 -\allowbreak{} 2*\allowbreak{}Sqrt[\allowbreak{}2]\allowbreak{})\allowbreak{}\textasciicircum{}\allowbreak{}(\allowbreak{}3/\allowbreak{}2)\allowbreak{}]\allowbreak{} & (\allowbreak{}-I)\allowbreak{}*\allowbreak{}(\allowbreak{}7 +\allowbreak{} 5*\allowbreak{}Sqrt[\allowbreak{}2]\allowbreak{})\allowbreak{} & yes & 2$\to$1 & 0.054\\
C05 & Strad[\allowbreak{}Sqrt[\allowbreak{}-2 -\allowbreak{} 2*\allowbreak{}Sqrt[\allowbreak{}2]\allowbreak{}]\allowbreak{}*\allowbreak{}Sqrt[\allowbreak{}-3 -\allowbreak{} 2*\allowbreak{}Sqrt[\allowbreak{}2]\allowbreak{}]\allowbreak{}]\allowbreak{} & -Sqrt[\allowbreak{}2*\allowbreak{}(\allowbreak{}7 +\allowbreak{} 5*\allowbreak{}Sqrt[\allowbreak{}2]\allowbreak{})\allowbreak{}]\allowbreak{} & yes & 2$\to$2 & 0.050\\
C06 & Strad[\allowbreak{}Sqrt[\allowbreak{}-3 -\allowbreak{} 2*\allowbreak{}Sqrt[\allowbreak{}2]\allowbreak{}]\allowbreak{}*\allowbreak{}Sqrt[\allowbreak{}-3 -\allowbreak{} 2*\allowbreak{}Sqrt[\allowbreak{}2]\allowbreak{}]\allowbreak{}]\allowbreak{} & -3 -\allowbreak{} 2*\allowbreak{}Sqrt[\allowbreak{}2]\allowbreak{} & yes & 1$\to$1 & 0.009\\
C07 & Strad[\allowbreak{}(\allowbreak{}-1 +\allowbreak{} 2*\allowbreak{}I*\allowbreak{}Sqrt[\allowbreak{}2]\allowbreak{})\allowbreak{}\textasciicircum{}\allowbreak{}(\allowbreak{}5/\allowbreak{}2)\allowbreak{}]\allowbreak{} & 1 -\allowbreak{} (\allowbreak{}11*\allowbreak{}I)\allowbreak{}*\allowbreak{}Sqrt[\allowbreak{}2]\allowbreak{} & yes & 2$\to$1 & 0.045\\
C08 & Strad[\allowbreak{}(\allowbreak{}-5 -\allowbreak{} 2*\allowbreak{}Sqrt[\allowbreak{}6]\allowbreak{})\allowbreak{}\textasciicircum{}\allowbreak{}(\allowbreak{}3/\allowbreak{}2)\allowbreak{}]\allowbreak{} & (\allowbreak{}-I)\allowbreak{}*\allowbreak{}(\allowbreak{}11*\allowbreak{}Sqrt[\allowbreak{}2]\allowbreak{} +\allowbreak{} 9*\allowbreak{}Sqrt[\allowbreak{}3]\allowbreak{})\allowbreak{} & yes & 2$\to$1 & 0.091\\
\end{longtable}
\begin{longtable}{@{}p{0.055\textwidth}>{\ttfamily\footnotesize\raggedright\arraybackslash}p{0.33\textwidth}>{\ttfamily\footnotesize\raggedright\arraybackslash}p{0.33\textwidth}p{0.05\textwidth}p{0.05\textwidth}p{0.05\textwidth}@{}}
\caption{Corrected version: direct core calls with forced multipliers (D) and symbolic input (E).}\label{tab:fixedDE}\\
\toprule ID & Input & Output & Equal & Depth & Time (s)\\\midrule\endfirsthead
\toprule ID & Input & Output & Equal & Depth & Time (s)\\\midrule\endhead
\bottomrule\endfoot
D01 & DenestCore[\allowbreak{}-Sqrt[\allowbreak{}3 +\allowbreak{} 2*\allowbreak{}Sqrt[\allowbreak{}2]\allowbreak{}]\allowbreak{},\allowbreak{} "Multipliers" -> \{1\}]\allowbreak{} & -1 -\allowbreak{} Sqrt[\allowbreak{}2]\allowbreak{} & yes & 2$\to$1 & 0.029\\
D02 & DenestCore[\allowbreak{}Sqrt[\allowbreak{}3 +\allowbreak{} 2*\allowbreak{}Sqrt[\allowbreak{}2]\allowbreak{}]\allowbreak{},\allowbreak{} "Multipliers" -> \{0\}]\allowbreak{} & 1 +\allowbreak{} Sqrt[\allowbreak{}2]\allowbreak{} & yes & 2$\to$1 & 0.012\\
D03 & DenestCore[\allowbreak{}Sqrt[\allowbreak{}3 +\allowbreak{} 2*\allowbreak{}Sqrt[\allowbreak{}2]\allowbreak{}]\allowbreak{},\allowbreak{} "Multipliers" -> \{\}]\allowbreak{} & 1 +\allowbreak{} Sqrt[\allowbreak{}2]\allowbreak{} & yes & 2$\to$1 & 0.006\\
D04 & DenestCore[\allowbreak{}(\allowbreak{}-1 +\allowbreak{} 2*\allowbreak{}I*\allowbreak{}Sqrt[\allowbreak{}2]\allowbreak{})\allowbreak{}\textasciicircum{}\allowbreak{}(\allowbreak{}3/\allowbreak{}2)\allowbreak{},\allowbreak{} "Multipliers" -> \{1\}]\allowbreak{} & -5 +\allowbreak{} I*\allowbreak{}Sqrt[\allowbreak{}2]\allowbreak{} & yes & 2$\to$1 & 0.028\\
D05 & DenestCore[\allowbreak{}Sqrt[\allowbreak{}-1 +\allowbreak{} 2*\allowbreak{}I*\allowbreak{}Sqrt[\allowbreak{}2]\allowbreak{}]\allowbreak{}*\allowbreak{}Sqrt[\allowbreak{}-2 +\allowbreak{} 2*\allowbreak{}I*\allowbreak{}Sqrt[\allowbreak{}3]\allowbreak{}]\allowbreak{},\allowbreak{} "Multipliers" -> \{1\}]\allowbreak{} & -(\allowbreak{}(\allowbreak{}-I +\allowbreak{} Sqrt[\allowbreak{}2]\allowbreak{})\allowbreak{}*\allowbreak{}(\allowbreak{}-I +\allowbreak{} Sqrt[\allowbreak{}3]\allowbreak{})\allowbreak{})\allowbreak{} & yes & 2$\to$1 & 0.089\\
D06 & DenestCore[\allowbreak{}Sqrt[\allowbreak{}5 +\allowbreak{} 2*\allowbreak{}Sqrt[\allowbreak{}6]\allowbreak{}]\allowbreak{},\allowbreak{} "Multipliers" -> \{1,\allowbreak{} 6 -\allowbreak{} 2*\allowbreak{}Sqrt[\allowbreak{}6]\allowbreak{} -\allowbreak{} 2*\allowbreak{}Sqrt[\allowbreak{}2]\allowbreak{} +\allowbreak{} 2*\allowbreak{}Sqrt[\allowbreak{}3]\allowbreak{}\}]\allowbreak{} & Sqrt[\allowbreak{}2]\allowbreak{} +\allowbreak{} Sqrt[\allowbreak{}3]\allowbreak{} & yes & 2$\to$1 & 0.016\\
D07 & DenestCore[\allowbreak{}Sqrt[\allowbreak{}5 +\allowbreak{} 2*\allowbreak{}Sqrt[\allowbreak{}6]\allowbreak{}]\allowbreak{},\allowbreak{} "Multipliers" -> \{6 -\allowbreak{} 2*\allowbreak{}Sqrt[\allowbreak{}6]\allowbreak{} -\allowbreak{} 2*\allowbreak{}Sqrt[\allowbreak{}2]\allowbreak{} +\allowbreak{} 2*\allowbreak{}Sqrt[\allowbreak{}3]\allowbreak{}\}]\allowbreak{} & Sqrt[\allowbreak{}2]\allowbreak{} +\allowbreak{} Sqrt[\allowbreak{}3]\allowbreak{} & yes & 2$\to$1 & 0.016\\
D08 & DenestCore[\allowbreak{}Sqrt[\allowbreak{}3 +\allowbreak{} 2*\allowbreak{}Sqrt[\allowbreak{}3 +\allowbreak{} 2*\allowbreak{}Sqrt[\allowbreak{}2]\allowbreak{}]\allowbreak{}]\allowbreak{}]\allowbreak{} & Sqrt[\allowbreak{}3 +\allowbreak{} 2*\allowbreak{}Sqrt[\allowbreak{}3 +\allowbreak{} 2*\allowbreak{}Sqrt[\allowbreak{}2]\allowbreak{}]\allowbreak{}]\allowbreak{} & yes & 3$\to$3 & 0.114\\
D09 & DenestCore[\allowbreak{}2 +\allowbreak{} Sqrt[\allowbreak{}3]\allowbreak{}]\allowbreak{} & 2 +\allowbreak{} Sqrt[\allowbreak{}3]\allowbreak{} & yes & 1$\to$1 & 0.001\\
D10 & DenestCore[\allowbreak{}Sqrt[\allowbreak{}3 +\allowbreak{} 2*\allowbreak{}Sqrt[\allowbreak{}2]\allowbreak{}]\allowbreak{},\allowbreak{} "Multipliers" -> \{2.\}]\allowbreak{} & 1 +\allowbreak{} Sqrt[\allowbreak{}2]\allowbreak{} & yes & 2$\to$1 & 0.013\\
D11 & DenestCore[\allowbreak{}Sqrt[\allowbreak{}3 +\allowbreak{} 2*\allowbreak{}Sqrt[\allowbreak{}2]\allowbreak{}]\allowbreak{},\allowbreak{} "Multipliers" -> \{Pi\}]\allowbreak{} & 1 +\allowbreak{} Sqrt[\allowbreak{}2]\allowbreak{} & yes & 2$\to$1 & 0.013\\
D12 & DenestCore[\allowbreak{}Sqrt[\allowbreak{}5 +\allowbreak{} 2*\allowbreak{}Sqrt[\allowbreak{}6]\allowbreak{}]\allowbreak{},\allowbreak{} "Multipliers" -> \{1,\allowbreak{} Sqrt[\allowbreak{}5]\allowbreak{}\}]\allowbreak{} & Sqrt[\allowbreak{}2]\allowbreak{} +\allowbreak{} Sqrt[\allowbreak{}3]\allowbreak{} & yes & 2$\to$1 & 0.013\\
D13 & DenestCore[\allowbreak{}Sqrt[\allowbreak{}3 +\allowbreak{} 2*\allowbreak{}Sqrt[\allowbreak{}2]\allowbreak{}]\allowbreak{},\allowbreak{} "Multipliers" -> \{-1\}]\allowbreak{} & 1 +\allowbreak{} Sqrt[\allowbreak{}2]\allowbreak{} & yes & 2$\to$1 & 0.012\\
D14 & DenestCore[\allowbreak{}Sqrt[\allowbreak{}3 +\allowbreak{} 2*\allowbreak{}Sqrt[\allowbreak{}2]\allowbreak{}]\allowbreak{},\allowbreak{} "Multipliers" -> \{I\}]\allowbreak{} & 1 +\allowbreak{} Sqrt[\allowbreak{}2]\allowbreak{} & yes & 2$\to$1 & 0.012\\
D15 & DenestCore[\allowbreak{}Sqrt[\allowbreak{}5 +\allowbreak{} 2*\allowbreak{}Sqrt[\allowbreak{}6]\allowbreak{}]\allowbreak{},\allowbreak{} "Multipliers" -> \{1\}]\allowbreak{} & Sqrt[\allowbreak{}2]\allowbreak{} +\allowbreak{} Sqrt[\allowbreak{}3]\allowbreak{} & yes & 2$\to$1 & 0.019\\
E01 & Strad[\allowbreak{}Sqrt[\allowbreak{}z\textasciicircum{}\allowbreak{}2]\allowbreak{}]\allowbreak{} & Sqrt[\allowbreak{}z\textasciicircum{}\allowbreak{}2]\allowbreak{} & -- & 1$\to$1 & 0.017\\
E02 & Strad[\allowbreak{}Sqrt[\allowbreak{}t +\allowbreak{} t\textasciicircum{}\allowbreak{}2]\allowbreak{}]\allowbreak{} & Sqrt[\allowbreak{}t +\allowbreak{} t\textasciicircum{}\allowbreak{}2]\allowbreak{} & -- & 1$\to$1 & 0.002\\
E03 & Strad[\allowbreak{}Sqrt[\allowbreak{}-x -\allowbreak{} y]\allowbreak{}]\allowbreak{} & Sqrt[\allowbreak{}-x -\allowbreak{} y]\allowbreak{} & -- & 1$\to$1 & 0.002\\
E04 & Strad[\allowbreak{}Sqrt[\allowbreak{}a +\allowbreak{} b]\allowbreak{}]\allowbreak{} & Sqrt[\allowbreak{}a +\allowbreak{} b]\allowbreak{} & -- & 1$\to$1 & 0.002\\
E05 & Strad[\allowbreak{}x]\allowbreak{} & x & -- & 0$\to$0 & 0.000\\
E06 & Strad[\allowbreak{}Sqrt[\allowbreak{}x]\allowbreak{}*\allowbreak{}Sqrt[\allowbreak{}3 +\allowbreak{} 2*\allowbreak{}Sqrt[\allowbreak{}2]\allowbreak{}]\allowbreak{}]\allowbreak{} & (\allowbreak{}1 +\allowbreak{} Sqrt[\allowbreak{}2]\allowbreak{})\allowbreak{}*\allowbreak{}Sqrt[\allowbreak{}x]\allowbreak{} & -- & 2$\to$1 & 0.018\\
E07 & Strad[\allowbreak{}f[\allowbreak{}1]\allowbreak{}]\allowbreak{} & f[\allowbreak{}1]\allowbreak{} & -- & 0$\to$0 & 0.001\\
E08 & Strad[\allowbreak{}f[\allowbreak{}Sqrt[\allowbreak{}3 +\allowbreak{} 2*\allowbreak{}Sqrt[\allowbreak{}2]\allowbreak{}]\allowbreak{}]\allowbreak{}]\allowbreak{} & f[\allowbreak{}1 +\allowbreak{} Sqrt[\allowbreak{}2]\allowbreak{}]\allowbreak{} & -- & 2$\to$1 & 0.014\\
E09 & Strad[\allowbreak{}RadicalDenest`Private`mark[\allowbreak{}1]\allowbreak{}]\allowbreak{} & 1 & -- & 0$\to$0 & 0.001\\
\end{longtable}
\begin{longtable}{@{}p{0.055\textwidth}>{\ttfamily\footnotesize\raggedright\arraybackslash}p{0.33\textwidth}>{\ttfamily\footnotesize\raggedright\arraybackslash}p{0.33\textwidth}p{0.05\textwidth}p{0.05\textwidth}p{0.05\textwidth}@{}}
\caption{Corrected version: factorizer and rationalizer probes (J) and miscellaneous inputs (K).}\label{tab:fixedJK}\\
\toprule ID & Input & Output & Equal & Depth & Time (s)\\\midrule\endfirsthead
\toprule ID & Input & Output & Equal & Depth & Time (s)\\\midrule\endhead
\bottomrule\endfoot
J01 & Strad[\allowbreak{}Sqrt[\allowbreak{}-(\allowbreak{}Sqrt[\allowbreak{}2]\allowbreak{} +\allowbreak{} (\allowbreak{}-1)\allowbreak{}\textasciicircum{}\allowbreak{}(\allowbreak{}1/\allowbreak{}3)\allowbreak{})\allowbreak{}]\allowbreak{}]\allowbreak{} & Sqrt[\allowbreak{}-(\allowbreak{}-1)\allowbreak{}\textasciicircum{}\allowbreak{}(\allowbreak{}1/\allowbreak{}3)\allowbreak{} -\allowbreak{} Sqrt[\allowbreak{}2]\allowbreak{}]\allowbreak{} & yes & 2$\to$2 & 0.480\\
J02 & Strad[\allowbreak{}Sqrt[\allowbreak{}-(\allowbreak{}Sqrt[\allowbreak{}2]\allowbreak{} +\allowbreak{} (\allowbreak{}-1)\allowbreak{}\textasciicircum{}\allowbreak{}(\allowbreak{}1/\allowbreak{}3)\allowbreak{})\allowbreak{}]\allowbreak{} +\allowbreak{} 1]\allowbreak{} & 1 +\allowbreak{} Sqrt[\allowbreak{}-(\allowbreak{}-1)\allowbreak{}\textasciicircum{}\allowbreak{}(\allowbreak{}1/\allowbreak{}3)\allowbreak{} -\allowbreak{} Sqrt[\allowbreak{}2]\allowbreak{}]\allowbreak{} & yes & 2$\to$2 & 0.437\\
J06 & Strad[\allowbreak{}(\allowbreak{}-(\allowbreak{}Sqrt[\allowbreak{}2]\allowbreak{} +\allowbreak{} (\allowbreak{}-1)\allowbreak{}\textasciicircum{}\allowbreak{}(\allowbreak{}1/\allowbreak{}3)\allowbreak{})\allowbreak{})\allowbreak{}\textasciicircum{}\allowbreak{}(\allowbreak{}1/\allowbreak{}3)\allowbreak{}]\allowbreak{} & (\allowbreak{}-(\allowbreak{}-1)\allowbreak{}\textasciicircum{}\allowbreak{}(\allowbreak{}1/\allowbreak{}3)\allowbreak{} -\allowbreak{} Sqrt[\allowbreak{}2]\allowbreak{})\allowbreak{}\textasciicircum{}\allowbreak{}(\allowbreak{}1/\allowbreak{}3)\allowbreak{} & yes & 2$\to$2 & 1.826\\
J09 & Factorc[\allowbreak{}Sqrt[\allowbreak{}-(\allowbreak{}Sqrt[\allowbreak{}2]\allowbreak{} +\allowbreak{} (\allowbreak{}-1)\allowbreak{}\textasciicircum{}\allowbreak{}(\allowbreak{}1/\allowbreak{}3)\allowbreak{})\allowbreak{}]\allowbreak{}]\allowbreak{} & Sqrt[\allowbreak{}-(\allowbreak{}-1)\allowbreak{}\textasciicircum{}\allowbreak{}(\allowbreak{}1/\allowbreak{}3)\allowbreak{} -\allowbreak{} Sqrt[\allowbreak{}2]\allowbreak{}]\allowbreak{} & yes & 2$\to$2 & 0.002\\
J10 & Factorc[\allowbreak{}Sqrt[\allowbreak{}-3 -\allowbreak{} 2*\allowbreak{}Sqrt[\allowbreak{}2]\allowbreak{}]\allowbreak{}]\allowbreak{} & I*\allowbreak{}Sqrt[\allowbreak{}3 +\allowbreak{} 2*\allowbreak{}Sqrt[\allowbreak{}2]\allowbreak{}]\allowbreak{} & yes & 2$\to$2 & 0.001\\
J11 & Factorc[\allowbreak{}Sqrt[\allowbreak{}8 +\allowbreak{} 4*\allowbreak{}Sqrt[\allowbreak{}3]\allowbreak{}]\allowbreak{}]\allowbreak{} & 2*\allowbreak{}Sqrt[\allowbreak{}2 +\allowbreak{} Sqrt[\allowbreak{}3]\allowbreak{}]\allowbreak{} & yes & 2$\to$2 & 0.006\\
J12 & Factorc[\allowbreak{}Sqrt[\allowbreak{}z\textasciicircum{}\allowbreak{}2]\allowbreak{}]\allowbreak{} & Sqrt[\allowbreak{}z\textasciicircum{}\allowbreak{}2]\allowbreak{} & -- & 1$\to$1 & 0.001\\
J13 & Factorc[\allowbreak{}Sqrt[\allowbreak{}-x -\allowbreak{} y]\allowbreak{}]\allowbreak{} & Sqrt[\allowbreak{}-x -\allowbreak{} y]\allowbreak{} & -- & 1$\to$1 & 0.002\\
J14 & Factorc[\allowbreak{}Sqrt[\allowbreak{}t +\allowbreak{} t\textasciicircum{}\allowbreak{}2]\allowbreak{}]\allowbreak{} & Sqrt[\allowbreak{}t*\allowbreak{}(\allowbreak{}1 +\allowbreak{} t)\allowbreak{}]\allowbreak{} & -- & 1$\to$1 & 0.001\\
J15 & RationalizeDenominator[\allowbreak{}x/\allowbreak{}(\allowbreak{}1 +\allowbreak{} Sqrt[\allowbreak{}2]\allowbreak{})\allowbreak{}]\allowbreak{} & -x +\allowbreak{} Sqrt[\allowbreak{}2]\allowbreak{}*\allowbreak{}x & -- & 1$\to$1 & 0.001\\
J16 & RationalizeDenominator[\allowbreak{}1/\allowbreak{}2 +\allowbreak{} 1/\allowbreak{}Sqrt[\allowbreak{}2]\allowbreak{}]\allowbreak{} & (\allowbreak{}1 +\allowbreak{} Sqrt[\allowbreak{}2]\allowbreak{})\allowbreak{}/\allowbreak{}2 & yes & 1$\to$1 & 0.001\\
J17 & RationalizeDenominator[\allowbreak{}1/\allowbreak{}x]\allowbreak{} & x\textasciicircum{}\allowbreak{}(\allowbreak{}-1)\allowbreak{} & -- & 0$\to$0 & 0.000\\
K01 & Strad[\allowbreak{}Sqrt[\allowbreak{}28\textasciicircum{}\allowbreak{}(\allowbreak{}1/\allowbreak{}3)\allowbreak{} -\allowbreak{} 3]\allowbreak{}]\allowbreak{} & (\allowbreak{}-1 -\allowbreak{} 2\textasciicircum{}\allowbreak{}(\allowbreak{}2/\allowbreak{}3)\allowbreak{}*\allowbreak{}7\textasciicircum{}\allowbreak{}(\allowbreak{}1/\allowbreak{}3)\allowbreak{} +\allowbreak{} 2\textasciicircum{}\allowbreak{}(\allowbreak{}1/\allowbreak{}3)\allowbreak{}*\allowbreak{}7\textasciicircum{}\allowbreak{}(\allowbreak{}2/\allowbreak{}3)\allowbreak{})\allowbreak{}/\allowbreak{}3 & yes & 2$\to$1 & 0.121\\
K02 & Strad[\allowbreak{}(\allowbreak{}1 +\allowbreak{} 2\textasciicircum{}\allowbreak{}(\allowbreak{}1/\allowbreak{}3)\allowbreak{})\allowbreak{}\textasciicircum{}\allowbreak{}3]\allowbreak{} & (\allowbreak{}1 +\allowbreak{} 2\textasciicircum{}\allowbreak{}(\allowbreak{}1/\allowbreak{}3)\allowbreak{})\allowbreak{}\textasciicircum{}\allowbreak{}3 & yes & 1$\to$1 & 0.007\\
K04 & Strad[\allowbreak{}Strad[\allowbreak{}Sqrt[\allowbreak{}16 -\allowbreak{} 2*\allowbreak{}Sqrt[\allowbreak{}29]\allowbreak{} +\allowbreak{} 2*\allowbreak{}Sqrt[\allowbreak{}55 -\allowbreak{} 10*\allowbreak{}Sqrt[\allowbreak{}29]\allowbreak{}]\allowbreak{}]\allowbreak{},\allowbreak{} True]\allowbreak{},\allowbreak{} True]\allowbreak{} & Sqrt[\allowbreak{}5]\allowbreak{} +\allowbreak{} Sqrt[\allowbreak{}11 -\allowbreak{} 2*\allowbreak{}Sqrt[\allowbreak{}29]\allowbreak{}]\allowbreak{} & yes & 2$\to$2 & 1.762\\
K09 & Strad[\allowbreak{}Sqrt[\allowbreak{}5 +\allowbreak{} 2*\allowbreak{}Sqrt[\allowbreak{}6]\allowbreak{}]\allowbreak{}/\allowbreak{}Sqrt[\allowbreak{}3 +\allowbreak{} 2*\allowbreak{}Sqrt[\allowbreak{}2]\allowbreak{}]\allowbreak{}]\allowbreak{} & 2 -\allowbreak{} Sqrt[\allowbreak{}2]\allowbreak{} -\allowbreak{} Sqrt[\allowbreak{}3]\allowbreak{} +\allowbreak{} Sqrt[\allowbreak{}6]\allowbreak{} & yes & 2$\to$1 & 0.284\\
K10 & Strad[\allowbreak{}Exp[\allowbreak{}I*\allowbreak{}(\allowbreak{}Pi/\allowbreak{}7)\allowbreak{}]\allowbreak{}*\allowbreak{}Sqrt[\allowbreak{}3 +\allowbreak{} 2*\allowbreak{}Sqrt[\allowbreak{}2]\allowbreak{}]\allowbreak{}]\allowbreak{} & (\allowbreak{}-1)\allowbreak{}\textasciicircum{}\allowbreak{}(\allowbreak{}1/\allowbreak{}7)\allowbreak{}*\allowbreak{}(\allowbreak{}1 +\allowbreak{} Sqrt[\allowbreak{}2]\allowbreak{})\allowbreak{} & yes & 2$\to$1 & 0.107\\
K11 & Strad[\allowbreak{}Sqrt[\allowbreak{}3 +\allowbreak{} 2*\allowbreak{}Sqrt[\allowbreak{}2]\allowbreak{}]\allowbreak{} +\allowbreak{} Pi]\allowbreak{} & 1 +\allowbreak{} Sqrt[\allowbreak{}2]\allowbreak{} +\allowbreak{} Pi & yes & 2$\to$1 & 0.035\\
K12 & Strad[\allowbreak{}Sqrt[\allowbreak{}3. +\allowbreak{} 2*\allowbreak{}Sqrt[\allowbreak{}2]\allowbreak{}]\allowbreak{}]\allowbreak{} & 2.414213562373095 & yes & 0$\to$0 & 0.000\\
K14 & Strad[\allowbreak{}\{Sqrt[\allowbreak{}3 +\allowbreak{} 2*\allowbreak{}Sqrt[\allowbreak{}2]\allowbreak{}]\allowbreak{},\allowbreak{} Sqrt[\allowbreak{}5 +\allowbreak{} 2*\allowbreak{}Sqrt[\allowbreak{}6]\allowbreak{}]\allowbreak{}\}]\allowbreak{} & \{1 +\allowbreak{} Sqrt[\allowbreak{}2]\allowbreak{},\allowbreak{} Sqrt[\allowbreak{}2]\allowbreak{} +\allowbreak{} Sqrt[\allowbreak{}3]\allowbreak{}\} & -- & 2$\to$1 & 0.103\\
K15 & Strad[\allowbreak{}Sqrt[\allowbreak{}3 +\allowbreak{} 2*\allowbreak{}Sqrt[\allowbreak{}2]\allowbreak{}]\allowbreak{} == 1 +\allowbreak{} Sqrt[\allowbreak{}2]\allowbreak{}]\allowbreak{} & True & -- & 0$\to$0 & 0.002\\
K17 & Strad[\allowbreak{}Sqrt[\allowbreak{}(\allowbreak{}3 +\allowbreak{} 2*\allowbreak{}Sqrt[\allowbreak{}2]\allowbreak{})\allowbreak{}/\allowbreak{}(\allowbreak{}5 +\allowbreak{} 2*\allowbreak{}Sqrt[\allowbreak{}6]\allowbreak{})\allowbreak{}]\allowbreak{}]\allowbreak{} & -2 -\allowbreak{} Sqrt[\allowbreak{}2]\allowbreak{} +\allowbreak{} Sqrt[\allowbreak{}3]\allowbreak{} +\allowbreak{} Sqrt[\allowbreak{}6]\allowbreak{} & yes & 2$\to$1 & 0.206\\
K19 & Strad[\allowbreak{}Sqrt[\allowbreak{}Sqrt[\allowbreak{}2]\allowbreak{} +\allowbreak{} 1]\allowbreak{}\textasciicircum{}\allowbreak{}3]\allowbreak{} & (\allowbreak{}1 +\allowbreak{} Sqrt[\allowbreak{}2]\allowbreak{})\allowbreak{}\textasciicircum{}\allowbreak{}(\allowbreak{}3/\allowbreak{}2)\allowbreak{} & yes & 2$\to$2 & 0.065\\
K20 & Strad[\allowbreak{}Sqrt[\allowbreak{}2 +\allowbreak{} Sqrt[\allowbreak{}3]\allowbreak{}]\allowbreak{}*\allowbreak{}Sqrt[\allowbreak{}3 +\allowbreak{} 2*\allowbreak{}Sqrt[\allowbreak{}2]\allowbreak{}]\allowbreak{}]\allowbreak{} & (\allowbreak{}(\allowbreak{}2 +\allowbreak{} Sqrt[\allowbreak{}2]\allowbreak{})\allowbreak{}*\allowbreak{}(\allowbreak{}1 +\allowbreak{} Sqrt[\allowbreak{}3]\allowbreak{})\allowbreak{})\allowbreak{}/\allowbreak{}2 & yes & 2$\to$1 & 0.148\\
K21 & Strad[\allowbreak{}Sqrt[\allowbreak{}3 +\allowbreak{} 2*\allowbreak{}Sqrt[\allowbreak{}2]\allowbreak{}]\allowbreak{},\allowbreak{} "Verbose" -> False,\allowbreak{} "MaxTrials" -> 1]\allowbreak{} & 1 +\allowbreak{} Sqrt[\allowbreak{}2]\allowbreak{} & yes & 2$\to$1 & 0.035\\
K22 & Strad[\allowbreak{}Sqrt[\allowbreak{}5 +\allowbreak{} 2*\allowbreak{}Sqrt[\allowbreak{}6]\allowbreak{}]\allowbreak{},\allowbreak{} "Factor" -> False]\allowbreak{} & Sqrt[\allowbreak{}2]\allowbreak{} +\allowbreak{} Sqrt[\allowbreak{}3]\allowbreak{} & yes & 2$\to$1 & 0.041\\
\end{longtable}
\begin{longtable}{@{}p{0.055\textwidth}>{\ttfamily\footnotesize\raggedright\arraybackslash}p{0.33\textwidth}>{\ttfamily\footnotesize\raggedright\arraybackslash}p{0.33\textwidth}p{0.05\textwidth}p{0.05\textwidth}p{0.05\textwidth}@{}}
\caption{Corrected version: higher-order roots (L) and roots of negative bases (Q).}\label{tab:fixedLQ}\\
\toprule ID & Input & Output & Equal & Depth & Time (s)\\\midrule\endfirsthead
\toprule ID & Input & Output & Equal & Depth & Time (s)\\\midrule\endhead
\bottomrule\endfoot
L01 & Strad[\allowbreak{}(\allowbreak{}239 +\allowbreak{} 169*\allowbreak{}Sqrt[\allowbreak{}2]\allowbreak{})\allowbreak{}\textasciicircum{}\allowbreak{}(\allowbreak{}1/\allowbreak{}7)\allowbreak{}]\allowbreak{} & 1 +\allowbreak{} Sqrt[\allowbreak{}2]\allowbreak{} & yes & 2$\to$1 & 0.067\\
L02 & Strad[\allowbreak{}(\allowbreak{}2 +\allowbreak{} Sqrt[\allowbreak{}2]\allowbreak{})\allowbreak{}\textasciicircum{}\allowbreak{}(\allowbreak{}1/\allowbreak{}5)\allowbreak{}]\allowbreak{} & (\allowbreak{}2 +\allowbreak{} Sqrt[\allowbreak{}2]\allowbreak{})\allowbreak{}\textasciicircum{}\allowbreak{}(\allowbreak{}1/\allowbreak{}5)\allowbreak{} & yes & 2$\to$2 & 0.430\\
L03 & Strad[\allowbreak{}(\allowbreak{}1 +\allowbreak{} Sqrt[\allowbreak{}2]\allowbreak{})\allowbreak{}\textasciicircum{}\allowbreak{}(\allowbreak{}1/\allowbreak{}7)\allowbreak{}]\allowbreak{} & (\allowbreak{}1 +\allowbreak{} Sqrt[\allowbreak{}2]\allowbreak{})\allowbreak{}\textasciicircum{}\allowbreak{}(\allowbreak{}1/\allowbreak{}7)\allowbreak{} & yes & 2$\to$2 & 0.773\\
L04 & Strad[\allowbreak{}Sqrt[\allowbreak{}17 +\allowbreak{} 2*\allowbreak{}Sqrt[\allowbreak{}6]\allowbreak{} +\allowbreak{} 2*\allowbreak{}Sqrt[\allowbreak{}10]\allowbreak{} +\allowbreak{} 2*\allowbreak{}Sqrt[\allowbreak{}14]\allowbreak{} +\allowbreak{} 2*\allowbreak{}Sqrt[\allowbreak{}15]\allowbreak{} +\allowbreak{} 2*\allowbreak{}Sqrt[\allowbreak{}21]\allowbreak{} +\allowbreak{} 2*\allowbreak{}Sqrt[\allowbreak{}35]\allowbreak{}]\allowbreak{}]\allowbreak{} & Sqrt[\allowbreak{}2]\allowbreak{} +\allowbreak{} Sqrt[\allowbreak{}3]\allowbreak{} +\allowbreak{} Sqrt[\allowbreak{}5]\allowbreak{} +\allowbreak{} Sqrt[\allowbreak{}7]\allowbreak{} & yes & 2$\to$1 & 2.916\\
L05 & Strad[\allowbreak{}Expand[\allowbreak{}(\allowbreak{}1 +\allowbreak{} Sqrt[\allowbreak{}2]\allowbreak{} +\allowbreak{} Sqrt[\allowbreak{}3]\allowbreak{})\allowbreak{}\textasciicircum{}\allowbreak{}3]\allowbreak{}\textasciicircum{}\allowbreak{}(\allowbreak{}1/\allowbreak{}3)\allowbreak{}]\allowbreak{} & 1 +\allowbreak{} Sqrt[\allowbreak{}2]\allowbreak{} +\allowbreak{} Sqrt[\allowbreak{}3]\allowbreak{} & yes & 2$\to$1 & 0.118\\
L06 & Strad[\allowbreak{}Expand[\allowbreak{}(\allowbreak{}Sqrt[\allowbreak{}2]\allowbreak{} +\allowbreak{} Sqrt[\allowbreak{}3]\allowbreak{})\allowbreak{}\textasciicircum{}\allowbreak{}4]\allowbreak{}\textasciicircum{}\allowbreak{}(\allowbreak{}1/\allowbreak{}4)\allowbreak{}]\allowbreak{} & Sqrt[\allowbreak{}2]\allowbreak{} +\allowbreak{} Sqrt[\allowbreak{}3]\allowbreak{} & yes & 2$\to$1 & 0.057\\
L07 & Strad[\allowbreak{}Sqrt[\allowbreak{}1000003 +\allowbreak{} 2*\allowbreak{}Sqrt[\allowbreak{}999983]\allowbreak{}]\allowbreak{}]\allowbreak{} & Sqrt[\allowbreak{}1000003 +\allowbreak{} 2*\allowbreak{}Sqrt[\allowbreak{}999983]\allowbreak{}]\allowbreak{} & yes & 2$\to$2 & 0.212\\
L08 & Strad[\allowbreak{}Sqrt[\allowbreak{}2 +\allowbreak{} Sqrt[\allowbreak{}3]\allowbreak{} +\allowbreak{} Sqrt[\allowbreak{}5]\allowbreak{}]\allowbreak{}]\allowbreak{} & Sqrt[\allowbreak{}2 +\allowbreak{} Sqrt[\allowbreak{}3]\allowbreak{} +\allowbreak{} Sqrt[\allowbreak{}5]\allowbreak{}]\allowbreak{} & yes & 2$\to$2 & 0.376\\
L09 & Strad[\allowbreak{}(\allowbreak{}5 +\allowbreak{} 2*\allowbreak{}Sqrt[\allowbreak{}6]\allowbreak{})\allowbreak{}\textasciicircum{}\allowbreak{}(\allowbreak{}1/\allowbreak{}6)\allowbreak{}]\allowbreak{} & (\allowbreak{}5 +\allowbreak{} 2*\allowbreak{}Sqrt[\allowbreak{}6]\allowbreak{})\allowbreak{}\textasciicircum{}\allowbreak{}(\allowbreak{}1/\allowbreak{}6)\allowbreak{} & yes & 2$\to$2 & 1.554\\
L10 & Strad[\allowbreak{}(\allowbreak{}11*\allowbreak{}Sqrt[\allowbreak{}2]\allowbreak{} +\allowbreak{} 9*\allowbreak{}Sqrt[\allowbreak{}3]\allowbreak{})\allowbreak{}\textasciicircum{}\allowbreak{}(\allowbreak{}1/\allowbreak{}3)\allowbreak{}]\allowbreak{} & Sqrt[\allowbreak{}2]\allowbreak{} +\allowbreak{} Sqrt[\allowbreak{}3]\allowbreak{} & yes & 2$\to$1 & 0.054\\
L11 & Strad[\allowbreak{}Expand[\allowbreak{}(\allowbreak{}2\textasciicircum{}\allowbreak{}(\allowbreak{}1/\allowbreak{}3)\allowbreak{} +\allowbreak{} 3\textasciicircum{}\allowbreak{}(\allowbreak{}1/\allowbreak{}3)\allowbreak{})\allowbreak{}\textasciicircum{}\allowbreak{}3]\allowbreak{}\textasciicircum{}\allowbreak{}(\allowbreak{}1/\allowbreak{}3)\allowbreak{}]\allowbreak{} & 2\textasciicircum{}\allowbreak{}(\allowbreak{}1/\allowbreak{}3)\allowbreak{} +\allowbreak{} 3\textasciicircum{}\allowbreak{}(\allowbreak{}1/\allowbreak{}3)\allowbreak{} & yes & 2$\to$1 & 0.266\\
L12 & Strad[\allowbreak{}Sqrt[\allowbreak{}Expand[\allowbreak{}(\allowbreak{}1 +\allowbreak{} 2\textasciicircum{}\allowbreak{}(\allowbreak{}1/\allowbreak{}3)\allowbreak{} +\allowbreak{} 3\textasciicircum{}\allowbreak{}(\allowbreak{}1/\allowbreak{}3)\allowbreak{})\allowbreak{}\textasciicircum{}\allowbreak{}2]\allowbreak{}]\allowbreak{}]\allowbreak{} & 1 +\allowbreak{} 2\textasciicircum{}\allowbreak{}(\allowbreak{}1/\allowbreak{}3)\allowbreak{} +\allowbreak{} 3\textasciicircum{}\allowbreak{}(\allowbreak{}1/\allowbreak{}3)\allowbreak{} & yes & 2$\to$1 & 0.058\\
L13 & Strad[\allowbreak{}Expand[\allowbreak{}(\allowbreak{}1 +\allowbreak{} 2\textasciicircum{}\allowbreak{}(\allowbreak{}1/\allowbreak{}5)\allowbreak{})\allowbreak{}\textasciicircum{}\allowbreak{}5]\allowbreak{}\textasciicircum{}\allowbreak{}(\allowbreak{}1/\allowbreak{}5)\allowbreak{}]\allowbreak{} & 1 +\allowbreak{} 2\textasciicircum{}\allowbreak{}(\allowbreak{}1/\allowbreak{}5)\allowbreak{} & yes & 2$\to$1 & 0.055\\
L14 & Strad[\allowbreak{}Sqrt[\allowbreak{}5 +\allowbreak{} 2*\allowbreak{}Sqrt[\allowbreak{}6]\allowbreak{}]\allowbreak{}\textasciicircum{}\allowbreak{}(\allowbreak{}1/\allowbreak{}5)\allowbreak{}]\allowbreak{} & (\allowbreak{}5 +\allowbreak{} 2*\allowbreak{}Sqrt[\allowbreak{}6]\allowbreak{})\allowbreak{}\textasciicircum{}\allowbreak{}(\allowbreak{}1/\allowbreak{}10)\allowbreak{} & yes & 2$\to$2 & 6.605\\
Q01 & Strad[\allowbreak{}Expand[\allowbreak{}(\allowbreak{}-1 -\allowbreak{} Sqrt[\allowbreak{}2]\allowbreak{})\allowbreak{}\textasciicircum{}\allowbreak{}3]\allowbreak{}\textasciicircum{}\allowbreak{}(\allowbreak{}1/\allowbreak{}3)\allowbreak{}]\allowbreak{} & (\allowbreak{}(\allowbreak{}1 +\allowbreak{} Sqrt[\allowbreak{}2]\allowbreak{})\allowbreak{}*\allowbreak{}(\allowbreak{}1 +\allowbreak{} I*\allowbreak{}Sqrt[\allowbreak{}3]\allowbreak{})\allowbreak{})\allowbreak{}/\allowbreak{}2 & yes & 2$\to$1 & 0.401\\
Q02 & Strad[\allowbreak{}Expand[\allowbreak{}(\allowbreak{}1 -\allowbreak{} Sqrt[\allowbreak{}2]\allowbreak{})\allowbreak{}\textasciicircum{}\allowbreak{}3]\allowbreak{}\textasciicircum{}\allowbreak{}(\allowbreak{}1/\allowbreak{}3)\allowbreak{}]\allowbreak{} & (\allowbreak{}(\allowbreak{}-1 +\allowbreak{} Sqrt[\allowbreak{}2]\allowbreak{})\allowbreak{}*\allowbreak{}(\allowbreak{}1 +\allowbreak{} I*\allowbreak{}Sqrt[\allowbreak{}3]\allowbreak{})\allowbreak{})\allowbreak{}/\allowbreak{}2 & yes & 2$\to$1 & 0.366\\
Q07 & Strad[\allowbreak{}Expand[\allowbreak{}(\allowbreak{}2 -\allowbreak{} 3*\allowbreak{}Sqrt[\allowbreak{}2]\allowbreak{})\allowbreak{}\textasciicircum{}\allowbreak{}3]\allowbreak{}\textasciicircum{}\allowbreak{}(\allowbreak{}1/\allowbreak{}3)\allowbreak{}]\allowbreak{} & (\allowbreak{}(\allowbreak{}-I)\allowbreak{}*\allowbreak{}(\allowbreak{}-3 +\allowbreak{} Sqrt[\allowbreak{}2]\allowbreak{})\allowbreak{}*\allowbreak{}(\allowbreak{}-I +\allowbreak{} Sqrt[\allowbreak{}3]\allowbreak{})\allowbreak{})\allowbreak{}/\allowbreak{}Sqrt[\allowbreak{}2]\allowbreak{} & yes & 2$\to$1 & 0.475\\
Q08 & Strad[\allowbreak{}(\allowbreak{}-7 -\allowbreak{} 5*\allowbreak{}Sqrt[\allowbreak{}2]\allowbreak{})\allowbreak{}\textasciicircum{}\allowbreak{}(\allowbreak{}1/\allowbreak{}3)\allowbreak{}]\allowbreak{} & (\allowbreak{}(\allowbreak{}1 +\allowbreak{} Sqrt[\allowbreak{}2]\allowbreak{})\allowbreak{}*\allowbreak{}(\allowbreak{}1 +\allowbreak{} I*\allowbreak{}Sqrt[\allowbreak{}3]\allowbreak{})\allowbreak{})\allowbreak{}/\allowbreak{}2 & yes & 2$\to$1 & 0.212\\
\end{longtable}
